\section{Étude du changement de base dans les foncteurs homologiques covariants de Modules}
\label{section:III.7.fr}

\subsection{Foncteurs de $A$-modules}
\label{subsection:III.7.1.fr}

\begin{env}[7.1.1]
Étant donné un anneau $A$ (non nécessairement commutatif), nous noterons
$\mathbf{Ab}_A$ la catégorie des $A$-modules à gauche, et nous noterons
simplement $\mathbf{Ab}$ la catégorie des $\mathbf Z$-modules, identiques aux
groupes commutatifs. Soit
\[
  T:\mathbf{Ab}_A\longrightarrow\mathbf{Ab}
\]
un foncteur covariant additif, et soit $M$ un $(A,A)$-bimodule; $T(M)$ est
alors muni de façon naturelle d'une structure de $A$-module à droite. En
effet, pour tout $a\in A$, notons $h_{a,M}$ (ou simplement $h_a$)
l'endomorphisme $x\mapsto xa$ du $A$-module à gauche $M$. Par hypothèse,
$T(h_a)$ est un endomorphisme du $\mathbf Z$-module $T(M)$; en outre, comme
$T$ est un foncteur covariant additif, on a, pour $a\in A$, $b\in A$,
\[
  T(h_{ab})
  =
  T(h_b\circ h_a)
  =
  T(h_b)\circ T(h_a),
  \qquad\text{et}\qquad
  T(h_{a+b})
  =
  T(h_a+h_b)
  =
  T(h_a)+T(h_b);
\]
cela prouve que l'application
\[
  (a,y)\longmapsto T(h_a)(y)
\]
est une loi externe de $A$-module à droite sur $T(M)$. En particulier,
$T(A_s)$ est un $A$-module à droite.
\end{env}

\begin{env}[7.1.2]
Lorsque $A$ est un anneau commutatif, il résulte de (7.1.1) que pour tout
$A$-module $M$, $T(M)$ est naturellement muni d'une structure de $A$-module;
en outre, si $u:M\to N$ est un homomorphisme de $A$-modules, on a, pour tout
$a\in A$,
\[
  u\circ h_{a,M}=h_{a,N}\circ u,
\]
d'où
\[
  T(u)\circ T(h_{a,M})
  =
  T(h_{a,N})\circ T(u),
\]
ce qui prouve que $T(u):T(M)\to T(N)$ est un homomorphisme de $A$-modules; on
voit donc que $T$ peut être considéré comme un foncteur covariant additif de
la catégorie $\mathbf{Ab}_A$ dans elle-même. De façon précise, on a ainsi
défini une équivalence canonique entre la catégorie des foncteurs covariants
additifs $\mathbf{Ab}_A\to\mathbf{Ab}$ et la catégorie des foncteurs
covariants $A$-linéaires
\[
  T:\mathbf{Ab}_A\longrightarrow\mathbf{Ab}_A,
\]
c'est-à-dire tels que
\[
  T(h_{a,M})=h_{a,T(M)}
  \qquad\text{pour tout }a\in A.
\]
Comme le foncteur d'inclusion
\[
  I:\mathbf{Ab}_A\longrightarrow\mathbf{Ab},
\]
faisant correspondre à tout $A$-module le $\mathbf Z$-module sous-jacent, est
exact et fidèle, les propriétés d'exactitude des deux foncteurs associés par
l'équivalence précédente sont les mêmes.
\end{env}

\begin{env}[7.1.3]
L'anneau $A$ étant toujours supposé commutatif, soit $B$ une $A$-algèbre (non
nécessairement commutative), et soit $\rho:A\to B$ l'homomorphisme d'anneaux
correspondant à cette structure d'algèbre; cet homomorphisme définit un
foncteur covariant
\oldpage[III]{176}
additif
\[
  \rho_*:M\longmapsto M_{[\rho]}
\]
de la catégorie $\mathbf{Ab}_B$ des $B$-modules à gauche dans la catégorie
$\mathbf{Ab}_A$ des $A$-modules. Par composition, on en déduit un foncteur
\[
  T_{(B)}:
  \mathbf{Ab}_B
  \xrightarrow{\rho_*}
  \mathbf{Ab}_A
  \xrightarrow{T}
  \mathbf{Ab},
\]
évidemment covariant et additif, que nous noterons aussi $T^{(B)}$ (pour des
raisons typographiques) ou $T\otimes_A B$, et que nous dirons obtenu à partir
de $T$ par extension des scalaires de $A$ à $B$. Bien entendu, si $B$ est
commutative, on peut considérer $T_{(B)}$ comme un foncteur de
$\mathbf{Ab}_B$ dans elle-même (7.1.2). Lorsque $B$ est commutative, et $C$
une $B$-algèbre, on voit aussitôt que
\[
  T_{(C)}=(T_{(B)})_{(C)};
\]
il est immédiat que l'extension des scalaires est fonctorielle et additive en
$T$; en outre, lorsque $T$ commute aux limites inductives ou aux sommes
directes (resp. est exact à gauche, exact à droite, exact), il en est de même
de $T_{(B)}$: en effet, $\rho_*$ est exact et commute aux limites inductives
et aux sommes directes.
\end{env}

\begin{env}[7.1.4]
Supposons toujours $A$ commutatif, et soit
$T:\mathbf{Ab}_A\to\mathbf{Ab}_A$ un foncteur covariant additif
$A$-linéaire, commutant aux limites inductives. Alors, pour toute partie
multiplicative $S$ de $A$ et tout $A$-module $M$, on a un isomorphisme
canonique fonctoriel de $A$-modules
\[
  T(S^{-1}M)\xrightarrow{\sim}S^{-1}T(M).
  \tag{7.1.4.1}\label{III.7.1.4.1.fr}
\]
Supposons en effet d'abord que $S$ soit l'ensemble des puissances $f^n$
$(n\geqslant0)$ d'un élément $f\in A$. On sait alors que
\[
  M_f=\varinjlim M_n,
\]
où $(M_n,\varphi_{nm})$ est le système inductif des $A$-modules $M_n=M$,
avec
\[
  \varphi_{nm}:z\longmapsto f^{\,n-m}z
  \qquad(0_{\mathrm I}, 1.6.1);
\]
d'où dans ce cas l'isomorphisme (7.1.4.1) en vertu de l'hypothèse sur $T$; si
ensuite $S$ est quelconque, $S^{-1}M$ est limite inductive des $M_f$, pour
$f\in S$ ($0_{\mathrm I}$, 1.4.5) et l'on conclut de la même façon. En outre,
la fonctorialité de l'isomorphisme (7.1.4.1) montre que c'est un isomorphisme
de $S^{-1}A$-modules, et qu'on peut donc écrire, à un isomorphisme canonique
près
\[
  T_{(S^{-1}A)}(S^{-1}M)
  =
  S^{-1}T(M)
  =
  T(S^{-1}M).
  \tag{7.1.4.2}\label{III.7.1.4.2.fr}
\]
Lorsque $S=A-\mathfrak p$ est le complémentaire d'un idéal premier
$\mathfrak p$ de $A$, on écrit $T_{\mathfrak p}$ au lieu de
$T_{(A_{\mathfrak p})}$.
\end{env}

\begin{proposition}[7.1.5]
Sous les hypothèses de (7.1.4), si $T_{\mathfrak m}$ est exact à gauche
(resp. exact à droite, exact) pour tout idéal maximal $\mathfrak m$ de $A$,
$T$ est exact à gauche (resp. exact à droite, exact).
\end{proposition}

\begin{proof}
On sait en effet que lorsque deux sous-modules $N$, $P$ d'un $A$-module $M$
sont tels que $N_{\mathfrak m}=P_{\mathfrak m}$ pour tout idéal maximal
$\mathfrak m$ de $A$, on a $N=P$ (Bourbaki, \emph{Alg. comm.}, chap. II,
\S~3, no~3, th.~1).
\end{proof}

\subsection{Caractérisations du foncteur produit tensoriel}
\label{subsection:III.7.2.fr}

\begin{env}[7.2.1]
Soient $A$ un anneau (non nécessairement commutatif), $M$ (resp. $N$) un
$A$-module à gauche (resp. à droite), $P$ un $\mathbf Z$-module. Rappelons
que la donnée d'un $\mathbf Z$-homomorphisme
\[
  v:N\otimes_A M\longrightarrow P
\]
équivaut à la donnée d'une application $\mathbf Z$-bilinéaire
\[
  u:N\times M\longrightarrow P
\]
telle que
\[
  u(ta,x)=u(t,ax)
  \qquad\text{pour }a\in A,\ t\in N,\ x\in M,
\]
les deux applications étant reliées par $v(t\otimes x)=u(t,x)$. D'autre part,
la donnée de $u$ équivaut à celle d'un
\oldpage[III]{177}
$\mathbf Z$-homomorphisme $x\mapsto f_x$ de $M$ dans
$\operatorname{Hom}_{\mathbf Z}(N,P)$ tel que
\[
  f_{ax}(t)=f_x(ta)
  \qquad\text{pour }a\in A,\ t\in N,\ x\in M,
\]
les deux applications étant reliées par $u(t,x)=f_x(t)$.
\end{env}

\begin{env}[7.2.2]
Soit $T:\mathbf{Ab}_A\to\mathbf{Ab}$ un foncteur covariant additif. Nous
allons définir pour tout $A$-module à gauche $M$, un homomorphisme canonique
fonctoriel en $M$, de $\mathbf Z$-modules
\[
  t_M:T(A_s)\otimes_A M\longrightarrow T(M).
  \tag{7.2.2.1}\label{III.7.2.2.1.fr}
\]
Il suffira pour cela, en vertu de (7.2.1), de définir un
$\mathbf Z$-homomorphisme $x\mapsto t'_M(x)$ de $M$ dans
$\operatorname{Hom}_{\mathbf Z}(T(A_s),T(M))$, tel que l'on ait
\[
  t'_M(ax)(y)=t'_M(x)(ya)
  \qquad\text{pour }a\in A,\ x\in M,\ y\in T(A_s).
\]
Notons pour cela que
$\operatorname{Hom}_{\mathbf Z}(T(A_s),T(M))$ est canoniquement muni d'une
structure de $A$-module à gauche provenant de la structure de $A$-module à
droite de $T(A_s)$, la loi externe étant telle que si
$a\in A$, $v\in\operatorname{Hom}_{\mathbf Z}(T(A_s),T(M))$,
\[
  (a\cdot v)(y)=v(ya)
  \qquad\text{pour }y\in T(A_s).
\]
Cela étant, nous définirons $t'_M$ comme composé des deux homomorphismes
canoniques
\[
  M\xrightarrow{\sim}\operatorname{Hom}_A(A_s,M)
  \xrightarrow{T}
  \operatorname{Hom}_{\mathbf Z}(T(A_s),T(M)),
\]
la seconde flèche étant l'application $u\mapsto T(u)$, la première
l'isomorphisme canonique de $A$-modules $x\mapsto\theta_x$ tel que
$\theta_x(\xi)=\xi x$ pour $\xi\in A$, $x\in M$. On a
$\theta_{ax}=\theta_x\circ h_a$, donc
\[
  T(\theta_{ax})
  =
  T(\theta_x\circ h_a)
  =
  T(\theta_x)\circ T(h_a)
\]
et par suite, pour $y\in T(A_s)$,
\[
  T(\theta_{ax})(y)
  =
  T(\theta_x)(T(h_a)(y))
  =
  T(\theta_x)(ya)
\]
par définition de la loi externe sur $T(A_s)$, ce qui prouve l'existence de
$t_M$; il est immédiat de vérifier que cet homomorphisme est fonctoriel en
$M$, c'est-à-dire que pour tout homomorphisme $w:M\to M'$ de $A$-modules à
gauche, le diagramme
\[
\xymatrix@C=30pt@R=30pt{
  T(A_s)\otimes_A M
  \ar[r]^-{t_M}\ar[d]_{1\otimes w}
  &
  T(M)\ar[d]^{T(w)}
  \\
  T(A_s)\otimes_A M'
  \ar[r]_-{t_{M'}}
  &
  T(M')
}
\tag{7.2.2.2}\label{III.7.2.2.2.fr}
\]
est commutatif.

La fonctorialité de l'homomorphisme (7.2.2.1) montre que lorsque $A$ est
commutatif, c'est un homomorphisme de $A$-modules (cf. (7.1.2)).
\end{env}

\begin{env}[7.2.3]
Lorsque $A$ est commutatif, on peut plus généralement définir un
homomorphisme canonique de $A$-modules
\[
  T(N)\otimes_A M\longrightarrow T(N\otimes_A M)
  \tag{7.2.3.1}\label{III.7.2.3.1.fr}
\]
pour tout $A$-module $N$; il suffit dans la construction de (7.2.2) de
remplacer l'homomorphisme $\theta_x$ par l'homomorphisme de $A$-modules
\[
  N\longrightarrow N\otimes_A M
\]
qui à tout $y\in N$ fait correspondre $y\otimes x$. Il est immédiat que cet
homomorphisme est fonctoriel en $M$ et $N$.

\oldpage[III]{178}
En particulier, si $B$ est une $A$-algèbre (non nécessairement commutative),
on a un homomorphisme fonctoriel en $M$
\[
  (T(M))_{(B)}
  =
  T(M)\otimes_A B
  \longrightarrow
  T(M\otimes_A B)
  =
  T_{(B)}(M_{(B)})
  \tag{7.2.3.2}\label{III.7.2.3.2.fr}
\]
qui, en vertu de la fonctorialité de (7.2.3.1) en $M$, est un homomorphisme
de $B$-modules.

On a en outre le diagramme commutatif
\[
\xymatrix@C=30pt@R=30pt{
  T(A)\otimes_A M
  \ar[r]^-{t_M}\ar[d]
  &
  T(M)\ar[d]
  \\
  T_{(B)}(B_s)\otimes_B M_{(B)}
  \ar[r]_-{t_{M_{(B)}}}
  &
  T_{(B)}(M_{(B)})
}
\tag{7.2.3.3}\label{III.7.2.3.3.fr}
\]
où la flèche verticale de droite est l'homomorphisme composé
\[
  T(M)
  \longrightarrow
  T(M)\otimes_A B
  \longrightarrow
  T(M\otimes_A B)
  =
  T_{(B)}(M_{(B)})
\]
de (7.2.3.2) et de l'homomorphisme canonique; quant à la flèche verticale de
gauche de (7.2.3.3), c'est l'homomorphisme
\[
  T(A)\otimes_A M
  \longrightarrow
  T_{(B)}(B_s)\otimes_B(B\otimes_A M)
  =
  T_{(B)}(B_s)\otimes_A M
\]
où
\[
  T(A)\longrightarrow T_{(B)}(B_s)=T(B)
\]
est $T(\rho)$, $\rho$ étant considéré comme homomorphisme de $A$-modules
$A\to B_{[\rho]}$.
\end{env}

\begin{lemma}[7.2.4]
Si $T$ est un foncteur covariant additif de $\mathbf{Ab}_A$ dans
$\mathbf{Ab}$, commutant avec les sommes directes, l'homomorphisme canonique
$t_L$ (7.2.2.1) est un isomorphisme pour tout $A$-module libre $L$.
\end{lemma}

\begin{proof}
En effet, on a
\[
  L=\bigoplus_{\alpha\in I}L_\alpha
\]
où $L_\alpha$ est isomorphe à $A_s$ pour tout $\alpha\in I$; la définition de
$t_M$ donnée dans (7.2.2) montre que
\[
  t_L=\bigoplus_{\alpha\in I}t_{L_\alpha},
\]
puisque
\[
  T:
  \operatorname{Hom}_A(A_s,L)
  \longrightarrow
  \operatorname{Hom}_{\mathbf Z}(T(A_s),T(L))
\]
est somme directe des applications $\mathbf Z$-linéaires
\[
  T_\alpha:
  \operatorname{Hom}_A(A_s,L_\alpha)
  \longrightarrow
  \operatorname{Hom}_{\mathbf Z}(T(A_s),T(L_\alpha))
\]
en vertu de l'hypothèse sur $T$. On est donc ramené à prouver le lemme pour
$L=A_s$; mais $t_L$ n'est autre alors que l'isomorphisme canonique
\[
  T(A_s)\otimes_A A_s\xrightarrow{\sim}T(A_s)
\]
valable pour tout $A$-module à droite.
\end{proof}

\begin{proposition}[7.2.5]
Soit $T$ un foncteur covariant additif de $\mathbf{Ab}_A$ dans
$\mathbf{Ab}$, commutant aux sommes directes. Les conditions suivantes sont
équivalentes :
\begin{enumerate}
  \item[a)] $T$ est exact à droite.
  \item[b)] L'homomorphisme canonique $t_M$ (7.2.2.1) est un isomorphisme pour
    tout $A$-module à gauche $M$.
  \item[$b'$)] $T$ est semi-exact et l'homomorphisme $t_M$ est surjectif pour
    tout $A$-module à gauche $M$.
  \item[c)] $T$ est isomorphe à un foncteur en $M$ de la forme
    $N\otimes_A M$, où $N$ est un $A$-module à droite.
\end{enumerate}
\end{proposition}

\begin{proof}
Il est clair que $b)$ implique $c)$ et que $c)$ implique $a)$; montrons que
$a)$ implique $b)$.
\oldpage[III]{179}
Posons
\[
  T'(M)=T(A_s)\otimes_A M
\]
pour tout $A$-module à gauche $M$. Il existe une suite exacte
\[
  L'\longrightarrow L\longrightarrow M\longrightarrow0,
\]
où $L$ et $L'$ sont deux $A$-modules à gauche libres; comme $T$ et $T'$ sont
exacts à droite, on a donc le diagramme commutatif
\[
\xymatrix@C=22pt@R=25pt{
  T'(L')\ar[r]\ar[d]^{t_{L'}}
  &
  T'(L)\ar[r]\ar[d]^{t_L}
  &
  T'(M)\ar[r]\ar[d]^{t_M}
  &
  0
  \\
  T(L')\ar[r]
  &
  T(L)\ar[r]
  &
  T(M)\ar[r]
  &
  0
}
\]
où les deux lignes sont exactes; comme $t_L$ et $t_{L'}$ sont des
isomorphismes en vertu de (7.2.4), il en est de même de $t_M$ par le lemme des
cinq. Enfin, il est clair que $b)$ entraîne $b')$. Pour montrer que $b')$
entraîne $a)$, il suffit de prouver le lemme suivant.
\end{proof}

\begin{lemma}[7.2.5.1]
Soient $\mathcal K$, $\mathcal K'$ deux catégories abéliennes, $F$, $G$ deux
foncteurs covariants additifs de $\mathcal K$ dans $\mathcal K'$,
$f:F\to G$ un morphisme fonctoriel (T, I, 1.2) tel que, pour tout objet $E$
de la catégorie $\mathcal K$,
\[
  f_E:F(E)\longrightarrow G(E)
\]
soit un épimorphisme. Alors, si $F$ est exact à droite et $G$ semi-exact,
$G$ est exact à droite.
\end{lemma}

\begin{proof}
Tout revient en effet à démontrer que pour tout épimorphisme
$v:E'\to E$ dans $\mathcal K$,
\[
  G(v):G(E')\longrightarrow G(E)
\]
est un épimorphisme; or, on a le diagramme commutatif
\[
\xymatrix@C=35pt@R=30pt{
  F(E')\ar[r]^{F(v)}\ar[d]_{f_{E'}}
  &
  F(E)\ar[d]^{f_E}
  \\
  G(E')\ar[r]_{G(v)}
  &
  G(E)
}
\]
dans lequel $F(v)$, $f_{E'}$ et $f_E$ sont des épimorphismes; il en est donc
de même de $G(v)$.
\end{proof}

\begin{remark}[7.2.6]
Pour tout $A$-module à droite $N$, posons
\[
  T_N(M)=N\otimes_A M
\]
pour tout $A$-module à gauche $M$, de sorte que $T_N$ est un foncteur
covariant additif de $\mathbf{Ab}_A$ dans $\mathbf{Ab}$, exact à droite et
commutant aux sommes directes. Si on identifie canoniquement $T_N(A_s)$ à
$N$, on vérifie aussitôt que l'homomorphisme correspondant (7.2.2.1) devient
l'identité. On en conclut que le $A$-module à droite $N$ de l'énoncé de
(7.2.5, $c)$) est déterminé à un isomorphisme unique près et est
canoniquement isomorphe à $T(A_s)$. On peut encore dire que les morphismes
fonctoriels
\[
  T\longmapsto T(A_s),
  \qquad
  N\longmapsto T_N
\]
constituent une équivalence (T, I, 1.2) de la catégorie des $A$-modules à
droite et de la catégorie des foncteurs covariants additifs
$\mathbf{Ab}_A\to\mathbf{Ab}$ qui sont exacts à droite et commutent aux
sommes directes.
\end{remark}

\begin{proposition}[7.2.7]
Soient $A$ un anneau artinien à gauche, dont le quotient par son radical
$\mathfrak m$ est un corps $k$. Soit $T$ un foncteur covariant additif de
$\mathbf{Ab}_A$ dans $\mathbf{Ab}$, commutant aux sommes directes. Les
conditions de (7.2.5) sont alors aussi équivalentes à
\begin{enumerate}
  \item[d)] $T$ est semi-exact et l'homomorphisme
    \[
      T(\varepsilon):T(A_s)\longrightarrow T(k)
    \]
    déduit de l'homomorphisme canonique $\varepsilon:A_s\to k$ est surjectif.
\end{enumerate}
\end{proposition}

\begin{proof}
\oldpage[III]{180}
Il est clair que la condition $b')$ de (7.2.5) entraîne $d)$; prouvons que
$d)$ entraîne $b')$.

Il existe un entier $n$ tel que $\mathfrak m^n=0$; posons, pour tout
$A$-module $M$, $M_h=\mathfrak m^hM$; nous prouverons par récurrence
descendante sur $h$ que $t_{M_h}$ est surjectif. La proposition est évidente
pour $h=n$; pour $h<n$, on a une suite exacte
\[
  0\longrightarrow M_{h+1}
  \longrightarrow M_h
  \longrightarrow M_h/M_{h+1}
  \longrightarrow 0
\]
et l'hypothèse de récurrence entraîne que $t_{M_{h+1}}$ est surjectif.
D'autre part, $M_h/M_{h+1}$ est annulé par $\mathfrak m$ et est donc un
$(A/\mathfrak m)$-module, autrement dit est somme directe de $A$-modules
isomorphes à $k$. Pour prouver que $t_{M_h/M_{h+1}}$ est surjectif, il suffit
donc de prouver que $t_k$ l'est, puisque $T$ commute aux sommes directes. Or,
en vertu de la commutativité du diagramme
\[
\xymatrix@C=35pt@R=30pt{
  T(A_s)\otimes_A A_s
  \ar[r]^-{t_{A_s}}\ar[d]_{1\otimes\varepsilon}
  &
  T(A_s)\ar[d]^{T(\varepsilon)}
  \\
  T(A_s)\otimes_A k
  \ar[r]_-{t_k}
  &
  T(k)
}
\]
et de (7.2.4), l'hypothèse $d)$ entraîne que $t_k$ est bien surjectif. Pour
terminer la démonstration, il suffira de prouver que si l'on a une suite
exacte
\[
  0\longrightarrow M'
  \longrightarrow M
  \xrightarrow{v}M''
  \longrightarrow 0
\]
de $A$-modules, telle que $t_{M'}$ et $t_{M''}$ sont surjectifs, alors $t_M$
est surjectif. Or, on a un diagramme commutatif
\[
\xymatrix@C=27pt@R=30pt{
  T'(M')\ar[r]\ar[d]_{t_{M'}}
  &
  T'(M)\ar[r]\ar[d]_{t_M}
  &
  T'(M'')\ar[r]\ar[d]_{t_{M''}}
  &
  0\ar[d]
  \\
  T(M')\ar[r]
  &
  T(M)\ar[r]_-{T(v)}
  &
  T(M'')\ar[r]
  &
  \operatorname{Coker}(T(v))
}
\]
dans lequel les deux lignes sont exactes, en vertu de l'hypothèse que $T$ est
semi-exact. Comme par l'hypothèse de récurrence $t_{M'}$ et $t_{M''}$ sont
des épimorphismes et que la dernière flèche verticale est un monomorphisme,
le lemme des cinq (M, I, 1.1) montre que $t_M$ est un épimorphisme.
\end{proof}

\subsection{Critères d'exactitude des foncteurs homologiques de modules}
\label{subsection:III.7.3.fr}

\begin{proposition}[7.3.1]
Soient $A$ un anneau (non nécessairement commutatif), $T_\bullet$ un foncteur
homologique covariant (T, II, 2.1) de la catégorie $\mathbf{Ab}_A$ dans la
catégorie $\mathbf{Ab}$, commutant aux sommes directes. Soit $p$ un entier tel
que $T_p$ et $T_{p-1}$ soient définis. Les conditions suivantes sont
équivalentes :
\begin{enumerate}
  \item[a)] $T_p$ est exact à droite.
  \item[b)] $T_{p-1}$ est exact à gauche.
  \oldpage[III]{181}
  \item[c)] Pour tout $A$-module à gauche $M$, l'homomorphisme fonctoriel
    canonique (7.2.2.1)
    \[
      T_p(A_s)\otimes_A M\longrightarrow T_p(M)
      \tag{7.3.1.1}\label{III.7.3.1.1.fr}
    \]
    est un isomorphisme.
  \item[d)] Pour tout $A$-module à gauche $M$, l'homomorphisme (7.3.1.1) est
    un épimorphisme.
  \item[e)] $T_p$ est isomorphe à un foncteur
    $M\mapsto N\otimes_A M$, où $N$ est un $A$-module à droite.
\end{enumerate}
Si en outre les conditions de (7.2.7) sur $A$ et $\mathfrak m$ sont vérifiées,
les conditions précédentes sont aussi équivalentes à
\begin{enumerate}
  \item[f)] L'homomorphisme canonique
    \[
      T_p(\varepsilon):T_p(A_s)\longrightarrow T_p(k)
    \]
    est un épimorphisme.
\end{enumerate}
\end{proposition}

\begin{proof}
Comme par définition d'un foncteur homologique, $T_i$ est semi-exact pour tous
les $i$ tels que $T_i$ soit défini, et que, pour toute suite exacte
\[
  0\longrightarrow M'\xrightarrow{u}M\xrightarrow{v}M''\longrightarrow0,
\]
on a
\[
  \operatorname{Ker}(T_{i-1}(u))
  =
  \operatorname{Coker}(T_i(v)),
\]
il est clair que $a)$ et $b)$ sont équivalentes et les autres assertions
résultent trivialement de (7.2.5) et (7.2.7).
\end{proof}

\begin{corollary}[7.3.2]
Soit $A$ un anneau commutatif. Avec les notations de (7.3.1), supposons $T_p$
exact à droite. Si $f\in A$ n'appartient à l'annulateur d'aucun élément
$\ne0$ d'un $A$-module $M$, alors $f$ n'appartient à l'annulateur d'aucun
élément $\ne0$ de $T_{p-1}(M)$. En particulier, si $A$ est intègre, le
$A$-module $T_{p-1}(A)$ est sans torsion.
\end{corollary}

\begin{proof}
En effet, si $h_f$ désigne l'homothétie $x\mapsto fx$ de $M$, l'hypothèse
signifie que $h_f$ est injectif; il en est donc de même de $T_{p-1}(h_f)$ par
la condition $b)$ de (7.3.1).
\end{proof}

\begin{proposition}[7.3.3]
Soient $A$ un anneau, $T_\bullet$ un foncteur homologique covariant de
$\mathbf{Ab}_A$ dans $\mathbf{Ab}$, commutant aux sommes directes. Soit $p$ un
entier tel que $T_{p-1}$, $T_p$ et $T_{p+1}$ soient définis. Les conditions
suivantes sont équivalentes :
\begin{enumerate}
  \item[a)] $T_p$ est exact.
  \item[b)] $T_{p+1}$ et $T_p$ sont exacts à droite.
  \item[c)] $T_p$ et $T_{p-1}$ sont exacts à gauche.
  \item[d)] $T_{p+1}$ est exact à droite et $T_{p-1}$ est exact à gauche.
  \item[e)] Pour tout $A$-module $M$, les homomorphismes canoniques
    \[
      T_i(A_s)\otimes_A M\longrightarrow T_i(M)
      \tag{7.3.3.1}\label{III.7.3.3.1.fr}
    \]
    sont des isomorphismes pour $i=p$ et $i=p+1$.
  \item[$e'$)] Pour tout $A$-module $M$, les homomorphismes canoniques
    (7.3.3.1) sont des épimorphismes pour $i=p$ et $i=p+1$.
  \item[f)] Pour tout $A$-module $M$, l'homomorphisme (7.3.3.1) est un
    isomorphisme pour $i=p$ et $T_p(A_s)$ est un $A$-module à droite plat.
  \item[$f'$)] Pour tout $A$-module $M$, l'homomorphisme (7.3.3.1) est un
    épimorphisme pour $i=p$ et $T_p(A_s)$ est un $A$-module à droite plat.
\end{enumerate}
\end{proposition}

\begin{proof}
L'équivalence des conditions $a)$, $b)$, $c)$, $d)$ résulte de l'équivalence
des conditions $a)$ et $b)$ de (7.3.1). L'équivalence de $b)$, $e)$ et $e')$
résulte de l'équivalence de $a)$, $c)$ et $d)$ dans (7.3.1). Enfin, dire que
$T_p(A_s)$ est plat signifie que le foncteur
\[
  M\mapsto T_p(A_s)\otimes_A M
\]
est exact à gauche; l'équivalence de $a)$, $f)$ et $f')$ résulte encore de
l'équivalence de $a)$, $c)$, $d)$ dans (7.3.1).
\end{proof}

\oldpage[III]{182}

\begin{corollary}[7.3.4]
Supposons $A$ commutatif, $T_p$ exact et supposons en outre que $T_p(A)$ soit
un $A$-module de présentation finie. Alors la fonction
\[
  x\mapsto\operatorname{rang}_{\kappa(x)}(T_p(\kappa(x)))
\]
est localement constante dans $X=\operatorname{Spec}(A)$, donc constante si
$\operatorname{Spec}(A)$ est connexe.
\end{corollary}

\begin{proof}
En effet, comme $T_p(A)$ est un $A$-module plat en vertu de (7.3.3, $f)$), il
est projectif de type fini, et $\bigl(T_p(A)\bigr)^{\sim}$ est donc un
$\mathscr O_X$-Module localement libre (Bourbaki, \emph{Alg. comm.}, chap. II,
\S 5, n\textsuperscript{o} 2, th. 1); on a en outre
\[
  T_p(\kappa(x))=T_p(A)\otimes_A\kappa(x)
\]
(7.3.3, $e)$) et l'on sait que le rang au point $x$ du $A$-module $T_p(A)$ est
localement constant (\emph{loc. cit.}), d'où le corollaire.
\end{proof}

\begin{proposition}[7.3.5]
Supposons que $A$ soit un anneau artinien à gauche dont le quotient par son
radical $\mathfrak m$ est un corps $k$. Alors les conditions de (7.3.3) sont
encore équivalentes à chacune des suivantes :
\begin{enumerate}
  \item[g)] L'homomorphisme canonique
    \[
      T_i(\varepsilon):T_i(A_s)\longrightarrow T_i(k)
    \]
    est un épimorphisme pour $i=p$ et $i=p+1$.
  \item[h)] $T_p(\varepsilon)$ est un épimorphisme et $T_p(A_s)$ est un
    $A$-module à droite plat (ou, ce qui revient au même (Bourbaki,
    \emph{Alg. comm.}, chap. II, \S 3, n\textsuperscript{o} 2, cor. 2 de la
    prop. 5) un $A$-module libre).
\end{enumerate}
Supposons en outre que $A$ soit commutatif et le $A$-module $T_p(k)$ de
longueur finie $d$. Alors les conditions précédentes équivalent aussi à
chacune des suivantes :
\begin{enumerate}
  \item[i)] Pour tout $A$-module $M$ de longueur finie, on a
    \[
      \operatorname{long}(T_p(M))=d\mathbin{.}\operatorname{long}(M).
      \tag{7.3.5.1}\label{III.7.3.5.1.fr}
    \]
  \item[j)] On a
    \[
      \operatorname{long}(T_p(A))=d\mathbin{.}\operatorname{long}(A).
      \tag{7.3.5.2}\label{III.7.3.5.2.fr}
    \]
\end{enumerate}
\end{proposition}

L'équivalence de $g)$ et $h)$ avec les conditions de (7.3.3) découle aussitôt
de (7.2.7). Pour démontrer les autres assertions, nous utiliserons le lemme
suivant :

\begin{lemma}[7.3.5.3]
Soient $\mathcal K$, $\mathcal K'$ deux catégories abéliennes,
$F:\mathcal K\to\mathcal K'$ un foncteur
covariant additif; on suppose que $F$ soit semi-exact, et que, pour tout objet
simple $S$ de $\mathcal K$, $F(S)$ soit un objet de longueur finie dans
$\mathcal K'$. Alors, pour tout objet $E$ de longueur finie dans $\mathcal K$,
$F(E)$ est de longueur finie dans $\mathcal K'$. Pour toute suite exacte
\[
  0\longrightarrow E'\xrightarrow{u}E\xrightarrow{v}E''\longrightarrow0
\]
d'objets de longueur finie dans $\mathcal K$, on a
\[
  \operatorname{long}F(E)
  \leq
  \operatorname{long}F(E')+\operatorname{long}F(E'')
  \tag{7.3.5.4}\label{III.7.3.5.4.fr}
\]
et pour que les deux membres de (7.3.5.4) soient égaux, il faut et il suffit
que la suite
\[
  0\longrightarrow F(E')\longrightarrow F(E)\longrightarrow F(E'')
  \longrightarrow0
\]
soit exacte.
\end{lemma}

\begin{proof}
En effet, la suite
\[
  F(E')\xrightarrow{F(u)}F(E)\xrightarrow{F(v)}F(E'')
\]
est exacte par hypothèse; si l'on suppose $F(E')$ et $F(E'')$ de longueur
finie, il en est de même de $\operatorname{Im}(F(u))$ et de
$\operatorname{Im}(F(v))$ et comme
$\operatorname{Ker}(F(v))=\operatorname{Im}(F(u))$, $F(E)$ est de longueur
finie et l'on a
\[
  \operatorname{long}F(E)
  =
  \operatorname{long}\operatorname{Im}(F(u))
  +
  \operatorname{long}\operatorname{Im}(F(v))
  \leq
  \operatorname{long}F(E')+\operatorname{long}F(E'').
  \tag{7.3.5.5}\label{III.7.3.5.5.fr}
\]
\oldpage[III]{183}
Par récurrence sur la longueur de $E$, cela prouve déjà la première assertion;
en outre, les deux membres de (7.3.5.5) ne peuvent être égaux que si
\[
  \operatorname{long}\operatorname{Im}(F(u))=\operatorname{long}F(E')
\]
(ce qui équivaut à
$\operatorname{long}\operatorname{Ker}(F(u))=0$, ou
$\operatorname{Ker}(F(u))=0$) et
\[
  \operatorname{long}\operatorname{Im}(F(v))=\operatorname{long}F(E'')
\]
(ce qui équivaut à
$\operatorname{long}\operatorname{Coker}(F(v))=0$, ou
$\operatorname{Coker}(F(v))=0$).
\end{proof}

Notons maintenant que si $M$ est un $A$-module de longueur finie ($A$ étant
commutatif), les quotients d'une suite de Jordan-Hölder de $M$ sont
nécessairement isomorphes au $A$-module $k$, donc on déduit de (7.3.5.4), par
récurrence sur la longueur de $M$
\[
  \operatorname{long}T_p(M)\leq d\mathbin{.}\operatorname{long}(M).
  \tag{7.3.5.6}\label{III.7.3.5.6.fr}
\]
En outre, il résulte de (7.3.5.3) que si $T_p$ est exact, on a l'égalité
(7.3.5.1); donc la condition $a)$ de (7.3.3) entraîne $i)$; il est clair que
$i)$ entraîne $j)$, et il reste à prouver le

\begin{lemma}[7.3.5.7]
La relation
\[
  \operatorname{long}T_p(A)=d\mathbin{.}\operatorname{long}A
\]
entraîne que $T_p(\varepsilon)$ est un épimorphisme et que $T_p(A)$ est un
$A$-module plat.
\end{lemma}

\begin{proof}
En effet, partant de la suite exacte
\[
  0\longrightarrow\mathfrak m\longrightarrow A\longrightarrow k
  \longrightarrow0,
\]
il résulte de (7.3.5.4) et (7.3.5.6) que l'on a
\[
  \operatorname{long}T_p(A)
  \leq
  \operatorname{long}T_p(\mathfrak m)+\operatorname{long}T_p(k)
  \leq
  d\bigl(\operatorname{long}\mathfrak m+\operatorname{long}k\bigr)
  =d\mathbin{.}\operatorname{long}A
\]
et que l'égalité ne peut avoir lieu (7.3.5.3) que si la suite
\[
  0\longrightarrow T_p(\mathfrak m)\longrightarrow T_p(A)
  \longrightarrow T_p(k)\longrightarrow0
  \tag{7.3.5.8}\label{III.7.3.5.8.fr}
\]
est exacte. En vertu de (7.2.7) et (7.2.5), $T_p$ est isomorphe à un foncteur
$M\mapsto N\otimes_A M$, et l'exactitude de la suite (7.3.5.8) montre, en vertu
de la suite exacte des $\operatorname{Tor}$, que l'on a
$\operatorname{Tor}_1^A(N,k)=0$. On en conclut que $N=T_p(A)$ est un
$A$-module plat (0, 10.1.3).
\end{proof}

\begin{lemma}[7.3.6]
Soient $A$ un anneau, $T_\bullet$ un foncteur homologique covariant de
$\mathbf{Ab}_A$ dans $\mathbf{Ab}$, commutant aux sommes directes. Supposons
$T_p$ et $T_{p+1}$ définis, et $T_p$ exact à gauche. Pour que $T_{p+1}$ soit
exact, il faut et il suffit que $T_{p+1}(A_s)$ soit un $A$-module à droite
plat.
\end{lemma}

\begin{proof}
En effet, on sait d'après (7.3.1) que l'homomorphisme canonique
\[
  T_{p+1}(A_s)\otimes_A M\longrightarrow T_{p+1}(M)
\]
est un isomorphisme de foncteurs; il suffit d'appliquer la définition d'un
$A$-module plat.
\end{proof}

\begin{proposition}[7.3.7]
Soient $A$ un anneau, $T_\bullet$ un foncteur homologique covariant de
$\mathbf{Ab}_A$ dans $\mathbf{Ab}$, commutant aux sommes directes. On suppose
qu'il existe $i_0$ tel que $T_i$ soit exact pour $i\leq i_0$. Alors, pour tout
entier $p>i_0$, les conditions suivantes sont équivalentes :
\begin{enumerate}
  \item[a)] $T_q$ est exact pour $q\leq p$.
  \item[b)] $T_q(A_s)$ est un $A$-module à droite plat pour $q\leq p$.
  \item[c)] Pour tout $A$-module $M$, l'homomorphisme canonique
    \[
      T_q(A_s)\otimes_A M\longrightarrow T_q(M)
    \]
    est surjectif pour $q\leq p+1$.
\end{enumerate}
\end{proposition}

\begin{proof}
L'équivalence de $a)$ et $b)$ résulte de (7.3.6) par récurrence sur $q$,
puisque $T_{i_0}$ est exact par hypothèse; l'équivalence de $a)$ et $c)$
résulte de l'équivalence des conditions $a)$ et $e')$ dans (7.3.3).
\end{proof}

\oldpage[III]{184}

\begin{env}[7.3.8]
Si $A$ est un anneau commutatif, $B$ une $A$-algèbre (non nécessairement
commutative), $T_\bullet$ un foncteur homologique covariant de
$\mathbf{Ab}_A$ dans $\mathbf{Ab}$, il résulte des définitions (7.1.3) que le
foncteur de $\mathbf{Ab}_B$ dans $\mathbf{Ab}$ obtenu par extension des
scalaires de $A$ à $B$, et que nous noterons
\[
  T_\bullet^{(B)}=(T_i^{(B)}),
\]
est encore un foncteur homologique.
\end{env}

\begin{corollary}[7.3.9]
Supposons que $T_\bullet$ vérifie les conditions générales de (7.3.7) et
commute aux limites inductives, et en outre que $A$ soit un anneau intègre et
tous les $T_n(A)$ des $A$-modules de présentation finie. Alors, pour tout
entier $N$, il existe un $f\in A-\{0\}$ tel que le foncteur
\[
  T_p^{(A_f)}:\mathbf{Ab}_{A_f}\longrightarrow\mathbf{Ab}
\]
soit exact pour $p\leq N$.
\end{corollary}

\begin{proof}
Par hypothèse, $T_i$ est exact pour $i\leq i_0$, donc $T_i(A)$ est plat pour
ces valeurs de $i$. En vertu de (7.3.7, $b)$), il suffit de prendre $f$ tel que
\[
  T_p^{(A_f)}(A_f)=T_p(A_f)
\]
soit un $A_f$-module libre pour $i_0<p\leq N$. Or, on a
$T_p(A_f)=(T_p(A))_f$ puisque $T_p$ commute aux limites inductives (7.1.4). Si
$x$ est le point générique de $\operatorname{Spec}(A)$, $(T_p(A))_x$ est un
espace vectoriel de dimension finie sur le corps des fractions de $A$. Comme
chaque $T_p(A)$ est de présentation finie, il existe bien un $f$ ayant la
propriété voulue (Bourbaki, \emph{Alg. comm.}, chap. II, \S 5,
n\textsuperscript{o} 1, cor. de la prop. 2).

On notera que s'il n'y a qu'un nombre fini d'indices $i$ tels que $T_i\ne0$,
il existe $f\in A-\{0\}$ tel que tous les $T_p^{(A_f)}$ soient exacts.
\end{proof}

\begin{corollary}[7.3.10]
Supposons que $T_\bullet$ vérifie les conditions générales de (7.3.7) et
commute aux limites inductives, que $A$ soit commutatif et noethérien et les
$T_n(A)$ des $A$-modules de type fini. Alors, pour tout entier $N$, il existe
un ouvert dense $U$ de $\operatorname{Spec}(A)$ tel que, pour tout $p\leq N$,
la fonction
\[
  x\mapsto\operatorname{rang}_{\kappa(x)}(T_p(\kappa(x)))
\]
soit constante dans $U$.
\end{corollary}

\begin{proof}
Soit $\mathfrak p$ un idéal premier minimal de $A$; par hypothèse, l'anneau
$B=A/\mathfrak p$ est intègre et $\operatorname{Spec}(B)$ s'identifie à une
composante irréductible de l'espace topologique $\operatorname{Spec}(A)$. Nous
allons montrer, par récurrence sur $p\leq N$ qu'il existe
$f_p\in B-\{0\}$ tel que si on pose $B'=B_{f_p}$, $T_i^{(B')}$ soit exact et
les $T_i(B')$ des $B'$-modules de type fini pour $i\leq p$. La proposition est
vraie pour $p\leq i_0$ en vertu de l'hypothèse, en prenant $f_p=1$ (donc
$B'=B=A/\mathfrak p$) car $T_p$ étant alors exact, $T_p(B)$ est isomorphe à
$T_p(A)/T_p(\mathfrak p)$, donc est un $A$-module (et \emph{a fortiori} un
$B$-module) de type fini. Raisonnons par récurrence sur $p$; $f_p$ est l'image
canonique dans $B$ d'un élément $g_p\in A$, et si l'on pose $A'=A_{g_p}$, on a
$B'=A'/\mathfrak p'$ où $\mathfrak p'$ est un idéal premier minimal de $A'$,
égal à $\mathfrak p_{g_p}$. Comme on a
\[
  T_i(A_{g_p})=(T_i(A))_{g_p},
\]
les $T_i(A')$ sont des $A'$-modules de type fini, donc le foncteur
$T_\bullet^{(A')}$ vérifie les mêmes hypothèses que $T_\bullet$, mais en
remplaçant $i_0$ par $p$. On peut donc se borner au cas où $A'=A$, et où $T_p$
est exact; la suite exacte
\[
  0\longrightarrow\mathfrak p\longrightarrow A\longrightarrow A/\mathfrak p
  \longrightarrow0
\]
donne alors la suite exacte
\[
  T_{p+1}(A)\longrightarrow T_{p+1}(A/\mathfrak p)
  \xrightarrow{\partial}T_p(\mathfrak p)\longrightarrow T_p(A),
\]
et comme $T_p$ est exact, la dernière flèche de droite est injective, donc
$T_{p+1}(A/\mathfrak p)$ est un quotient de $T_{p+1}(A)$ et est par suite de
type fini. Notons maintenant que le raisonnement de (7.3.9) n'a utilisé le fait
que les $T_p(A)$ sont de type fini que pour $p\leq N$; on peut donc l'appliquer
à l'anneau intègre $B$, au foncteur $T_\bullet^{(B)}$ et à $N=p+1$, ce qui
achève le raisonnement par récurrence. Cela étant, il y a un
$f_N\in B-\{0\}$ tel que $\operatorname{Spec}(B_{f_N})$ soit un ouvert $V$
partout dense dans $\operatorname{Spec}(B)$ ne rencontrant aucune autre
composante irréductible de $\operatorname{Spec}(A)$. Si la proposition est
\oldpage[III]{185}
démontrée pour $B_{f_N}$, on aura un ouvert $W$ partout dense dans $V$ dans
lequel les fonctions de l'énoncé seront constantes, puisque
$A_x=(B_{f_N})_x$ pour tout $x\in W$. En faisant le même raisonnement pour
toute composante irréductible de $\operatorname{Spec}(A)$, le corollaire sera
démontré. On peut donc se borner au cas où $A$ est intègre; le raisonnement de
(7.3.9) prouve alors l'existence d'un $f\in A-\{0\}$ tel que les $T_p(A_f)$
soient des $A_f$-modules libres de type fini pour $p\leq N$, ce qui entraîne
la conclusion de (7.3.10) en vertu de (7.3.4).
\end{proof}

\begin{proposition}[7.3.11]
Soient $A$ un anneau commutatif local, $k$ son corps résiduel, $T_\bullet$ un
foncteur homologique covariant de $\mathbf{Ab}_A$ dans $\mathbf{Ab}$,
commutant aux sommes directes. On suppose qu'il existe $i_0$ tel que $T_i$ soit
exact pour $i\leq i_0$, et que tous les $T_n(A)$ soient des $A$-modules de
présentation finie. Alors les conditions équivalentes $a)$, $b)$, $c)$ de
(7.3.7) impliquent les deux suivantes, et leur sont équivalentes lorsque
l'anneau est en outre réduit :
\begin{enumerate}
  \item[d)] Pour tout $x\in\operatorname{Spec}(A)$, on a
    \[
      \operatorname{rang}_{\kappa(x)}T_q(\kappa(x))
      =
      \operatorname{rang}_kT_q(k)
      \quad\text{pour }q\leq p.
    \]
  \item[$d'$)] Pour tout point générique $x_j$ d'une composante irréductible
    de $\operatorname{Spec}(A)$, on a
    \[
      \operatorname{rang}_{\kappa(x_j)}T_q(\kappa(x_j))
      =
      \operatorname{rang}_kT_q(k)
      \quad\text{pour }q\leq p.
    \]
\end{enumerate}
\end{proposition}

\begin{proof}
Comme $T_q(A)$ est un $A$-module de présentation finie, la condition $b)$ de
(7.3.7) équivaut à dire que $T_q(A)$ est un $A$-module libre pour $q\leq p$
(Bourbaki, \emph{Alg. comm.}, chap. II, \S 3, n\textsuperscript{o} 2, cor. 2
de la prop. 5); la condition $c)$ implique que
\[
  T_q(\kappa(x))=T_q(A)\otimes_A\kappa(x)
  \quad\text{pour }q\leq p,
\]
donc les conditions équivalentes de (7.3.7) impliquent $d)$, et il est trivial
que $d)$ entraîne $d')$. Reste à prouver que $d')$ implique $a)$ lorsque $A$
est réduit. Raisonnons par récurrence sur $q\leq p$, puisque $T_q$ est exact
pour $q\leq i_0$. Supposons donc $T_k$ exact pour $k\leq q<p$ et montrons que
$T_{q+1}(A)$ est un $A$-module libre. En vertu de l'hypothèse de récurrence,
$T_{q+1}(A)\otimes_A M$ est isomorphe à $T_{q+1}(M)$ pour tout $A$-module $M$,
par la condition $c)$ de (7.3.7) et (7.3.3); appliquant cette propriété à
$M=\kappa(x_j)$ et $M=k$, on trouve, en vertu de l'hypothèse $d')$, que
\[
  \operatorname{rang}_{\kappa(x_j)}
  \bigl(T_{q+1}(A)\otimes_A\kappa(x_j)\bigr)
  =
  \operatorname{rang}_kT_{q+1}(k)
\]
pour tout $i$; mais cela implique que $T_{q+1}(A)$ est libre (Bourbaki,
\emph{Alg. comm.}, chap. II, \S 3, n\textsuperscript{o} 2, prop. 7), ce qui
achève la démonstration.
\end{proof}

Les résultats précédents seront considérablement améliorés pour les foncteurs
homologiques de type particulier que nous allons étudier dans (7.4); on
obtiendra en effet des critères d'exactitude ne faisant intervenir qu'un seul
des $T_p$.

\subsection{Critères d'exactitude pour les foncteurs
$H_\bullet(P_\bullet\otimes_A M)$}
\label{subsection:III.7.4.fr}

\begin{env}[7.4.1]
Soient $A$ un anneau (non nécessairement commutatif), $P_\bullet$ un complexe
de $A$-modules à droite plats. Comme le foncteur
$M\mapsto P_k\otimes_A M$ est alors exact dans $\mathbf{Ab}_A$ pour tout $k$,
le $\partial$-foncteur
\[
  T_\bullet(M)=H_\bullet(P_\bullet\otimes_A M)
  \tag{7.4.1.1}\label{III.7.4.1.1.fr}
\]
\oldpage[III]{186}
est un foncteur homologique de $\mathbf{Ab}_A$ dans $\mathbf{Ab}$,
évidemment $A$-linéaire lorsque $A$ est commutatif (7.1.2), et commutant aux
limites inductives.

Si $A$ est commutatif, alors, pour toute $A$-algèbre $B$, le foncteur
homologique $T_\bullet^{(B)}$ (7.3.8) est donné par définition par
\[
  T_\bullet^{(B)}(N)
  =
  H_\bullet(P_\bullet\otimes_A N_{[\rho]})
  \tag{7.4.1.2}\label{III.7.4.1.2.fr}
\]
où $\rho:A\to B$ est l'homomorphisme définissant la structure d'algèbre de
$B$; comme on peut aussi écrire
\[
  P_\bullet\otimes_A N_{[\rho]}
  =P_\bullet\otimes_A(B\otimes_B N)_{[\rho]}
  =(P_\bullet\otimes_A B)\otimes_B N,
\]
on voit que l'on a
\[
  T_\bullet^{(B)}(N)=H_\bullet(P'_\bullet\otimes_B N)
  \tag{7.4.1.3}\label{III.7.4.1.3.fr}
\]
pour tout $B$-module $N$, $P'_\bullet$ étant le complexe
$P_\bullet\otimes_A B$ de $B$-modules plats (0\textsubscript{I}, 6.2.1).
\end{env}

\begin{proposition}[7.4.2]
Sous les conditions générales de (7.4.1), et pour un entier $p\in\mathbf Z$
donné, les propriétés suivantes sont équivalentes :
\begin{enumerate}
  \item[a)] $T_p$ est exact à gauche (ou, ce qui revient au même, $T_{p+1}$
    est exact à droite).
  \item[b)]
    \[
      Z'_p(P_\bullet)
      =
      \operatorname{Coker}(P_{p+1}\longrightarrow P_p)
    \]
    est un $A$-module à droite plat.
  \item[c)] Il existe un complexe $P'_\bullet$ de $A$-modules à droite plats
    tel que la différentielle
    \[
      d'_{p+1}:P'_{p+1}\longrightarrow P'_p
    \]
    soit nulle, et un isomorphisme de foncteurs homologiques de
    $H_\bullet(P_\bullet\otimes_A M)$ sur
    $H_\bullet(P'_\bullet\otimes_A M)$.
\end{enumerate}
\end{proposition}

\begin{proof}
Par définition, on a une suite exacte fonctorielle en $M$
\[
  0\longrightarrow T_p(M)
  \longrightarrow Z'_p(P_\bullet\otimes M)
  \longrightarrow P_{p-1}\otimes M
\]
où
\[
  Z'_p(P_\bullet\otimes M)
  =
  \operatorname{Coker}(P_{p+1}\otimes M\longrightarrow P_p\otimes M)
  =
  Z'_p(P_\bullet)\otimes M
\]
en vertu de l'exactitude à droite du produit tensoriel. Pour tout homomorphisme
$f:M\to N$, on a donc un diagramme commutatif
\[
\xymatrix@C=30pt@R=38pt{
  0\ar[r]
  &T_p(M)\ar[r]\ar[d]_u
  &Z'_p(P_\bullet)\otimes M\ar[r]\ar[d]_v
  &P_{p-1}\otimes M\ar[d]_w
  \\
  0\ar[r]
  &T_p(N)\ar[r]
  &Z'_p(P_\bullet)\otimes N\ar[r]
  &P_{p-1}\otimes N
}
\tag{7.4.2.1}\label{III.7.4.2.1.fr}
\]
dont les lignes sont exactes. Si $f$ est un monomorphisme, il en est de même
de $w$ puisque $P_{p-1}$ est plat; si $T_p$ est exact à gauche, $u$ est lui
aussi un monomorphisme; on en conclut que $v$ est un monomorphisme, ce qui
entraîne que $Z'_p(P_\bullet)$ est plat. Inversement, s'il en est ainsi, $v$
est un monomorphisme pour tout monomorphisme $f:M\to N$, donc le diagramme
(7.4.2.1) montre que $u$ est un monomorphisme, et par suite $T_p$ (qui est déjà
semi-exact) est exact à gauche. Ainsi $a)$ et $b)$ sont équivalentes. Il est
immédiat que $c)$ entraîne $a)$, car si
$d'_{p+1}:P'_{p+1}\to P'_p$ est nulle, et
$0\to M'\to M\to M''\to0$ est une suite exacte de $A$-modules, l'opérateur
bord dans la suite exacte
\[
  H_{p+1}(P'_\bullet\otimes M'')
  \xrightarrow{\partial}H_p(P'_\bullet\otimes M')
  \longrightarrow H_p(P'_\bullet\otimes M)
\]
\oldpage[III]{187}
est nul par définition (M, IV, 1), donc $T_p$ est exact à gauche. Montrons
inversement que $b)$ implique $c)$. Si
\[
  Z_{p+1}(P_\bullet)=\operatorname{Ker}(P_{p+1}\longrightarrow P_p),
\]
on a une suite exacte
\[
  0\longrightarrow Z_{p+1}(P_\bullet)\longrightarrow P_{p+1}
  \longrightarrow Z'_p(P_\bullet)\longrightarrow0
\]
dans laquelle $P_{p+1}$ et $Z'_p(P_\bullet)$ sont plats, donc
$Z_{p+1}(P_\bullet)$ est plat ($0_{\mathrm I}$, 6.1.2). On va prendre
\[
  P'_i=P_i\quad\text{pour }i\ne p\text{ et }i\ne p+1,
  \qquad
  P'_p=Z'_p(P_\bullet),
  \qquad
  P'_{p+1}=Z_{p+1}(P_\bullet);
\]
pour différentielle $d'_i:P'_i\to P'_{i-1}$, on prendra celle du complexe
$P_\bullet$ pour $i\ne p$ et $i\ne p+1$, $0$ pour $i=p+1$ et pour $i=p$
l'homomorphisme $Z'_p(P_\bullet)\to P_{p-1}$ déduit de $d_p$ par passage au
quotient. Comme les $P_i$ sont plats, on a
\[
  Z'_i(P_\bullet\otimes M)=Z'_i(P_\bullet)\otimes M,
  \qquad
  Z_i(P_\bullet\otimes M)=Z_i(P_\bullet)\otimes M
\]
et
\[
  B_i(P_\bullet\otimes M)=B_i(P_\bullet)\otimes M
\]
(en posant $B_i(P_\bullet)=\operatorname{Im}(P_{i+1}\to P_i)$); on en conclut
aussitôt pour tout $M$ des isomorphismes fonctoriels
\[
  H_i(P_\bullet\otimes M)\xrightarrow{\sim}H_i(P'_\bullet\otimes M)
\]
pour tout $i$, et la vérification du fait qu'il s'agit d'un isomorphisme de
$\partial$-foncteurs découle sans peine de la définition de $\partial$
(M, IV, 1).
\end{proof}

On remarquera que les conditions de (7.4.2) impliquent aussi que
$B_p(P_\bullet)$ est plat, car on a une suite exacte
\[
  0\longrightarrow B_p(P_\bullet)\longrightarrow P_p
  \longrightarrow Z'_p(P_\bullet)\longrightarrow0,
\]
dans laquelle $P_p$ et $Z'_p(P_\bullet)$ sont plats ($0_{\mathrm I}$, 6.1.2).

\begin{corollary}[7.4.3]
Supposons que $A$ soit un anneau noethérien régulier de dimension $1$
(autrement dit un produit d'anneaux de Dedekind
($0_{\mathrm{IV}}$, 17.1.3 et 17.3.7), par exemple un anneau principal). Alors,
pour que $T_p$ soit exact à gauche, il faut et il suffit que $T_p(A)$ soit un
$A$-module plat. Pour que $T_p$ soit exact, il faut et il suffit que $T_p(A)$
et $T_{p-1}(A)$ soient des $A$-modules plats.
\end{corollary}

\begin{proof}
Rappelons que pour un module $M$ sur un anneau de Dedekind, il revient au même
de dire que $M$ est plat ou qu'il est sans torsion ($0_{\mathrm I}$, 6.3.3 et
6.3.4); sous les hypothèses de (7.4.3), tout sous-module d'un $A$-module plat
est donc plat.

La seconde assertion de (7.4.3) résulte de la première, puisque dire que $T_p$
est exact signifie que $T_p$ et $T_{p-1}$ sont exacts à gauche. Pour démontrer
la première assertion, notons que l'on a une suite exacte
\[
  0\longrightarrow H_p(P_\bullet)\longrightarrow Z'_p(P_\bullet)
  \longrightarrow B_{p-1}(P_\bullet)\longrightarrow0
\]
dans laquelle $B_{p-1}(P_\bullet)$ est un $A$-module plat, en tant que
sous-module du $A$-module plat $P_{p-1}$. Il est donc équivalent de dire que
$H_p(P_\bullet)$ est plat ou que $Z'_p(P_\bullet)$ est plat
($0_{\mathrm I}$, 6.1.2).
\end{proof}

Les applications les plus importantes de (7.4.2) sont les suivantes :

\begin{proposition}[7.4.4]
Soient $A$ un anneau noethérien, $P_\bullet$ un complexe de $A$-modules plats;
on suppose, soit que les $P_i$ sont de type fini, soit que les
$H_i(P_\bullet)$ sont des $A$-modules de type fini et qu'il existe $i_0$ tel
que $H_i(P_\bullet)=0$ pour $i<i_0$. Soit $T$ le foncteur homologique défini
par (7.4.1.1). Alors l'ensemble $U$ des
$y\in\operatorname{Spec}(A)$ tels que $(T_p)_y$ (7.1.4) soit exact à droite
(resp. exact à gauche, exact) est ouvert dans $\operatorname{Spec}(A)$.
\end{proposition}

\begin{proof}
Dans la seconde hypothèse sur $P_\bullet$, on peut remplacer $P_\bullet$ par
un complexe $P'_\bullet$ de
\oldpage[III]{188}
$A$-modules libres de type fini pour lequel le foncteur
$H_\bullet(P'_\bullet\otimes_A M)$ est isomorphe (en tant que
$\partial$-foncteur) à $T_\bullet(M)$ (0, 11.9.3). On peut donc toujours se
ramener à la première hypothèse et dans ce cas les $Z'_i(P_\bullet)$ sont de
type fini; en outre, on peut se borner à démontrer les assertions relatives à
l'exactitude à gauche (cf. (7.4.2, $a)$). Cela étant, soit $x\in U$; comme le
foncteur $M\mapsto M_x$ est exact, on a
\[
  (Z'_p(P_\bullet))_x=Z'_p((P_\bullet)_x),
\]
et (en tenant compte de (7.4.1.3)) l'hypothèse entraîne, en vertu de
(7.4.2, $b)$) que $(Z'_p(P_\bullet))_x$ est un $A_x$-module plat, donc libre
puisqu'il est de type fini et que $A_x$ est un anneau local noethérien
(Bourbaki, \emph{Alg. comm.}, chap. II, \S 3, n\textsuperscript{o} 2, cor. 2
de la prop. 5). On en conclut qu'il y a un $f\in A$ tel que
$(Z'_p(P_\bullet))_f$ soit libre sur $A_f$ (Bourbaki, \emph{Alg. comm.},
chap. II, \S 5, n\textsuperscript{o} 1, cor. de la prop. 2), et \emph{a
fortiori} $(Z'_p(P_\bullet))_y$ est libre sur $A_y$ pour tout $y\in D(f)$, ce
qui achève la démonstration en vertu de (7.4.2, $b)$).
\end{proof}

\begin{corollary}[7.4.5]
Sous les hypothèses de (7.4.4), supposons en outre que $A$ soit intègre. Alors
l'ensemble $U$ des $x\in\operatorname{Spec}(A)$ tels que $(T_p)_x$ soit exact
est ouvert non vide.
\end{corollary}

\begin{proof}
Il suffit en effet de prouver que $(T_p)_x$ est exact pour le point générique
$x$ de $\operatorname{Spec}(A)$, ce qui est immédiat puisque $A_x$ est un
corps, donc tout foncteur additif sur $\mathbf{Ab}_{A_x}$ est exact.
\end{proof}

\begin{proposition}[7.4.6]
Sous les hypothèses générales de (7.4.4), les conditions $a)$, $b)$ et $c)$ de
(7.4.2) sont aussi équivalentes à :
\begin{enumerate}
  \item[d)] Il existe un $A$-module $Q$ et un isomorphisme fonctoriel
    \[
      T_p(M)\xrightarrow{\sim}\operatorname{Hom}_A(Q,M).
      \tag{7.4.6.1}\label{III.7.4.6.1.fr}
    \]
    En outre, le $A$-module $Q$ est déterminé à isomorphisme unique près par
    cette propriété, et il est de type fini.
\end{enumerate}
\end{proposition}

\begin{proof}
L'unicité de $Q$ est un cas particulier de l'unicité d'un objet représentatif
d'un foncteur représentable (0, 8.1.5). Il est clair que le second membre de
(7.4.6.1) est exact à gauche. Inversement, pour démontrer l'existence de $Q$,
lorsque $T_p$ est exact à gauche, on peut d'abord, comme dans (7.4.4), se
ramener au cas où les $P_i$ sont plats et de type fini, donc (puisque $A$ est
noethérien) projectifs de type fini (Bourbaki, \emph{Alg. comm.}, chap. II,
\S 5, n\textsuperscript{o} 2, cor. du th. 1). Le dual $\check P_i$ de $P_i$
est alors aussi un $A$-module projectif de type fini, $P_i$ est canoniquement
isomorphe au dual de $\check P_i$, et l'homomorphisme canonique
\[
  P_i\otimes_A M\longrightarrow\operatorname{Hom}_A(\check P_i,M)
\]
est bijectif (Bourbaki, \emph{Alg.}, chap. II, 3\textsuperscript{e} éd.,
\S 4, n\textsuperscript{o} 2, prop. 2). On sait d'autre part (7.4.2, $c)$)
qu'on peut supposer que $d_{p+1}:P_{p+1}\to P_p$ est nulle, donc que l'on a une
suite exacte
\[
  0\longrightarrow T_p(M)\xrightarrow{u}P_p\otimes M
  \xrightarrow{v}P_{p-1}\otimes M
\]
où $v=d_p\otimes1$. Posons alors $Q'=\operatorname{Ker}(d_p)$, de sorte qu'on
a la suite exacte
\[
  0\longrightarrow Q'\xrightarrow{w}P_p\xrightarrow{d_p}P_{p-1},
\]
d'où, par transposition, la suite exacte
\[
  \check P_{p-1}\xrightarrow{{}^t d_p}\check P_p
  \xrightarrow{{}^t w}\check Q'\longrightarrow0.
\]
Nous allons voir que
\[
  Q=\check Q'=\operatorname{Coker}({}^t d_p)
\]
répond à la question. En effet, on a la suite exacte
\[
  0\longrightarrow\operatorname{Hom}_A(Q,M)
  \longrightarrow\operatorname{Hom}_A(\check P_p,M)
  \xrightarrow{v'}\operatorname{Hom}_A(\check P_{p-1},M)
\]
\oldpage[III]{189}
où $v'=\operatorname{Hom}({}^t d_p,1)$; lorsqu'on identifie canoniquement
$P_i\otimes M$ à $\operatorname{Hom}(\check P_i,M)$, $v'$ s'identifie donc à
$v=d_p\otimes1$, et l'on a par suite l'isomorphisme fonctoriel
\[
  T_p(M)\xrightarrow{\sim}\operatorname{Hom}_A(Q,M)
\]
cherché. En outre, $Q$, étant un quotient de $\check P_p$, est de type fini.
\end{proof}

\begin{proposition}[7.4.7]
Supposons vérifiées les conditions générales de (7.4.4). Alors, pour tout
$A$-module $M$ de type fini :
\begin{enumerate}
  \item[(i)] Les $T_i(M)$ sont des $A$-modules de type fini.
  \item[(ii)] Pour tout idéal $\mathfrak m$ de $A$, l'homomorphisme canonique
    \[
      \widehat{T_i(M)}
      \longrightarrow
      \varprojlim_n T_i\bigl(M\otimes_A(A/\mathfrak m^{n+1})\bigr)
      \tag{7.4.7.1}\label{III.7.4.7.1.fr}
    \]
    (où le premier membre est le séparé complété de $T_i(M)$ pour la topologie
    $\mathfrak m$-préadique) est bijectif.
\end{enumerate}
\end{proposition}

\begin{proof}
Comme dans (7.4.4), on se ramène d'abord au cas où les $P_i$ sont de type fini;
$A$ étant noethérien, les sous-modules des $P_i\otimes_A M$ sont de type fini,
d'où trivialement l'assertion (i). Quant à l'assertion (ii), elle résulte plus
généralement du lemme suivant :
\end{proof}

\begin{lemma}[7.4.7.2]
Soient $A$ un anneau noethérien, $u:E\to F$ un homomorphisme de $A$-modules de
type fini. Pour tout $A$-module de type fini, posons
\[
  K(M)=\operatorname{Ker}(u\otimes1_M),
  \qquad
  C(M)=\operatorname{Coker}(u\otimes1_M);
\]
alors les homomorphismes canoniques
\[
  \widehat{K(M)}\longrightarrow\varprojlim_n K(M_n),
  \qquad
  \widehat{C(M)}\longrightarrow\varprojlim_n C(M_n)
  \tag{7.4.7.3}\label{III.7.4.7.3.fr}
\]
(où l'on a posé
$M_n=M\otimes_A(A/\mathfrak m^{n+1})=M/\mathfrak m^{n+1}M$) sont bijectifs
pour tout idéal $\mathfrak m$ de $A$.
\end{lemma}

\begin{proof}
Comme $E\otimes M$ et $F\otimes M$ sont de type fini, et que le foncteur
$M\mapsto\widehat M$ est exact dans la catégorie des $A$-modules de type fini
($0_{\mathrm I}$, 7.3.3), $\widehat{K(M)}$ et $\widehat{C(M)}$ sont
respectivement le noyau et le conoyau de
\[
  \widehat{(u\otimes1)}:
  \widehat{(E\otimes M)}\longrightarrow\widehat{(F\otimes M)}.
\]
L'exactitude à gauche du foncteur $\varprojlim$ montre donc que
$\widehat{K(M)}=\varprojlim_n K(M_n)$; d'autre part, l'exactitude à droite du
produit tensoriel prouve que
$C(M_n)=C(M)\otimes_A(A/\mathfrak m^{n+1})$, donc
$\widehat{C(M)}=\varprojlim_n C(M_n)$ par définition.
\end{proof}

\begin{remark}[7.4.8]
Compte tenu de (6.10.5) et (6.10.6), on voit que, moyennant une hypothèse de
platitude supplémentaire, (7.4.7) redonne le fait que (4.1.7.1) est un
isomorphisme, c'est-à-dire l'essentiel du « premier théorème de comparaison »
pour les morphismes propres; en outre, l'énoncé s'applique non plus seulement
à un $\mathscr O_X$-Module cohérent, mais à un complexe de tels Modules. Il
serait intéressant d'obtenir un énoncé comprenant à la fois (7.4.7) et
(4.1.7.1) comme cas particuliers. On notera que lorsque les $P_i$ sont de type
fini, la démonstration de (7.4.7) n'utilise pas le fait que ce sont des modules
plats; il y aurait lieu d'examiner si la conclusion de (7.4.7) est encore
valable lorsque les $P_i$ ne sont pas supposés plats ni de type fini, mais que
les $H_i(P_\bullet)$ sont supposés de type fini pour tout $i$ et nuls pour
$i<i_0$. Est-il alors possible de remplacer $P_\bullet$ par un complexe
$P'_\bullet$ de $A$-modules de type fini tel que les foncteurs
$H_\bullet(P_\bullet\otimes M)$ et $H_\bullet(P'_\bullet\otimes M)$ (qui ne
sont plus homologiques) soient encore isomorphes ?
\end{remark}

\oldpage[III]{190}
\subsection{Cas des anneaux locaux noethériens}
\label{subsection:III.7.5.fr}

\begin{env}[7.5.1]
Soient $A$ un anneau local noethérien, $\mathfrak m$ son idéal maximal, et pour
tout $A$-module $M$, désignons par $\widehat M$ son séparé complété pour la
topologie $\mathfrak m$-préadique, isomorphe à
\[
  \varprojlim_n\bigl(M\otimes_A(A/\mathfrak m^{n+1})\bigr)
  =
  \varprojlim_n\bigl(M/\mathfrak m^{n+1}M\bigr).
\]
Soit $T$ un foncteur covariant additif de $\mathbf{Ab}_A$ dans $\mathbf{Ab}$;
les homomorphismes canoniques (7.2.3.1)
\[
  T(M)\otimes_A(A/\mathfrak m^{n+1})
  \longrightarrow
  T\bigl(M\otimes_A(A/\mathfrak m^{n+1})\bigr)
\]
forment évidemment un système projectif de $A$-homomorphismes, qui donnent donc
à la limite un $\widehat A$-homomorphisme fonctoriel en $M$
\[
  \widehat{T(M)}
  \longrightarrow
  \varprojlim_n T(M_n)
  \tag{7.5.1.1}\label{III.7.5.1.1.fr}
\]
où l'on a posé
\[
  M_n=M\otimes_A(A/\mathfrak m^{n+1}),
  \qquad
  A_n=A/\mathfrak m^{n+1}.
\]
\end{env}

\begin{proposition}[7.5.2]
Soient $A$ un anneau local noethérien d'idéal maximal $\mathfrak m$,
$k=A/\mathfrak m$ son corps résiduel, $T$ un foncteur covariant additif de
$\mathbf{Ab}_A$ dans $\mathbf{Ab}$, semi-exact et commutant aux limites
inductives. On suppose en outre que pour tout $A$-module de type fini $M$,
$T(M)$ est un $A$-module de type fini et que l'homomorphisme canonique
(7.5.1.1) est un isomorphisme. Sous ces conditions, les propriétés suivantes
sont équivalentes :
\begin{enumerate}
  \item[a)] $T$ est exact à droite.
  \item[b)] Pour tout $n$, le foncteur $N\mapsto T(N)$ est exact à droite dans
    la catégorie des $A_n$-modules de type fini (ce qui revient à dire que $T$
    est exact à droite dans la catégorie des $A$-modules de longueur finie).
  \item[c)] L'homomorphisme canonique
    \[
      T(\varepsilon):T(A)\longrightarrow T(k)
    \]
    est surjectif.
  \item[d)] Pour tout $n$ assez grand, l'homomorphisme canonique
    \[
      T(A_n)\longrightarrow T(k)
    \]
    est surjectif.
\end{enumerate}
\end{proposition}

\begin{proof}
Il est clair que a) entraîne b). Montrons que b) entraîne a), c'est-à-dire que
si $u:M\to N$ est un épimorphisme de $A$-modules, $T(u)$ est un épimorphisme.
Comme $T$ commute aux limites inductives et que le foncteur $\varinjlim$ est
exact dans la catégorie des modules (pour des ensembles d'indices filtrants),
on peut se borner au cas où $M$ et $N$ sont de type fini. Comme $T(M)$ et
$T(N)$ sont alors de type fini, et que $A$ est un anneau local noethérien, il
suffit de montrer que
\[
  \widehat{T(u)}:
  \widehat{T(M)}
  \longrightarrow
  \widehat{T(N)}
\]
est surjectif ($0_{\mathrm I}$, 7.3.5 et $0_{\mathrm I}$, 6.4.1). Par hypothèse,
$\widehat{T(M)}$ et $\widehat{T(N)}$ sont respectivement $\varprojlim_n T(M_n)$
et $\varprojlim_n T(N_n)$, donc $\widehat{T(u)}$ est limite du système projectif
d'homomorphismes
\[
  T(u\otimes1_{A_n}):T(M_n)\longrightarrow T(N_n).
\]
Or, b) signifie que ces homomorphismes sont surjectifs; en outre, $T(M_n)$ est
un $A_n$-module de type fini, et $A_n$ est un anneau artinien par hypothèse; on
en conclut que $\widehat{T(u)}$ est surjectif ($0_{\mathrm I}$, 13.1.2 et
13.2.2). Il est clair que a) entraîne c), et comme $T(\varepsilon)$ se factorise
en
\[
  T(A)\longrightarrow T(A_n)\longrightarrow T(k),
\]
c) implique d); enfin, il résulte de (7.2.7) que b) et d) sont équivalentes
puisque $T$ est semi-exact dans $\mathbf{Ab}_{A_n}$, ce qui achève la
démonstration.
\end{proof}

\begin{corollary}[7.5.3]
Sous les hypothèses générales de (7.5.2), si $T(k)=0$, alors $T(M)=0$ pour tout
$A$-module $M$.
\end{corollary}

\begin{proof}
\oldpage[III]{191}
Comme $k$ est le seul $A$-module simple, on déduit de (7.3.5.4) que $T(E)=0$
pour tout $A$-module de longueur finie $E$. Si maintenant $M$ est de type fini,
$\widehat{T(M)}$ est isomorphe à $\varprojlim_n T(M_n)$, et comme les $M_n$ sont
de longueur finie, on a $\widehat{T(M)}=0$; comme $T(M)$ est de type fini par
hypothèse, il est isomorphe à un sous-module de $\widehat{T(M)}$
($0_{\mathrm I}$, 7.3.5), donc on a $T(M)=0$. Enfin, pour un $A$-module $M$
quelconque, $T(M)$ est limite inductive des $T(N_\alpha)$ pour les sous-modules
de type fini $N_\alpha$ de $M$, ce qui achève la démonstration.
\end{proof}

\begin{proposition}[7.5.4]
Soient $A$ un anneau local noethérien d'idéal maximal $\mathfrak m$,
$k=A/\mathfrak m$ son corps résiduel, $T_\bullet$ un foncteur homologique de
$\mathbf{Ab}_A$ dans $\mathbf{Ab}$, commutant aux limites inductives. On suppose
en outre que pour tout $i$ et tout $A$-module $M$ de type fini, $T_i(M)$ est de
type fini et l'homomorphisme canonique
\[
  \widehat{T_i(M)}
  \longrightarrow
  \varprojlim_n T_i(M_n)
\]
est bijectif. Pour un entier $p$ donné, les conditions suivantes sont alors
équivalentes :
\begin{enumerate}
  \item[a)] $T_p$ est exact.
  \item[b)] $T_p$ est exact à droite, et $T_p(A)$ est un $A$-module libre.
  \item[c)] Les homomorphismes canoniques
    \[
      T_{p+1}(A)\longrightarrow T_{p+1}(k)
      \quad\text{et}\quad
      T_p(A)\longrightarrow T_p(k)
    \]
    sont surjectifs.
  \item[d)] Pour tout $n$, les homomorphismes canoniques
    \[
      T_{p+1}(A_n)\longrightarrow T_{p+1}(k)
      \quad\text{et}\quad
      T_p(A_n)\longrightarrow T_p(k)
    \]
    sont surjectifs.
  \item[e)] Pour tout $n$, le foncteur $N\mapsto T_p(N)$ est exact dans la
    catégorie des $A_n$-modules de type fini.
\end{enumerate}
\end{proposition}

\begin{proof}
On sait (7.3.3) que a) équivaut à dire que $T_{p+1}$ et $T_p$ sont exacts à
droite; comme $T_\bullet$ est un foncteur homologique dans la catégorie
$\mathbf{Ab}_{A_n}$, le même raisonnement que dans (7.3.1) montre que e)
équivaut à dire que $T_p$ et $T_{p+1}$ sont exacts à droite dans la catégorie
des $A_n$-modules de type fini. On déduit donc de (7.5.2) que a) et e) sont
équivalentes; l'équivalence de a), c) et d) résulte aussi de (7.5.2). Enfin, on
sait que tout $A$-module plat de type fini est libre (Bourbaki,
\emph{Alg. comm.}, chap. II, \S 3, n\textsuperscript{o} 2, cor. 2 de la prop.
5); l'équivalence de a) et b) résulte alors de (7.3.1) et (7.3.3).
\end{proof}

\begin{corollary}[7.5.5]
Supposons vérifiées les conditions générales de (7.5.4).
\begin{enumerate}
  \item[(i)] Si $T_p(k)=0$, on a $T_p=0$, $T_{p+1}$ est exact à droite et
    $T_{p-1}$ est exact à gauche.
  \item[(ii)] Si $T_{p-1}(k)=T_{p+1}(k)=0$, $T_p$ est exact,
    l'homomorphisme canonique
    \[
      T_p(A)\otimes_A M\longrightarrow T_p(M)
    \]
    est bijectif et $T_p(A)$ est un $A$-module libre.
\end{enumerate}
\end{corollary}

\begin{proof}
(i) découle aussitôt de (7.5.3) puisque $T_p$ est semi-exact, la dernière
assertion résultant de la définition d'un foncteur homologique. On conclut
aussitôt de (i) les deux premières assertions de (ii), compte tenu de (7.3.3);
le fait que $T_p(A)$ soit libre résulte de (7.5.4).
\end{proof}

\begin{corollary}[7.5.6]
Supposons vérifiées les hypothèses générales de (7.5.4), et supposons de plus
que $A$ soit un anneau de valuation discrète.
\begin{enumerate}
  \item[(i)] Pour que $T_p$ soit exact à droite, il faut et il suffit que
    $T_{p-1}(A)$ soit un $A$-module libre.
  \item[(ii)] Pour que $T_p$ soit exact, il faut et il suffit que $T_p(A)$ et
    $T_{p-1}(A)$ soient des $A$-modules libres.
\end{enumerate}
\end{corollary}

\begin{proof}
\oldpage[III]{192}
Il est clair que (i) implique (ii) (cf. (7.3.3)). Pour prouver (i), remarquons
que si $f$ est un générateur de l'idéal maximal de $A$ (« uniformisante » de
$A$), pour qu'un $A$-module de type fini $M$ soit libre (ou plat, ce qui revient
au même), il faut et il suffit que l'homothétie $h_f:x\mapsto fx$ de $M$ soit
injective, car cela équivaut ici à dire que $M$ est sans torsion
($0_{\mathrm I}$, 6.3.4). Considérons alors la suite exacte
\[
  0\longrightarrow A\xrightarrow{h_f}A\longrightarrow k\longrightarrow0,
\]
qui fournit la suite exacte d'homologie
\[
  T_p(A)
  \longrightarrow T_p(k)
  \longrightarrow T_{p-1}(A)
  \xrightarrow{h_f}T_{p-1}(A).
\]
On voit que $T_{p-1}(A)$ est libre si et seulement si
$T_p(A)\to T_p(k)$ est surjectif; la conclusion résulte alors de (7.5.2).
\end{proof}

\begin{remark}[7.5.7]
On notera que, en vertu de (7.4.7), les hypothèses générales de (7.4.4)
entraînent que le foncteur homologique $T_\bullet$ défini par (7.4.1.1)
satisfait aux hypothèses générales de (7.5.4). Dans ce cas, (7.5.6) est donc
contenu dans (7.4.3).
\end{remark}

\subsection{Descente des propriétés d'exactitude. Théorème de semi-continuité
et critère d'exactitude de Grauert}
\label{subsection:III.7.6.fr}

\begin{proposition}[7.6.1]
Sous les conditions de (7.4.1), soit $B$ une $A$-algèbre commutative. Si $T_p$
est exact à droite (resp. exact à gauche, exact) il en est de même de
$T_p^{(B)}$; la réciproque est vraie lorsque $B$ est un $A$-module fidèlement
plat.
\end{proposition}

\begin{proof}
La première assertion est un cas particulier d'une assertion triviale de
(7.1.3). Inversement, supposons d'abord que $B$ soit un $A$-module plat. On a
alors, pour tout $A$-module $M$,
\[
  H_\bullet\bigl(P_\bullet\otimes_A(M\otimes_A B)\bigr)
  =
  \bigl(H_\bullet(P_\bullet\otimes_A M)\bigr)\otimes_A B,
\]
ce qui s'écrit aussi, pour tout $p$,
\[
  T_p(M)\otimes_A B
  =
  T_p^{(B)}(M_{(B)})
  \tag{7.6.1.1}\label{III.7.6.1.1.fr}
\]
à un isomorphisme canonique près. Supposons $T_p^{(B)}$ exact à droite (resp.
exact à gauche, exact); comme $M\mapsto M_{(B)}$ est un foncteur exact, le
premier membre de (7.6.1.1) est un foncteur exact à droite (resp. exact à
gauche, exact) en $M$; si maintenant $B$ est fidèlement plat sur $A$, on en
déduit que $T_p$ a la même propriété d'exactitude ($0_{\mathrm I}$, 6.4.1).
\end{proof}

\begin{proposition}[7.6.2]
Sous les conditions de (7.4.1), on suppose en outre que $A$ soit un anneau
noethérien réduit et que les $P_i$ soient des $A$-modules de type fini. Pour que
$T_p$ soit exact à droite (resp. exact à gauche, exact), il faut et il suffit
que, pour toute $A$-algèbre $B$ qui est un anneau de valuation discrète,
$T_p^{(B)}$ le soit.
\end{proposition}

\begin{proof}
En vertu de (7.3.1) et (7.3.3), on peut se borner à considérer l'exactitude à
droite, et il n'y a naturellement qu'à prouver la suffisance de la condition
(7.6.1). En vertu de (7.4.2), il suffit de montrer que
$Z'_{p-1}(P_\bullet)$ est un $A$-module plat; comme $P_{p-1}$ est de type fini,
$Z'_{p-1}(P_\bullet)$ est aussi de type fini; le critère ($0$, 10.2.8) montre
qu'il suffit alors que $Z'_{p-1}(P_\bullet)\otimes_A B$ soit un $B$-module plat
pour toute $A$-algèbre $B$ qui est un anneau de valuation discrète. Or, comme
$P_\bullet$ est un complexe de $A$-modules plats, on a
\[
  Z'_{p-1}(P_\bullet)\otimes_A B
  =
  Z'_{p-1}(P_\bullet\otimes_A B);
\]
\oldpage[III]{193}
$P_\bullet\otimes_A B$ est un complexe de $B$-modules plats
($0_{\mathrm I}$, 6.2.1), et pour tout $B$-module $N$, on a
\[
  H_\bullet(P_\bullet\otimes_A N)
  =
  H_\bullet\bigl((P_\bullet\otimes_A B)\otimes_B N\bigr),
\]
donc
\[
  T_p^{(B)}(N)
  =
  H_p\bigl((P_\bullet\otimes_A B)\otimes_B N\bigr);
\]
appliquant (7.4.2) à $T_p^{(B)}$, on voit que l'hypothèse que $T_p^{(B)}$ est
exact à droite est équivalente au fait que $Z'_{p-1}(P_\bullet\otimes_A B)$ est
un $B$-module plat.

Le critère précédent amène à étudier de plus près le cas des anneaux de
valuation discrète :
\end{proof}

\begin{proposition}[7.6.3]
Sous les conditions de (7.4.1), supposons que $A$ soit un anneau noethérien
régulier de dimension 1 (autrement dit, que $A$ est noethérien et que, pour tout
$x\in\operatorname{Spec}(A)$, $A_x$ est un corps ou un anneau de valuation
discrète). Alors, pour tout entier $p$ et tout $A$-module $M$, on a une suite
exacte canonique fonctorielle en $M$
\[
  0
  \longrightarrow T_p(A)\otimes_A M
  \xrightarrow{t_M}T_p(M)
  \longrightarrow \operatorname{Tor}_1^A\bigl(T_{p-1}(A),M\bigr)
  \longrightarrow0.
  \tag{7.6.3.1}\label{III.7.6.3.1.fr}
\]
\end{proposition}

\begin{proof}
Dans ce qui suit, nous supprimerons pour simplifier la mention du complexe
$P_\bullet$ dans les notations homologiques usuelles $H_p(P_\bullet)$,
$B_p(P_\bullet)$, $Z_p(P_\bullet)$ et $Z'_p(P_\bullet)$. On a les trois suites
exactes
\begin{gather*}
  0\longrightarrow H_p\longrightarrow Z'_p\longrightarrow B_{p-1}
  \longrightarrow0,\\
  0\longrightarrow B_{p-1}\longrightarrow Z_{p-1}\longrightarrow H_{p-1}
  \longrightarrow0,\\
  0\longrightarrow Z_{p-1}\longrightarrow P_{p-1}\longrightarrow B_{p-2}
  \longrightarrow0.
\end{gather*}
Comme $P_{p-1}$ et $P_{p-2}$ sont plats, il en est de même de leurs
sous-modules respectifs $B_{p-1}$, $Z_{p-1}$ et $B_{p-2}$ puisqu'il y a
identité entre $A_x$-modules plats et $A_x$-modules sans torsion (pour tout
$x\in\operatorname{Spec}(A)$); par tensorisation avec $M$, on a donc les suites
exactes
\[
  0=\operatorname{Tor}_1^A(B_{p-1},M)
  \longrightarrow H_p\otimes M
  \longrightarrow Z'_p\otimes M
  \xrightarrow{u}B_{p-1}\otimes M
  \longrightarrow0
  \tag{7.6.3.2}\label{III.7.6.3.2.fr}
\]
\[
  0=\operatorname{Tor}_1^A(Z_{p-1},M)
  \longrightarrow\operatorname{Tor}_1^A(H_{p-1},M)
  \longrightarrow B_{p-1}\otimes M
  \xrightarrow{v}Z_{p-1}\otimes M
  \tag{7.6.3.3}\label{III.7.6.3.3.fr}
\]
\[
  0=\operatorname{Tor}_1^A(B_{p-2},M)
  \longrightarrow Z_{p-1}\otimes M
  \xrightarrow{w}P_{p-1}\otimes M.
  \tag{7.6.3.4}\label{III.7.6.3.4.fr}
\]
Par définition,
$T_p(M)=\operatorname{Ker}(d_p\otimes1)/\operatorname{Im}(d_{p+1}\otimes1)$;
c'est donc le noyau de l'homomorphisme
\[
  (P_p\otimes M)/\operatorname{Im}(d_{p+1}\otimes1)
  \longrightarrow P_{p-1}\otimes M
\]
obtenu à partir de $d_p\otimes1$ par passage au quotient, homomorphisme qui
s'écrit aussi $Z'_p\otimes M\to P_{p-1}\otimes M$ par définition de
$Z'_p=P_p/B_p$; or, cet homomorphisme peut être considéré comme le composé
\[
  Z'_p\otimes M
  \xrightarrow{u}B_{p-1}\otimes M
  \xrightarrow{v}Z_{p-1}\otimes M
  \xrightarrow{w}P_{p-1}\otimes M.
\]
Comme $w$ est injectif d'après (7.6.3.4), on a une suite exacte
\[
  0\longrightarrow\operatorname{Ker}u
  \longrightarrow T_p(M)
  \longrightarrow\operatorname{Ker}v
  \longrightarrow0,
\]
qui n'est autre que (7.6.3.1), compte tenu de (7.6.3.2) et (7.6.3.3) et de ce
que $H_p=T_p(A)$ par définition.
\end{proof}

\begin{remarks}[7.6.4]
\begin{enumerate}
  \item[(i)] $H_\bullet(P_\bullet\otimes_A M)$ est l'homologie du bicomplexe
    $P_\bullet\otimes_A M$, où $M$ est considéré comme un complexe réduit à son
    terme de degré 0; elle est par suite (6.3.6 et 6.3.2) l'aboutissement de la
    suite spectrale régulière dont le terme $E_2$ est
    \[
      E^2_{pq}
      =
      \operatorname{Tor}_p^A\bigl(H_q(P_\bullet),M\bigr)
      =
      \operatorname{Tor}_p^A\bigl(T_q(A),M\bigr).
    \]
    \oldpage[III]{194}
    Or, l'hypothèse sur l'anneau $A$ entraîne que
    $\operatorname{Tor}_p^A(E,F)=0$ pour $p\geq2$ et pour des $A$-modules
    quelconques ($0_{\mathrm{IV}}$, 17.2.2); on sait (M, XV) que cela entraîne
    l'exactitude de la suite
    \[
      0\longrightarrow E^2_{0,q}
      \longrightarrow H_q(P_\bullet\otimes_A M)
      \longrightarrow E^2_{1,q-1}
      \longrightarrow0
    \]
    qui n'est autre que (7.6.3.1).
  \item[(ii)] Compte tenu de (7.3.1), la suite exacte (7.6.3.1) redonne comme
    cas particulier le résultat de (7.4.3).
\end{enumerate}
\end{remarks}

\begin{corollary}[7.6.5]
Sous les conditions de (7.4.1), supposons que $A$ soit un anneau de valuation
discrète, de corps des fractions $K$, de corps résiduel $k$, et que les $T_i(A)$
soient des $A$-modules de type fini. On a alors
\[
  \operatorname{rang}_k T_p(k)
  \geq
  \operatorname{rang}_k\bigl(T_p(A)\otimes_A k\bigr)
  \geq
  \operatorname{rang}_A T_p(A)
  =
  \operatorname{rang}_K T_p(K).
  \tag{7.6.5.1}\label{III.7.6.5.1.fr}
\]
En outre, pour que les termes extrêmes de cette inégalité soient égaux, il faut
et il suffit que $T_p$ soit exact, ou encore que $T_p(A)$ et $T_{p-1}(A)$ soient
des $A$-modules libres.
\end{corollary}

\begin{proof}
Faisons en effet $M=k$ dans la suite exacte (7.6.3.1); il vient, puisqu'il
s'agit d'espaces vectoriels sur $k$
\[
  \operatorname{rang}_k T_p(k)
  =
  \operatorname{rang}_k\bigl(T_p(A)\otimes_A k\bigr)
  +
  \operatorname{rang}_k\bigl(\operatorname{Tor}_1^A(T_{p-1}(A),k)\bigr).
\]
D'autre part, comme $T_p(A)$ est un module de type fini sur l'anneau local
intègre $A$, on a (Bourbaki, \emph{Alg. comm.}, chap. II, \S 3,
n\textsuperscript{o} 2, cor. 1 de la prop. 4)
\[
  \operatorname{rang}_k\bigl(T_p(A)\otimes_A k\bigr)
  \geq
  \operatorname{rang}_A T_p(A)
  =
  \operatorname{rang}_K\bigl(T_p(A)\otimes_A K\bigr)
  \tag{7.6.5.2}\label{III.7.6.5.2.fr}
\]
et en outre les deux membres de (7.6.5.2) sont égaux si et seulement si $T_p(A)$
est un $A$-module libre (\emph{loc. cit.}, prop. 7). On notera d'ailleurs que
puisque $K$ est un $A$-module plat, on a par définition
\[
  T_p(A)\otimes_A K
  =
  H_p(P_\bullet)\otimes_A K
  =
  H_p(P_\bullet\otimes_A K)
  =
  T_p(K).
\]
On a donc bien l'inégalité (7.6.5.1) et on voit en outre que l'égalité n'est
possible que si : 1\textsuperscript{o} $T_p(A)$ est libre;
2\textsuperscript{o} $\operatorname{Tor}_1^A(T_{p-1}(A),k)=0$, condition qui
équivaut, comme on sait ($0$, 10.1.3), au fait que $T_{p-1}(A)$ est un
$A$-module libre. Enfin, comme les $T_i(A)$ sont des $A$-modules de type fini,
il revient au même de dire qu'ils sont plats ou libres (Bourbaki,
\emph{Alg. comm.}, chap. II, \S 3, n\textsuperscript{o} 2, cor. 2 de la prop.
5), et on conclut par (7.4.3).
\end{proof}

\begin{env}[7.6.6]
Les hypothèses étant toujours celles de (7.4.1), nous poserons, pour tout
$x\in\operatorname{Spec}(A)$,
\[
  d_p(x)=d_p^T(x)=\operatorname{rang}_{\kappa(x)}T_p(\kappa(x)).
  \tag{7.6.6.1}\label{III.7.6.6.1.fr}
\]
\end{env}

\begin{lemma}[7.6.7]
Soit $\varphi:A\to A'$ un homomorphisme d'anneaux, et soit
\[
  f={}^a\!\varphi:
  \operatorname{Spec}(A')\longrightarrow\operatorname{Spec}(A)
\]
l'application correspondante (I, 1.2.1). Si l'on pose
$T'_\bullet=T_\bullet^{(A')}$ (7.1.3), on a
\[
  d_p^{T'}=d_p^T\circ f.
  \tag{7.6.7.1}\label{III.7.6.7.1.fr}
\]
\end{lemma}

\begin{proof}
\oldpage[III]{195}
En effet, pour tout $x'\in\operatorname{Spec}(A')$, on a, en posant
$x=f(x')$,
\[
  H_\bullet(P_\bullet\otimes_A\kappa(x'))
  =
  H_\bullet\bigl((P_\bullet\otimes_A\kappa(x))
    \otimes_{\kappa(x)}\kappa(x')\bigr)
  =
  H_\bullet(P_\bullet\otimes_A\kappa(x))
    \otimes_{\kappa(x)}\kappa(x'),
\]
puisque $\kappa(x')$ est plat sur $\kappa(x)$, d'où la relation (7.6.7.1).
\end{proof}

\begin{lemma}[7.6.8]
Si l'anneau $A$ est noethérien et le complexe $P_\bullet$ formé de $A$-modules
de type fini, la fonction $x\mapsto d_p^T(x)$ sur $\operatorname{Spec}(A)$ est
constructible.
\end{lemma}

\begin{proof}
Il faut prouver que pour toute partie fermée irréductible $Y$ de
$X=\operatorname{Spec}(A)$, il existe un ouvert non vide $U$ de $Y$ dans lequel
$d_p$ est constante ($0$, 9.2.2); comme
$Y=\operatorname{Spec}(A/\mathfrak a)$, où $\mathfrak a$ est un idéal de $A$ tel
que $A/\mathfrak a$ soit réduit, on peut, en vertu de (7.6.7), se borner au cas
où $Y=X$ et où $A$ est un anneau noethérien intègre; mais alors l'assertion
résulte de (7.4.5).
\end{proof}

\begin{theorem}[7.6.9]
Soient $A$ un anneau noethérien, $P_\bullet$ un complexe de $A$-modules plats de
type fini, $T_\bullet(M)=H_\bullet(P_\bullet\otimes_A M)$ le foncteur homologique
défini par $P_\bullet$; pour tout $x\in\operatorname{Spec}(A)$, soit
\[
  d_p(x)=\operatorname{rang}_{\kappa(x)}T_p(\kappa(x)).
\]
Alors :
\begin{enumerate}
  \item[(i)] La fonction $d_p$ est constructible et semi-continue
    supérieurement dans $\operatorname{Spec}(A)$.
  \item[(ii)] Si $T_p$ est exact, $d_p$ est continue (donc localement constante)
    dans $\operatorname{Spec}(A)$; la réciproque est vraie lorsque l'anneau $A$
    est réduit.
\end{enumerate}
\end{theorem}

\begin{proof}
\begin{enumerate}
  \item[(i)] La première assertion a été démontrée dans (7.6.8). Pour prouver
    la seconde, il suffit donc ($0$, 9.2.4) de montrer que si $x'\ne x$ est une
    générisation de $x$ dans $\operatorname{Spec}(A)$, on a
    $d_p(x')\leq d_p(x)$. Or, il existe alors un anneau de valuation discrète
    $B$ et un morphisme
    \[
      f:\operatorname{Spec}(B)\longrightarrow\operatorname{Spec}(A)
    \]
    tels que, si $a$ désigne le point fermé de $\operatorname{Spec}(B)$ et $b$
    son point générique, on ait $f(a)=x$ et $f(b)=y$ (II, 7.1.9). En vertu de
    la formule (7.6.7.1), on voit qu'on est ramené à démontrer l'inégalité
    $d_p(a)\geq d_p(b)$ dans $\operatorname{Spec}(B)$; mais cela n'est autre que
    l'inégalité (7.6.5.1)\footnote{Le principe de démonstration de (i) par
    réduction au cas d'un anneau de valuation discrète nous a été communiqué
    oralement par Hironaka.}.
  \item[(ii)] La première assertion a déjà été démontrée (7.3.4). Pour démontrer
    la réciproque, utilisons le critère valuatif (7.6.2); compte tenu de la
    formule (7.6.7.1), on est donc ramené au cas où $A$ est un anneau de
    valuation discrète; mais comme $\operatorname{Spec}(A)$ ne comporte alors
    que deux points, l'hypothèse que $d_p$ est constante implique bien que $T_p$
    est exact, en vertu de (7.6.5).
\end{enumerate}
\end{proof}

\begin{corollary}[7.6.10]
Soient $A$ un anneau noethérien, $\mathfrak p_i$ $(1\leq i\leq r)$ ses idéaux
premiers minimaux, $k_i$ le corps résiduel de $A_{\mathfrak p_i}$
$(1\leq i\leq r)$.
\begin{enumerate}
  \item[(i)] Pour tout $x\in\operatorname{Spec}(A)$, il existe un indice $i$
    tel que
    \[
      d_p(x)\geq\operatorname{rang}_{k_i}T_p(k_i).
      \tag{7.6.10.1}\label{III.7.6.10.1.fr}
    \]
\end{enumerate}
En particulier, si $A$ est intègre et si $K$ est son corps des fractions, on a
\[
  d_p(x)\geq\operatorname{rang}_K T_p(K)
  \tag{7.6.10.2}\label{III.7.6.10.2.fr}
\]
pour tout $x\in\operatorname{Spec}(A)$.
\oldpage[III]{196}
\begin{enumerate}
  \item[(ii)] Supposons en outre que $A$ soit local et réduit, et soit $k$ son
    corps résiduel. Alors, pour que $T_p$ soit exact, il faut et il suffit que
    l'on ait
    \[
      \operatorname{rang}_k T_p(k)
      =
      \operatorname{rang}_{k_i}T_p(k_i)
      \qquad\text{pour }1\leq i\leq r.
      \tag{7.6.10.3}\label{III.7.6.10.3.fr}
    \]
\end{enumerate}
\end{corollary}

\begin{proof}
(i) est immédiat puisque tout voisinage de $x$ contient un des $\mathfrak p_i$,
et il suffit d'appliquer la définition de la semi-continuité. D'autre part, si
$A$ est local, le seul voisinage dans $\operatorname{Spec}(A)$ de l'idéal
maximal $\mathfrak m$ est $\operatorname{Spec}(A)$ tout entier, donc on a
\[
  d_p(x)\leq\operatorname{rang}_k T_p(k)
\]
pour tout $x\in\operatorname{Spec}(A)$; cette relation, jointe à (i), montre que
la condition (7.6.10.3) entraîne que $d_p(x)$ est constante dans
$\operatorname{Spec}(A)$, et par suite que $T_p$ est exact en vertu de
(7.6.9, (ii)); la réciproque est évidente en vertu de (7.6.9, (ii)).
\end{proof}

\begin{remark}[7.6.11]
On peut se demander si l'assertion de (7.6.9, (i)) ne peut être renforcée par
l'inégalité
\[
  \operatorname{rang}_{\kappa(x)}T_p(\kappa(x))
  \geq
  \operatorname{rang}_{\kappa(x)}
    \bigl(T_p(A)\otimes_A\kappa(x)\bigr)
  \tag{7.6.11.1}\label{III.7.6.11.1.fr}
\]
pour tout $x\in\operatorname{Spec}(A)$, qui a lieu effectivement lorsque $A$
est un anneau de valuation discrète et $x$ son idéal maximal (7.6.5).
Bornons-nous au cas où $A$ est un anneau local noethérien, d'idéal maximal
$\mathfrak m$ et de corps résiduel $k$. Alors, les conditions suivantes sont
équivalentes :
\begin{enumerate}
  \item[a)] Pour tout complexe $P_\bullet$ de $A$-modules plats de type fini,
    on a
    \[
      \operatorname{rang}_k T_p(k)
      \geq
      \operatorname{rang}_k\bigl(T_p(A)\otimes_A k\bigr)
      \quad\text{pour tout entier }p.
      \tag{7.6.11.2}\label{III.7.6.11.2.fr}
    \]
  \item[b)] Pour tout $A$-module $M$ de type fini, on a
    \[
      \operatorname{rang}_k(M\otimes_A k)
      \geq
      \operatorname{rang}_k(\check M\otimes_A k).
      \tag{7.6.11.3}\label{III.7.6.11.3.fr}
    \]
  \item[c)] Pour tout $A$-module $N$ de type fini, on a
    \[
      \operatorname{rang}_k\bigl(\operatorname{Tor}_1^A(N,k)\bigr)
      \geq
      \operatorname{rang}_k\bigl(\operatorname{Tor}_2^A(N,k)\bigr).
      \tag{7.6.11.4}\label{III.7.6.11.4.fr}
    \]
\end{enumerate}
On notera qu'il revient au même, par décalage (M, V, 7.2), de dire que l'on a,
pour tout $i\geq1$,
\[
  \operatorname{rang}_k\bigl(\operatorname{Tor}_i^A(N,k)\bigr)
  \geq
  \operatorname{rang}_k\bigl(\operatorname{Tor}_{i+1}^A(N,k)\bigr).
  \tag{7.6.11.5}\label{III.7.6.11.5.fr}
\]
Donnons rapidement quelques indications sur la démonstration. Pour voir que a)
entraîne b), on considère une suite exacte
\[
  L_1\xrightarrow{d}L_0\longrightarrow M\longrightarrow0
\]
où $L_0$ et $L_1$ sont libres de type fini, et on applique a) au complexe
\[
  P_1\xrightarrow{{}^t d}P_0,
  \qquad
  P_0=\check L_1,
  \quad
  P_1=\check L_0,
\]
les autres termes étant nuls; on a alors $T_1(A)=\check M$ et
\[
  T_1(k)
  =
  \operatorname{Hom}_A(M,k)
  =
  \operatorname{Hom}_k(M/\mathfrak mM,k),
\]
autrement dit $T_1(k)$ est le dual de l'espace vectoriel $M\otimes_A k$, et a
donc même rang que ce dernier. Pour prouver que b) entraîne c), nous établissons
d'abord le lemme suivant :

\begin{lemma}[7.6.11.6]
Étant donné un complexe
\[
  \cdots\longrightarrow0\longrightarrow
  P_1\xrightarrow{d}P_0
  \longrightarrow0\longrightarrow\cdots
\]
de $A$-modules plats, on a une suite exacte
\[
  0\longrightarrow\operatorname{Tor}_2^A(Z'_0,k)
  \longrightarrow T_1(A)\otimes_A k
  \longrightarrow T_1(k)
  \longrightarrow\operatorname{Tor}_1^A(Z'_0,k)
  \longrightarrow0.
  \tag{7.6.11.7}\label{III.7.6.11.7.fr}
\]
\end{lemma}

\begin{proof}
\oldpage[III]{197}
En effet, partant de la suite exacte
$0\to Z_1\to P_1\to B_0\to0$, on en tire la suite exacte
\[
  0\longrightarrow\operatorname{Tor}_1^A(B_0,k)
  \longrightarrow Z_1\otimes k
  \longrightarrow P_1\otimes k
  \xrightarrow{u}B_0\otimes k
  \longrightarrow0.
\]
De la suite exacte $0\to B_0\to Z_0\to Z'_0\to0$, on tire, puisque $Z_0=P_0$
est plat,
$\operatorname{Tor}_1^A(B_0,k)=\operatorname{Tor}_2^A(Z'_0,k)$; par définition,
on a $Z_1=T_1(A)$; enfin, on a $T_1(k)=\operatorname{Ker}(d\otimes1)$, et
$d\otimes1$ se factorise en
\[
  P_1\otimes k
  \xrightarrow{u}B_0\otimes k
  \xrightarrow{v}Z_0\otimes k
  =P_0\otimes k;
\]
on a $T_1(k)=u^{-1}(R)$, où $R=\operatorname{Ker}v$, et comme $u$ est
surjectif, $R=u(T_1(k))$; enfin,
$R=\operatorname{Tor}_1^A(Z'_0,k)$ par définition de $v$, ce qui achève
d'établir la suite exacte (7.6.11.7).
\end{proof}

Pour déduire alors c) de b), on considère une suite exacte
\[
  L_1\xrightarrow{d}L_0\longrightarrow N\longrightarrow0,
\]
où $L_0$ et $L_1$ sont des modules libres de type fini; considérons le foncteur
$T$ associé au complexe formé de $L_1$ et $L_0$; comme $L_0$ et $L_1$ sont
libres, ils s'identifient à leurs biduals; donc si
\[
  M=\operatorname{Coker}({}^t d),
  \qquad
  T_1(A)=\operatorname{Ker}(d)=\check M;
\]
par ailleurs, $M\otimes_A k=\operatorname{Coker}({}^t d\otimes1_k)$ a même rang
sur $k$ que $\operatorname{Ker}(d\otimes1_k)$. L'hypothèse b) entraîne par suite
que
\[
  \operatorname{rang}_k\bigl(T_1(A)\otimes_A k\bigr)
  \leq
  \operatorname{rang}_k\bigl(T_1(k)\bigr);
\]
comme $Z'_0=N$, l'inégalité (7.6.11.4) résulte donc de la suite exacte
(7.6.11.7).

Enfin, pour prouver que c) entraîne a), appliquons (7.6.11.6) en remplaçant
$P_0$ et $P_1$ par $P_p$ et $P_{p-1}$; l'hypothèse c) appliquée au module
$Z'_p$ donne $\operatorname{rang}_kR\geq\operatorname{rang}_kS$, où
\[
  R=\operatorname{Ker}(P_p\otimes k\longrightarrow P_{p-1}\otimes k)
  \quad\text{et}\quad
  S=Z_p\otimes k.
\]
Or, si l'on factorise $d_{p+1}:P_{p+1}\to P_p$ en
\[
  P_{p+1}\xrightarrow{v}Z_p\xrightarrow{j}P_p,
\]
on a
\[
  \operatorname{Im}(d_{p+1}\otimes1)
  =
  (j\otimes1)\bigl(\operatorname{Im}(v\otimes1)\bigr).
\]
Comme
\[
  T_p(A)\otimes_A k
  =(Z_p/B_p)\otimes_A k
  =(Z_p\otimes k)/\operatorname{Im}(v\otimes1),
\]
et
\[
  T_p(k)=R/\operatorname{Im}(d_{p+1}\otimes1),
\]
on en conclut bien l'inégalité (7.6.11.2).

Cela étant, supposons que l'anneau local $A$ soit régulier de dimension $n$;
on sait alors [17] que le $A$-module
$\operatorname{Tor}_i^A(k,k)$ est isomorphe à la puissance extérieure
\[
  \bigwedge\nolimits^i(\mathfrak m/\mathfrak m^2);
\]
on voit donc que la condition (7.6.11.4) n'est pas vérifiée pour $N=k$, dès que
$n\geq4$.

Par contre, si l'anneau local intègre $A$ est tel que tout $A$-module réflexif
de type fini soit libre (ce qui est le cas lorsque $A$ est un anneau régulier
de dimension 2), la condition (7.6.11.3) est vérifiée : en effet, on sait que le
dual $\check M$ d'un $A$-module $M$ de type fini est réflexif, donc libre, et
par suite
\[
  \operatorname{rang}_k(\check M\otimes_A k)
  =
  \operatorname{rang}_K(\check M)
  =
  \operatorname{rang}_K(M),
\]
($K$ corps des fractions de $A$); par ailleurs, on sait que toute base sur $k$
de $M\otimes_A k$ est formée d'images d'un système de générateurs de $M$
(Bourbaki, \emph{Alg. comm.}, chap. II, \S 3, n\textsuperscript{o} 2, cor. 2 de
la prop. 4), donc
$\operatorname{rang}_K(M)\leq\operatorname{rang}_k(M\otimes_A k)$, ce qui
prouve notre assertion.
\end{remark}

\subsection{Application aux morphismes propres : I. La propriété d'échange}
\label{subsection:III.7.7.fr}

Les trois numéros qui suivent sont, pour l'essentiel, des traductions, dans le
langage des morphismes de préschémas, des résultats des numéros précédents.

\begin{env}[7.7.1]
Soit $f:X\to Y$ un morphisme quasi-compact et séparé de préschémas, et soit
$\mathscr P_\bullet$ un complexe de $\mathscr O_X$-Modules quasi-cohérents, dont
l'opérateur de dérivation est de degré $-1$; supposons en outre que les
$\mathscr O_X$-Modules $\mathscr P_i$ sont $Y$-plats ($0_{\mathrm I}$, 6.7.1).
\oldpage[III]{198}
Nous allons considérer le $\partial$-foncteur
$\mathscr M\mapsto\mathscr T_\bullet(\mathscr M)$ (aussi noté
$\mathscr T_\bullet(\mathscr P_\bullet,\mathscr M)$) dans la catégorie des
$\mathscr O_Y$-Modules quasi-cohérents, à valeurs dans la catégorie des
$\mathscr O_Y$-Modules quasi-cohérents (en vertu de (6.2.3)), défini par
\[
  \mathscr T_n(\mathscr P_\bullet,\mathscr M)
  =\mathscr T_n(\mathscr M)
  =\mathscr H^{-n}
    \bigl(f,\mathscr P^\bullet\otimes_{\mathscr O_Y}\mathscr M\bigr)
  \qquad\text{pour }n\in\mathbf Z,
  \tag{7.7.1.1}\label{III.7.7.1.1.fr}
\]
où $\mathscr P^\bullet$ est le complexe dont le terme de degré $j$ est
$\mathscr P_{-j}$, l'opérateur de dérivation étant donc alors de degré $+1$. Le foncteur
$\mathscr T_\bullet$ ainsi défini est un \emph{foncteur homologique} en
$\mathscr M$ (6.2.6).
\end{env}

\begin{env}[7.7.2]
Soit $g:Y'\to Y$ un morphisme, et posons
\[
  X'=X_{(Y')}=X\times_Y Y',
  \qquad
  f'=f_{(Y')}:X'\longrightarrow Y',
\]
qui est un morphisme quasi-compact et séparé; soit d'autre part
\[
  \mathscr P'_\bullet
  =\mathscr P_\bullet\otimes_{\mathscr O_Y}\mathscr O_{Y'};
\]
c'est un complexe de $\mathscr O_{X'}$-Modules quasi-cohérents qui sont
$Y'$-plats en vertu de (I, 9.1.12) et ($0_{\mathrm I}$, 6.2.1). Nous poserons
(avec les mêmes conventions sur les degrés)
\[
  \mathscr T^{Y'}_\bullet(\mathscr M')
  =\mathscr H^\bullet
    \bigl(f',\mathscr P'^\bullet\otimes_{\mathscr O_{Y'}}\mathscr M'\bigr)
  =\mathscr H^\bullet
    \bigl(f',\mathscr P^\bullet\otimes_{\mathscr O_Y}\mathscr M'\bigr)
  \tag{7.7.2.1}\label{III.7.7.2.1.fr}
\]
qui est un foncteur homologique en le $\mathscr O_{Y'}$-Module quasi-cohérent
$\mathscr M'$. Lorsque $Y'$ est un schéma affine d'anneau $A'$, on écrira
$\mathscr T^{A'}_\bullet$ au lieu de $\mathscr T^{Y'}_\bullet$; pour tout
$A'$-module $M'$, on a alors
\[
  \mathscr T^{A'}_\bullet(\widetilde{M'})
  =\bigl(\Gamma(Y',\mathscr T^{Y'}_\bullet(\widetilde{M'}))\bigr)^\sim;
\]
on posera
\[
  T^{A'}_\bullet(M')
  =\Gamma(Y',\mathscr T^{Y'}_\bullet(\widetilde{M'})),
\]
qui est un foncteur homologique de $A'$-modules, à valeurs dans la catégorie des
$A'$-modules.

On observera que si $Y=\operatorname{Spec}(A)$ est aussi affine, le foncteur de
$A'$-modules $T^{A'}_\bullet$ coïncide avec le foncteur obtenu par extension des
scalaires de $A$ à $A'$ à partir du foncteur homologique de $A$-modules
$T^A_\bullet$ (7.1.3) : en effet, soit $g:Y'\to Y$ le morphisme correspondant à
l'homomorphisme d'anneaux $A\to A'$, et soit $g':X'\to X$ le morphisme
correspondant qui est affine (II, 1.6.2); si $\mathfrak U$ est un recouvrement
ouvert affine de $X$, $\mathfrak U'=g'^{-1}(\mathfrak U)$ est un recouvrement
ouvert affine de $X'$; en vertu de (6.2.2), tout revient à voir que
\[
  C^\bullet
    \bigl(\mathfrak U,
      \mathscr P_\bullet\otimes_{\mathscr O_Y}g_*(\mathscr M')\bigr)
  =C^\bullet
    \bigl(\mathfrak U',
      \mathscr P_\bullet\otimes_{\mathscr O_Y}\mathscr M'\bigr),
\]
et finalement, que pour tout ouvert affine $U$ de $X$, en posant
$U'=g'^{-1}(U)$, on a
\[
  \Gamma
    \bigl(U,\mathscr P_\bullet\otimes_{\mathscr O_Y}g_*(\mathscr M')\bigr)
  =\Gamma
    \bigl(U',\mathscr P_\bullet\otimes_{\mathscr O_Y}\mathscr M'\bigr),
\]
ce qui est trivial (I, 1.3 et 3.2).

En particulier, si $U$ est un ouvert de $Y$, on a, pour tout
$\mathscr O_Y$-Module quasi-cohérent $\mathscr M$
\[
  \mathscr T^U_\bullet(\mathscr M|U)
  =\bigl(\mathscr T_\bullet(\mathscr M)\bigr)|U.
  \tag{7.7.2.2}\label{III.7.7.2.2.fr}
\]
\end{env}

\begin{env}[7.7.3]
Pour tout $\mathscr O_Y$-Module quasi-cohérent $\mathscr M$, on a un
homomorphisme canonique, fonctoriel en $\mathscr M$ :
\[
  \mathscr T_p(\mathscr O_Y)\otimes_{\mathscr O_Y}\mathscr M
  \longrightarrow
  \mathscr T_p(\mathscr M).
  \tag{7.7.3.1}\label{III.7.7.3.1.fr}
\]
En effet, si $Y$ est affine, cet homomorphisme a été défini en (7.2.2); cette
définition s'étend sans peine au cas général, en remarquant que si $U$, $V$ sont
deux ouverts affines de $Y$ tels que $V\subset U$, le diagramme
\oldpage[III]{199}
\[
\xymatrix@C=18pt@R=32pt{
  {\begin{gathered}
    \bigl(\mathscr T_p(\mathscr O_Y)
      \otimes_{\mathscr O_Y}\mathscr M\bigr)|U\\
    =\mathscr T^U_p(\mathscr O_Y|U)
      \otimes_{\mathscr O_Y|U}(\mathscr M|U)
  \end{gathered}}
  \ar[r]\ar[d]
  &
  {\begin{gathered}
    \mathscr T^U_p(\mathscr M|U)\\
    =\bigl(\mathscr T_p(\mathscr M)\bigr)|U
  \end{gathered}}
  \ar[d]
  \\
  {\begin{gathered}
    \bigl(\mathscr T_p(\mathscr O_Y)
      \otimes_{\mathscr O_Y}\mathscr M\bigr)|V\\
    =\mathscr T^V_p(\mathscr O_Y|V)
      \otimes_{\mathscr O_Y|V}(\mathscr M|V)
  \end{gathered}}
  \ar[r]
  &
  {\begin{gathered}
    \mathscr T^V_p(\mathscr M|V)\\
    =\bigl(\mathscr T_p(\mathscr M)\bigr)|V
  \end{gathered}}
}
\]
est commutatif par (7.2.3.3).

Pour tout morphisme $g:Y'\to Y$ on a un homomorphisme canonique
\[
  \mathscr T_p(\mathscr O_Y)\otimes_{\mathscr O_Y}\mathscr O_{Y'}
  \longrightarrow
  \mathscr T^{Y'}_p(\mathscr O_{Y'}),
  \tag{7.7.3.2}\label{III.7.7.3.2.fr}
\]
qui n'est autre que le cas particulier de (6.7.11.2) (pour les aboutissements)
dans le cas où $S=Y$, $v_1=f$, $v_2=1_Y$, $\mathscr P^{(2)}_\bullet$ réduit au
seul terme $\mathscr M$ de degré 0.

Lorsque $Y=\operatorname{Spec}(A)$, $Y'=\operatorname{Spec}(A')$ sont affines,
(7.7.3.2) n'est autre que l'homomorphisme de faisceaux correspondant à
l'homomorphisme canonique de $A'$-modules défini dans (7.2.2)
\[
  T^A_p(A)\otimes_A A'
  \longrightarrow
  T^{A'}_p(A')=T^A_p(A')
\]
comme il résulte aisément de (6.7.11) (car dans le cas envisagé, on peut prendre
$\mathfrak U'^{(i)}=u_i^{-1}(\mathfrak U^{(i)})$ dans (6.7.11)).
\end{env}

\begin{env}[7.7.4]
Lorsque $f$ est un morphisme \emph{propre}, $Y=\operatorname{Spec}(A)$ un schéma
affine noethérien et $\mathscr P_\bullet$ un complexe de
$\mathscr O_X$-Modules cohérents et $Y$-plats \emph{limité inférieurement}, on a
vu (6.10.5) que l'on peut écrire à un isomorphisme près,
\[
  \mathscr T_p(\mathscr M)
  =\mathscr H_p
    \bigl(\mathscr L_\bullet\otimes_{\mathscr O_Y}\mathscr M\bigr),
  \qquad
  \mathscr L_\bullet=\widetilde{L_\bullet},
\]
où $L_\bullet$ est un complexe de $A$-modules \emph{libres de type fini limité
inférieurement}; le foncteur $\mathscr T_\bullet$ est donc du type qui a été
étudié en détail dans (7.4) et (7.6). Nous allons traduire les résultats de cette
étude :
\end{env}

\begin{theorem}[7.7.5]
Soient $Y$ un préschéma localement noethérien, $(U_\alpha)$ un recouvrement de
$Y$ formé d'ouverts affines, $f:X\to Y$ un morphisme propre,
$\mathscr P_\bullet$ un complexe de $\mathscr O_X$-Modules cohérents et $Y$-plats
limité inférieurement. Le foncteur homologique $\mathscr T_\bullet(\mathscr M)$
défini par (7.7.1.1) possède alors les propriétés suivantes :

\medskip
\noindent
\textnormal{I) (La propriété de semi-continuité).}\footnote{Un cas particulier de
ce théorème se trouve déjà dans la note [3] de Chow-Igusa. La propriété de
semi-continuité a été découverte, dans le cadre des espaces analytiques (et sous
des hypothèses assez particulières), par Kodaira-Spencer (« On the variations of
almost-complex structures », \emph{Algebraic Geometry and Topology, A Symposium in
honor of S. Lefschetz}, Princeton Series n\textsuperscript{o} 12, p. 139--150,
Princeton, 1957) et la version générale démontrée par Grauert [5].}
\emph{La fonction}
\[
  y\longmapsto
  d_p(y)
  =
  \operatorname{rang}_{\kappa(y)}
    \mathscr T_p^{\kappa(y)}(\kappa(y))
  \tag{7.7.5.1}\label{III.7.7.5.1.fr}
\]
\emph{est semi-continue supérieurement.}

\oldpage[III]{200}
\medskip
\noindent
\textnormal{II) (La propriété d'échange).}
\emph{Pour un entier $p$ donné, les conditions suivantes sont équivalentes :}
\begin{enumerate}
  \item[a)] $\mathscr T_p$ est exact à droite.

  \item[a$'$)] $\mathscr T_p(\mathscr M)$ est isomorphe à un foncteur de la forme
    $\mathscr N\otimes_{\mathscr O_Y}\mathscr M$, $\mathscr N$ étant nécessairement
    isomorphe à
    \[
      \mathscr T_p(\mathscr O_Y)
      =
      \mathscr H^{-p}(f,\mathscr P^\bullet).
    \]

  \item[a$''$)] L'homomorphisme canonique fonctoriel (7.7.3.1) est un
    isomorphisme.

  \item[b)] $\mathscr T_{p-1}$ est exact à gauche.

  \item[b$'$)] Il existe un $\mathscr O_Y$-Module $\mathscr Q$ (nécessairement
    cohérent, et déterminé à un isomorphisme unique près) et un isomorphisme de
    foncteurs
    \[
      \mathscr T_{p-1}(\mathscr M)
      \xrightarrow{\sim}
      \mathscr{H}\!om_{\mathscr O_Y}(\mathscr Q,\mathscr M).
      \tag{7.7.5.2}\label{III.7.7.5.2.fr}
    \]

  \item[c)] En désignant par $A_\alpha$ l'anneau de l'ouvert affine $U_\alpha$,
    pour tout indice $\alpha$ le foncteur de $A_\alpha$-modules
    $T_p^{A_\alpha}$ est exact à droite.

  \item[d)] Pour tout morphisme $g:Y'\to Y$, l'homomorphisme canonique
    \[
      \mathscr T_p(\mathscr O_Y)
      \otimes_{\mathscr O_Y}\mathscr O_{Y'}
      \longrightarrow
      \mathscr T_p^{Y'}(\mathscr O_{Y'})
      \tag{7.7.5.3}\label{III.7.7.5.3.fr}
    \]
    est un isomorphisme.
\end{enumerate}
\end{theorem}

\begin{proof}
La propriété de semi-continuité est locale sur $Y$ et résulte donc de la
remarque (7.7.4) et de (7.6.9). Il est clair que a$''$) entraîne a$'$) et que
a$'$) entraîne a). L'équivalence de a), a$''$), b) et b$'$) a été démontrée
dans (7.3.1) et (7.4.6), compte tenu de la remarque (7.7.4), lorsque $Y$ est
affine. Pour passer au cas général, prouvons d'abord que a) est équivalent à c),
ce qui prouvera le caractère local sur $Y$ de la propriété a) ; la démonstration
s'appliquera également pour prouver le caractère local de a$''$) et b). Comme il
est clair que c) entraîne a), tout revient à prouver la réciproque. Il suffit
évidemment de montrer que pour tout ouvert affine $U$ de $Y$ et toute suite
exacte
\[
  0\longrightarrow\mathscr F'
  \longrightarrow\mathscr F
  \longrightarrow\mathscr F''
  \longrightarrow0
\]
de $(\mathscr O_Y|U)$-Modules quasi-cohérents, il existe une suite exacte
\[
  0\longrightarrow\mathscr G'
  \longrightarrow\mathscr G
  \longrightarrow\mathscr G''
  \longrightarrow0
\]
de $\mathscr O_Y$-Modules quasi-cohérents telle que
\[
  \mathscr F'=\mathscr G'|U,
  \qquad
  \mathscr F=\mathscr G|U,
  \qquad
  \mathscr F''=\mathscr G''|U;
\]
or, cela résulte aussitôt de l'hypothèse que $Y$ est localement noethérien, et
de (I, 9.4.2) : on prolonge en effet $\mathscr F$ en un $\mathscr O_Y$-Module
quasi-cohérent $\mathscr G$, $\mathscr F'$ en un sous-$\mathscr O_Y$-Module
$\mathscr G'$ de $\mathscr G$, et il suffit de prendre
$\mathscr G''=\mathscr G/\mathscr G'$.

Pour démontrer l'équivalence de b) et b$'$) dans le cas général, remarquons que
lorsque $Y$ est affine, on sait que $\mathscr Q$ est déterminé à un isomorphisme
unique près ; si alors $U$ est un ouvert affine du schéma affine $Y$, on en
déduit qu'il existe un isomorphisme fonctoriel
\[
  \mathscr T_{p-1}^{U}(\mathscr M|U)
  \xrightarrow{\sim}
  \mathscr{H}\!om_{\mathscr O_Y|U}
    (\mathscr Q|U,\mathscr M|U).
\]
Dans le cas général, pour tout ouvert affine $U$ de $Y$, il y a un
$(\mathscr O_Y|U)$-Module cohérent $\mathscr Q_U$ et un isomorphisme fonctoriel
\[
  \mathscr T_{p-1}^{U}(\mathscr M|U)
  \xrightarrow{\sim}
  \mathscr{H}\!om_{\mathscr O_Y|U}
    (\mathscr Q_U,\mathscr M|U);
\]
la remarque précédente montre que si $V$ est un ouvert affine contenu dans $U$,
on a $\mathscr Q_U|V=\mathscr Q_V$ ; d'où l'existence et l'unicité du
$\mathscr O_Y$-Module $\mathscr Q$ vérifiant (7.7.5.2).

Reste enfin à montrer l'équivalence de a) et d) ; il est clair que d) est de
caractère local sur $Y$, et l'on a vu ci-dessus qu'il en est de même de a) ; en
outre, d) est aussi local sur $Y'$. Or, lorsque $Y=\operatorname{Spec}(A)$,
$Y'=\operatorname{Spec}(A')$, on a vu que $T_\bullet^{A'}$ est le foncteur obtenu

\oldpage[III]{201}
à partir de $T_\bullet^A$ par extension des scalaires à $A'$, et il est clair
alors que a$'$) entraîne que (7.7.5.3) est un isomorphisme. Inversement, supposons
toujours $Y=\operatorname{Spec}(A)$ affine et soit $A'$ la $A$-algèbre
$A\oplus M$, où $M$ est un $A$-module quelconque, la multiplication dans $A'$
étant donnée par
\[
  (a_1,m_1)(a_2,m_2)
  =
  (a_1a_2,a_1m_2+a_2m_1);
\]
alors
\[
  T_p^{A'}(A')
  =
  T_p(A\oplus M)
  =
  T_p(A)\oplus T_p(M),
\]
et l'hypothèse que (7.7.5.3) soit bijectif entraîne qu'il en est de même de
l'application canonique
\[
  T_p(A)\otimes_A M\longrightarrow T_p(M),
\]
autrement dit d) entraîne a$''$), ce qui achève la démonstration.
\end{proof}

\begin{theorem}[7.7.6]
Soient $Y$ un préschéma localement noethérien, $f:X\to Y$ un morphisme propre,
$\mathscr F$ un $\mathscr O_X$-Module cohérent et $Y$-plat. Il existe alors un
$\mathscr O_Y$-Module cohérent $\mathscr Q$ (déterminé à isomorphisme unique près)
et un isomorphisme de foncteurs en le $\mathscr O_Y$-Module quasi-cohérent
$\mathscr M$ :
\[
  f_*(\mathscr F\otimes_{\mathscr O_Y}\mathscr M)
  \xrightarrow{\sim}
  \mathscr{H}\!om_{\mathscr O_Y}(\mathscr Q,\mathscr M)
  \tag{7.7.6.1}\label{III.7.7.6.1.fr}
\]
(d'où un isomorphisme de foncteurs
\[
  \Gamma(X,\mathscr F\otimes_{\mathscr O_Y}\mathscr M)
  \xrightarrow{\sim}
  \operatorname{Hom}_{\mathscr O_Y}(\mathscr Q,\mathscr M)).
  \tag{7.7.6.2}\label{III.7.7.6.2.fr}
\]
\end{theorem}

\begin{proof}
En effet, comme $\mathscr M\mapsto\mathscr F\otimes_{\mathscr O_Y}\mathscr M$ est
exact (0\textsubscript{I}, 6.7.4) et $f_*$ exact à gauche, le foncteur
$\mathscr M\mapsto f_*(\mathscr F\otimes_{\mathscr O_Y}\mathscr M)$ est exact à
gauche. Il suffit alors d'appliquer l'équivalence de (7.7.5, b)) et (7.7.5,
b$'$)) pour $p=1$.
\end{proof}

\begin{corollary}[7.7.7]
Soient $Y$ un préschéma localement noethérien, $f:X\to Y$ un morphisme propre,
$\mathscr F$, $\mathscr F'$ deux $\mathscr O_X$-Modules cohérents et $Y$-plats,
$u:\mathscr F\to\mathscr F'$ un homomorphisme. Considérons les deux foncteurs en
le $\mathscr O_Y$-Module quasi-cohérent $\mathscr M$ :
\begin{align*}
  \mathscr T(\mathscr M)
  &=
  \operatorname{Ker}
  \bigl(
    f_*(\mathscr F\otimes_{\mathscr O_Y}\mathscr M)
    \longrightarrow
    f_*(\mathscr F'\otimes_{\mathscr O_Y}\mathscr M)
  \bigr),\\
  T(\mathscr M)
  &=
  \Gamma(Y,\mathscr T(\mathscr M))\\
  &=
  \operatorname{Ker}
  \bigl(
    \Gamma(X,\mathscr F\otimes_{\mathscr O_Y}\mathscr M)
    \longrightarrow
    \Gamma(X,\mathscr F'\otimes_{\mathscr O_Y}\mathscr M)
  \bigr).
\end{align*}
Alors il existe un $\mathscr O_Y$-Module cohérent $\mathscr R$ (déterminé à
isomorphisme unique près) et des isomorphismes de foncteurs
\[
  \mathscr T(\mathscr M)
  \xrightarrow{\sim}
  \mathscr{H}\!om_{\mathscr O_Y}(\mathscr R,\mathscr M)
  \tag{7.7.7.1}\label{III.7.7.7.1.fr}
\]
\[
  T(\mathscr M)
  \xrightarrow{\sim}
  \operatorname{Hom}_{\mathscr O_Y}(\mathscr R,\mathscr M).
  \tag{7.7.7.2}\label{III.7.7.7.2.fr}
\]
\end{corollary}

\begin{proof}
On peut se borner à démontrer (7.7.7.2) ; cela prouvera en effet (7.7.7.1)
dans le cas où $Y$ est affine, et on passera de là au cas général en raisonnant
comme dans la démonstration de l'équivalence de (7.7.5, b) et b$'$)), grâce à
l'unicité à isomorphisme unique près d'un représentant d'un foncteur
représentable (0, 8.1.8). Il résulte de (7.7.6) qu'il existe deux
$\mathscr O_Y$-Modules cohérents $\mathscr Q$, $\mathscr Q'$ définissant des
isomorphismes fonctoriels
\[
  \Gamma(X,\mathscr F\otimes_{\mathscr O_Y}\mathscr M)
  \xrightarrow{\sim}
  \operatorname{Hom}_{\mathscr O_Y}(\mathscr Q,\mathscr M),
  \qquad
  \Gamma(X,\mathscr F'\otimes_{\mathscr O_Y}\mathscr M)
  \xrightarrow{\sim}
  \operatorname{Hom}_{\mathscr O_Y}(\mathscr Q',\mathscr M).
\]
Or, $u:\mathscr F\to\mathscr F'$ définit canoniquement un morphisme de foncteurs
\[
  \Gamma(X,\mathscr F\otimes_{\mathscr O_Y}\mathscr M)
  \longrightarrow
  \Gamma(X,\mathscr F'\otimes_{\mathscr O_Y}\mathscr M);
\]

\oldpage[III]{202}
il correspond à ce dernier un homomorphisme unique
$v:\mathscr Q'\to\mathscr Q$ de $\mathscr O_Y$-Modules tel que le diagramme
\[
\xymatrix@C=42pt@R=28pt{
  \Gamma(X,\mathscr F\otimes_{\mathscr O_Y}\mathscr M)
    \ar[r]\ar[d]_-{\sim}
  &
  \Gamma(X,\mathscr F'\otimes_{\mathscr O_Y}\mathscr M)
    \ar[d]^-{\sim}
  \\
  \operatorname{Hom}_{\mathscr O_Y}(\mathscr Q,\mathscr M)
    \ar[r]
  &
  \operatorname{Hom}_{\mathscr O_Y}(\mathscr Q',\mathscr M)
}
\]
soit commutatif (0, 8.1.4). Comme le foncteur contravariant
$\mathscr N\mapsto\operatorname{Hom}_{\mathscr O_Y}(\mathscr N,\mathscr M)$ est
exact à gauche dans la catégorie des $\mathscr O_Y$-Modules, il suffit de prendre
$\mathscr R=\operatorname{Coker}(v)$ pour obtenir l'isomorphisme (7.7.7.2)
cherché.
\end{proof}

\begin{corollary}[7.7.8]
Sous les hypothèses de (7.7.6) relatives à $X$, $Y$ et $f$, soient $\mathscr F$,
$\mathscr G$ deux $\mathscr O_X$-Modules cohérents vérifiant les conditions
suivantes : (i) $\mathscr F$ est $Y$-plat ; (ii) $\mathscr G$ est isomorphe au
conoyau d'un homomorphisme de $\mathscr O_X$-Modules localement libres de type
fini $\mathscr E_1\to\mathscr E_0$. Considérons les deux foncteurs en le
$\mathscr O_Y$-Module quasi-cohérent $\mathscr M$ :
\begin{align*}
  \mathscr T(\mathscr M)
  &=
  f_*
  \bigl(
    \mathscr{H}\!om_{\mathscr O_X}
      (\mathscr G,\mathscr F\otimes_{\mathscr O_Y}\mathscr M)
  \bigr),\\
  T(\mathscr M)
  &=
  \Gamma(Y,\mathscr T(\mathscr M))
  =
  \operatorname{Hom}_{\mathscr O_X}
    (\mathscr G,\mathscr F\otimes_{\mathscr O_Y}\mathscr M).
\end{align*}
Alors il existe un $\mathscr O_Y$-Module cohérent $\mathscr N$ (déterminé à
isomorphisme unique près) et des isomorphismes de foncteurs
\[
  \mathscr T(\mathscr M)
  \xrightarrow{\sim}
  \mathscr{H}\!om_{\mathscr O_Y}(\mathscr N,\mathscr M)
  \tag{7.7.8.1}\label{III.7.7.8.1.fr}
\]
\[
  T(\mathscr M)
  \xrightarrow{\sim}
  \operatorname{Hom}_{\mathscr O_Y}(\mathscr N,\mathscr M).
  \tag{7.7.8.2}\label{III.7.7.8.2.fr}
\]
\end{corollary}

\begin{proof}
En vertu de l'isomorphisme fonctoriel ($0_{\mathrm I}$, 5.4.2.1), on a des
isomorphismes fonctoriels en $\mathscr M$
\begin{align*}
  \mathscr{H}\!om_{\mathscr O_X}
    (\mathscr E_i,\mathscr F\otimes_{\mathscr O_Y}\mathscr M)
  &\xrightarrow{\sim}
  \check{\mathscr E}_i
    \otimes_{\mathscr O_X}
    (\mathscr F\otimes_{\mathscr O_Y}\mathscr M)\\
  &\xrightarrow{\sim}
  (\check{\mathscr E}_i\otimes_{\mathscr O_X}\mathscr F)
    \otimes_{\mathscr O_Y}\mathscr M\\
  &\xrightarrow{\sim}
  \mathscr{H}\!om_{\mathscr O_X}(\mathscr E_i,\mathscr F)
    \otimes_{\mathscr O_Y}\mathscr M
\end{align*}
pour $i=0,1$. Posons
$\mathscr F_i=\mathscr{H}\!om_{\mathscr O_X}(\mathscr E_i,\mathscr F)$ pour
$i=0,1$ ; ce sont des $\mathscr O_X$-Modules cohérents ($0_{\mathrm I}$, 5.3.5)
et $Y$-plats ($0_{\mathrm I}$, 5.4.2) ; soit
$u=\mathscr{H}\!om(v,1_{\mathscr F}):\mathscr F_0\to\mathscr F_1$. En vertu de
l'exactitude à gauche du foncteur
$\mathscr H\mapsto\mathscr{H}\!om_{\mathscr O_X}
(\mathscr H,\mathscr F\otimes_{\mathscr O_Y}\mathscr M)$, on a des isomorphismes
fonctoriels en $\mathscr M$
\begin{align*}
  \mathscr{H}\!om_{\mathscr O_X}
    (\mathscr G,\mathscr F\otimes_{\mathscr O_Y}\mathscr M)
  &\xrightarrow{\sim}
  \operatorname{Ker}
  \bigl(
    \mathscr{H}\!om_{\mathscr O_X}
      (\mathscr E_0,\mathscr F\otimes_{\mathscr O_Y}\mathscr M)
    \longrightarrow
    \mathscr{H}\!om_{\mathscr O_X}
      (\mathscr E_1,\mathscr F\otimes_{\mathscr O_Y}\mathscr M)
  \bigr)\\
  &\xrightarrow{\sim}
  \operatorname{Ker}
  \bigl(
    \mathscr F_0\otimes_{\mathscr O_Y}\mathscr M
    \longrightarrow
    \mathscr F_1\otimes_{\mathscr O_Y}\mathscr M
  \bigr).
\end{align*}
Puisque $f_*$ est exact à gauche, on en déduit un isomorphisme fonctoriel
\[
  f_*
  \bigl(
    \mathscr{H}\!om_{\mathscr O_X}
      (\mathscr G,\mathscr F\otimes_{\mathscr O_Y}\mathscr M)
  \bigr)
  \xrightarrow{\sim}
  \operatorname{Ker}
  \bigl(
    f_*(\mathscr F_0\otimes_{\mathscr O_Y}\mathscr M)
    \longrightarrow
    f_*(\mathscr F_1\otimes_{\mathscr O_Y}\mathscr M)
  \bigr)
\]
et il suffit alors d'appliquer (7.7.7).
\end{proof}

\oldpage[III]{203}
\begin{remarks}[7.7.9]
\textnormal{(i)} Dans (7.7.6), (7.7.7), (7.7.8), la formation des
$\mathscr O_Y$-Modules $\mathscr Q$, $\mathscr R$, $\mathscr N$ \emph{permute aux
changements de base}. Par exemple (gardant les notations de (7.7.2)), dans le
cas (7.7.6), on a, pour tout $\mathscr O_{Y'}$-Module quasi-cohérent
$\mathscr M'$, l'isomorphisme
\[
  f'_*(\mathscr F\otimes_{\mathscr O_Y}\mathscr M')
  \xrightarrow{\sim}
  \mathscr{H}\!om_{\mathscr O_{Y'}}(g^*(\mathscr Q),\mathscr M')
\]
car en vertu de la remarque faite dans (7.7.2), tout revient à voir que l'on a
\[
  \operatorname{Hom}_{\mathscr O_Y}(\mathscr Q,g_*(\mathscr M'))
  =
  \operatorname{Hom}_{\mathscr O_{Y'}}(g^*(\mathscr Q),\mathscr M')
\]
ce qui n'est autre que ($0_{\mathrm I}$, 4.4.3.1). De même, quand dans (7.7.7)
on remplace $Y$, $f$, $\mathscr M$, $\mathscr F$, $\mathscr F'$ par $Y'$, $f'$,
$\mathscr M'$, $\mathscr F\otimes_{\mathscr O_X}\mathscr O_{X'}$,
$\mathscr F'\otimes_{\mathscr O_X}\mathscr O_{X'}$, il faut remplacer
$\mathscr R$ par $g^*(\mathscr R)$. Enfin, dans (7.7.8), lorsqu'on remplace $X$,
$Y$, $f$, $\mathscr M$, $\mathscr F$ par $X'$, $Y'$, $f'$, $\mathscr M'$,
$\mathscr F\otimes_{\mathscr O_X}\mathscr O_{X'}$, et $\mathscr G$ par
$\mathscr G'=\mathscr G\otimes_{\mathscr O_Y}\mathscr O_{X'}$, il faut remplacer
$\mathscr N$ par $g^*(\mathscr N)$ : cela résulte de ce que l'on a encore une
suite exacte
\[
  \mathscr E'_1\longrightarrow
  \mathscr E'_0\longrightarrow
  \mathscr G'\longrightarrow0
  \qquad\text{avec}\qquad
  \mathscr E'_i=\mathscr E_i\otimes_{\mathscr O_Y}\mathscr O_{X'}
\]
et de ce que
\[
  \check{\mathscr E}'_i
  \otimes_{\mathscr O_{X'}}
  (\mathscr F\otimes_{\mathscr O_Y}\mathscr M')
  =
  \check{\mathscr E}_i
  \otimes_{\mathscr O_X}
  (\mathscr F\otimes_{\mathscr O_Y}\mathscr M')
  \qquad(i=0,1).
\]

\medskip
\textnormal{(ii)} La condition (ii) de l'énoncé de (7.7.8) relative à
$\mathscr G$ est toujours satisfaite pour un $\mathscr O_X$-Module cohérent
$\mathscr G$ quelconque lorsqu'il existe un $\mathscr O_X$-Module inversible
$Y$-ample, par exemple lorsque $Y$ est affine et $f:X\to Y$ est un morphisme
projectif. Il suffit alors de noter (compte tenu de (II, 5.5.1)) qu'il existe un
$\mathscr O_X$-Module localement libre de type fini $\mathscr E_0$ tel que
$\mathscr G$ soit isomorphe à un quotient de $\mathscr E_0$ (II, 2.7.10) ; comme
$\mathscr E_0$ et $\mathscr G$ sont cohérents, il en est de même du noyau
$\mathscr G_1$ de $\mathscr E_0\to\mathscr G$, et en appliquant le même résultat
à $\mathscr G_1$, on obtient bien une suite exacte
$\mathscr E_1\to\mathscr E_0\to\mathscr G\to0$ où $\mathscr E_0$ et
$\mathscr E_1$ sont localement libres de type fini.

\medskip
\textnormal{(iii)} Nous prouverons au chap. V que, dans (7.7.8), l'hypothèse
restrictive (ii) est \emph{surabondante}.
\end{remarks}

\begin{proposition}[7.7.10]
\textnormal{(Critères locaux pour la propriété d'échange).} Sous les conditions
générales de (7.7.5), soient $y$ un point de $Y$, $p$ un entier. Les propriétés
suivantes sont équivalentes :
\begin{enumerate}
  \item[a)] Le foncteur $T_p^{\mathscr O_y}$ est exact à droite.

  \item[b)] L'homomorphisme canonique
    \[
      T_p^{\mathscr O_y}(\mathscr O_y)
      \longrightarrow
      T_p^{\mathscr O_y}(\kappa(y))
    \]
    est surjectif.

  \item[c)] Pour tout entier $n$, l'homomorphisme canonique
    \[
      T_p^{\mathscr O_y}
      (\mathscr O_y/\mathfrak m_y^{n+1})
      \longrightarrow
      T_p^{\mathscr O_y}(\kappa(y))
    \]
    est surjectif.
\end{enumerate}
De plus, l'ensemble des $y\in Y$ vérifiant ces conditions est le plus grand
ensemble ouvert $U$ de $Y$ tel que $\mathscr T_p^U$ soit exact à droite.
\end{proposition}

\begin{proof}
Compte tenu de (7.7.4), l'équivalence de a), b) et c) résulte de (7.4.7) et
(7.5.2). Le fait que l'ensemble $U$ où $T_p^{\mathscr O_y}$ est exact à droite
soit ouvert est aussi conséquence de (7.4.4), et inversement si
$\mathscr T_p^V$ est exact à droite, il en est de même de
$T_p^{\mathscr O_y}$ pour tout $y\in V$, par la condition c) de (7.7.5) et
(7.6.1).
\end{proof}

\begin{corollary}[7.7.11]
Si $\mathscr T_p$ est exact à droite (resp. à gauche), alors, pour tout morphisme
$g:Y'\to Y$, $\mathscr T_p^{Y'}$ est exact à droite (resp. à gauche). La
réciproque est vraie lorsque le morphisme $g$ est fidèlement plat.
\end{corollary}

\oldpage[III]{204}
\begin{proof}
La première assertion est conséquence immédiate de (7.6.1) et du fait que la
question est locale sur $Y$ et $Y'$, d'après (7.7.5, c) et b)). Pour démontrer
la seconde assertion, il suffit de voir que pour tout $y\in Y$,
$T_p^{\mathscr O_y}$ est exact à droite (resp. à gauche), en vertu de (7.7.10).
Mais par hypothèse, il existe $y'\in Y'$ tel que $g(y')=y$, et
$\mathscr O_{y'}$ est un $\mathscr O_y$-module fidèlement plat ; la conclusion
résulte alors de l'hypothèse et de (7.6.1).
\end{proof}

\begin{remarks}[7.7.12]
\textnormal{(i)} Sous les hypothèses de (7.7.4), supposons en outre que
$\mathscr P_\bullet$ soit un complexe \emph{fini} ; alors il résulte de (7.7.1)
(puisqu'on peut prendre pour $\mathfrak U$ un recouvrement ouvert affine
\emph{fini} de $X$) que le bicomplexe
$C^\bullet(\mathfrak U,\mathscr P^\bullet\otimes_{\mathscr O_Y}\mathscr M)$ est
aussi fini, et de façon précise qu'il existe un ensemble fini $E$ de couples,
\emph{indépendant de $\mathscr M$}, tel que
\[
  C^h(\mathfrak U,\mathscr P^k\otimes_{\mathscr O_Y}\mathscr M)=0
  \qquad\text{pour tous les couples }(h,k)\notin E.
\]
On en conclut qu'il existe $i_1$ tel que, pour $i\geq i_1$, on ait
$\mathscr T_i(\mathscr M)=0$ pour \emph{tout} $\mathscr O_Y$-Module
quasi-cohérent $\mathscr M$. En particulier, $\mathscr T_i$ est trivialement un
foncteur exact en $\mathscr M$ pour ces valeurs de $i$, et par suite (7.4.1),
$Z'_i(L_\bullet)$ est un $A$-module \emph{plat} de type fini (donc
\emph{projectif} de type fini, puisque $A$ est noethérien) pour ces valeurs de
$i$. Considérons alors le complexe $(L'_\bullet)$ de $A$-modules tel que
$L'_i=L_i$ pour $i<i_1$, $L'_{i_1}=Z'_{i_1}(L_\bullet)$ et $L'_i=0$ pour
$i>i_1$ et soit $\mathscr L'_\bullet=\widetilde{L'_\bullet}$. Il est clair que
\[
  \mathscr H_i(\mathscr L'_\bullet\otimes_{\mathscr O_Y}\mathscr M)
  =
  \mathscr H_i(\mathscr L_\bullet\otimes_{\mathscr O_Y}\mathscr M)
\]
pour $i<i_1-1$ et aussi pour $i\geq i_1$ (les deux membres étant alors nuls) ;
enfin, comme
$\operatorname{Im}(Z'_{i_1}\otimes_A M)=\operatorname{Im}(L_{i_1}\otimes_A M)$
par définition, on a aussi
$\mathscr H_i(\mathscr L'_\bullet\otimes_{\mathscr O_Y}\mathscr M)
=\mathscr H_i(\mathscr L_\bullet\otimes_{\mathscr O_Y}\mathscr M)$ pour
$i=i_1-1$. On voit donc qu'on peut supposer dans (7.7.4) que
$\mathscr L_\bullet$ est aussi un complexe \emph{fini}, à condition d'exiger
seulement que les $\mathscr L_i$ soient des $\mathscr O_Y$-Modules
\emph{localement libres} (associés à des $A$-modules projectifs de type fini).

Ce raisonnement s'applique en particulier au cas où $\mathscr P_\bullet$ est
réduit à \emph{un seul} terme $\mathscr F\ne0$, de degré 0 (auquel cas
$\mathscr T_n(\mathscr M)=\mathrm R^{-n}f_*
(\mathscr F\otimes_{\mathscr O_Y}\mathscr M)$) ; on peut alors supposer que les
$\mathscr L_i$ sont nuls pour $i>0$ ; on utilisera de préférence dans ce cas les
notations cohomologiques, écrivant donc $\mathscr T^{-p}$ au lieu de
$\mathscr T_p$.

\medskip
\textnormal{(ii)} Lorsque dans l'énoncé de (7.7.5), on ne suppose plus que les
$\mathscr P_i$ sont $Y$-plats, les conclusions restent valables à condition de
poser cette fois
\[
  \mathscr T_p(\mathscr M)
  =
  \operatorname{Tor}^Y_p
    (f,1_Y;\mathscr P_\bullet,\mathscr M).
  \tag{7.7.12.1}\label{III.7.7.12.1.fr}
\]
En effet, $\operatorname{Tor}^Y_n(f,1_Y;\mathscr P_\bullet,\mathscr O_Y)$ est
alors un $\mathscr O_Y$-Module \emph{cohérent} en vertu de (6.7.9). La
démonstration de (6.10.5) s'applique sans changement, compte tenu de (6.10.1)
et montre encore que lorsque $Y=\operatorname{Spec}(A)$ est affine, l'on a
\[
  \mathscr T_p(\mathscr M)
  =
  \mathscr H_p(\mathscr L_\bullet\otimes_{\mathscr O_Y}\mathscr M),
  \qquad
  \mathscr L_\bullet=\widetilde{L_\bullet},
\]
où $L_\bullet$ est un complexe de $A$-modules libres de type fini ; cela prouve
notre assertion.
\end{remarks}

\subsection{Application aux morphismes propres : II. Critères de platitude
cohomologique}
\label{subsection:III.7.8.fr}

\begin{definition}[7.8.1]
Soient $X$, $Y$ deux préschémas, $f:X\to Y$ un morphisme quasi-compact et
séparé, $\mathscr P_\bullet$ un complexe de $\mathscr O_X$-Modules
quasi-cohérents et $Y$-plats, $\mathscr T_\bullet$ le foncteur homologique de
$\mathscr O_Y$-Modules quasi-cohérents défini par (7.7.1.1), $y$ un point de $Y$.
On dit que $\mathscr P_\bullet$ est

\oldpage[III]{205}
\emph{homologiquement plat sur $Y$ au point $y$, en dimension $p$} (ou
\emph{cohomologiquement plat sur $Y$ au point $y$, en dimension $-p$}) s'il
existe un voisinage ouvert $U$ de $y$ dans $Y$ tel que
$\mathscr T_p^U=\mathscr T_U^{-p}$ soit exact. On dit que $\mathscr P_\bullet$ est
\emph{homologiquement plat en dimension $p$ sur $Y$} (ou
\emph{cohomologiquement plat en dimension $-p$ sur $Y$}) s'il est
homologiquement plat sur $Y$ en tout point $y\in Y$, en dimension $p$.

Lorsque $\mathscr P_\bullet$ est homologiquement plat sur $Y$ (resp. sur $Y$ au
point $y$) pour \emph{toute dimension $p$}, on dit simplement que
$\mathscr P_\bullet$ est \emph{homologiquement plat sur $Y$} (resp. sur $Y$ au
point $y$) ou \emph{cohomologiquement plat sur $Y$} (resp. sur $Y$ au point
$y$).
\end{definition}

\begin{env}[7.8.2]
Par définition, la notion de platitude homologique sur $Y$ est \emph{locale sur
$Y$}. Si $Y$ est localement noethérien, ou un \emph{schéma}, pour que
$\mathscr P_\bullet$ soit homologiquement plat sur $Y$ en dimension $p$, il faut
et il suffit que le foncteur $\mathscr T_p$ soit \emph{exact} : la démonstration
a été faite dans le cas où $Y$ est localement noethérien au cours de la
démonstration de (7.7.5) ; le raisonnement est le même (basé sur (I, 9.4.2)
appliqué à un ouvert affine dans un schéma quasi-compact) lorsque $Y$ est un
schéma.
\end{env}

\begin{proposition}[7.8.3]
Les notations et hypothèses étant celles de (7.8.1), les conditions suivantes
sont équivalentes :
\begin{enumerate}
  \item[a)] $\mathscr P_\bullet$ est homologiquement plat sur $Y$ au point $y$
    en dimension $p$.
  \item[b)] Il existe un voisinage ouvert $U$ de $y$ dans $Y$ tel que
    $\mathscr T_p^U$ et $\mathscr T_{p+1}^U$ soient exacts à droite.
  \item[c)] Il existe un voisinage ouvert $U$ de $y$ dans $Y$ tel que
    $\mathscr T_p^U$ et $\mathscr T_{p-1}^U$ soient exacts à gauche.
  \item[d)] Il existe un voisinage ouvert $U$ de $y$ dans $Y$ tel que
    $\mathscr T_{p+1}^U$ soit exact à droite et $\mathscr T_{p-1}^U$ exact à
    gauche.
\end{enumerate}
\end{proposition}

\begin{proof}
Compte tenu de l'interprétation de $\mathscr T_p$ lorsque $Y$ est affine, cela
n'est qu'une traduction d'une partie de (7.3.3).
\end{proof}

\begin{proposition}[7.8.4]
Soient $Y$ un préschéma localement noethérien, $f:X\to Y$ un morphisme propre,
$\mathscr P_\bullet$ un complexe de $\mathscr O_X$-Modules cohérents et $Y$-plats
limité inférieurement, $\mathscr T_\bullet$ le foncteur défini par (7.7.1.1).
Pour tout $y\in Y$, les conditions suivantes sont équivalentes :
\begin{enumerate}
  \item[a)] $\mathscr P_\bullet$ est homologiquement plat sur $Y$ en $y$ en
    dimension $p$.
  \item[b)] Le foncteur $T_p^{\mathscr O_y}$ est exact.
  \item[c)] Il existe un entier $n_0$ tel que pour $n\geq n_0$, on ait
    \[
      \operatorname{long}T_p^{\mathscr O_y}
        (\mathscr O_y/\mathfrak m_y^{n+1})
      =
      \operatorname{long}T_p^{\mathscr O_y}(\kappa(y))
      \cdot
      \operatorname{long}(\mathscr O_y/\mathfrak m_y^{n+1})
      \tag{7.8.4.1}\label{III.7.8.4.1.fr}
    \]
    (où il s'agit de longueurs de $\mathscr O_y$-modules).
  \item[d)] Il y a un voisinage ouvert $U$ de $y$ tel que
    $(\mathscr H^{-p}(f,\mathscr P^\bullet))|U$ soit isomorphe à un
    $(\mathscr O_Y|U)$-Module de la forme $(\mathscr O_Y|U)^m$ et que, pour tout
    $(\mathscr O_Y|U)$-Module quasi-cohérent $\mathscr M$, l'homomorphisme
    canonique
    \[
      ((\mathscr H^{-p}(f,\mathscr P^\bullet))|U)
      \otimes_{\mathscr O_Y|U}\mathscr M
      \longrightarrow
      \mathscr H^{-p}
      \bigl(
        f,(\mathscr P^\bullet|U)
          \otimes_{\mathscr O_Y|U}\mathscr M
      \bigr)
      \tag{7.8.4.2}\label{III.7.8.4.2.fr}
    \]
    soit bijectif.
\end{enumerate}
Lorsque ces conditions sont vérifiées, on a en outre la propriété suivante :
\begin{enumerate}
  \item[e)] Il existe un voisinage de $y$ dans lequel la fonction
    $z\mapsto d_p(z)$ (définie dans (7.7.5.1)) est constante.
\end{enumerate}
De plus, si $Y$ est réduit au point $y$ ($0_{\mathrm I}$, 4.1.4), e) est
équivalente aux autres conditions.
\end{proposition}

\oldpage[III]{206}
\begin{proof}
En effet, la condition b) équivaut à dire que $T_p^{\mathscr O_y}$ et
$T_{p+1}^{\mathscr O_y}$ sont exacts à droite (7.3.3). L'équivalence de a) et b)
résulte alors de (7.7.10) et (7.8.3). Comme
$\mathscr O_y/\mathfrak m_y^{n+1}$ est artinien, et que
$T_p^{\mathscr O_y}(\mathscr O_y/\mathfrak m_y^{n+1})$ et
$T_p^{\mathscr O_y}(\kappa(y))$ sont des
$(\mathscr O_y/\mathfrak m_y^{n+1})$-modules de type fini (7.7.4), donc de
longueur finie, l'équivalence de b) et c) résulte encore de (7.7.10) et de
(7.3.5.7). Le fait que a) entraîne e), et lui est équivalente lorsque
$\mathscr O_y$ est réduit, est conséquence de (7.6.9). Enfin, a) entraîne que
(7.8.4.2) est bijectif en vertu de la définition (7.8.1) et de (7.7.5) ; d'autre
part, a) entraîne que $(\mathscr H^{-p}(f,\mathscr P^\bullet))_y$ est un
$\mathscr O_y$-module plat (7.3.3, f)), donc libre ($0_{\mathrm I}$, 10.1.3),
puisqu'il s'agit d'un $\mathscr O_y$-module de type fini en vertu de (7.7.4) ;
puisque $\mathscr H^{-p}(f,\mathscr P^\bullet)$ est un $\mathscr O_Y$-Module
cohérent (7.7.4), il est localement libre dans un voisinage de $y$
($0_{\mathrm I}$, 5.2.7). Inversement, il est clair que d) entraîne a) par
définition du foncteur $\mathscr T_p^U$ (7.7.2.2).
\end{proof}

\begin{proposition}[7.8.5]
Sous les hypothèses de (7.8.4), les conditions suivantes sont équivalentes :
\begin{enumerate}
  \item[a)] $\mathscr P_\bullet$ est homologiquement plat sur $Y$ en toute
    dimension $i\leq p$.
  \item[b)] Pour $i\leq p+1$ les foncteurs $\mathscr T_i$ sont exacts à droite.
  \item[c)] Pour $i\geq-p$, les $\mathscr O_Y$-Modules
    $\mathscr H^i(f,\mathscr P^\bullet)$ sont localement libres.
\end{enumerate}
\end{proposition}

\begin{proof}
L'équivalence de a) et b) est triviale (7.8.3) et a) entraîne c) en vertu de
(7.8.4). Inversement, supposons c) vérifiée ; notons d'autre part qu'on a
$\mathscr L_i=0$ pour $i<i_0$ (7.7.4), donc aussi $\mathscr T_i=0$ pour
$i\leq i_0$. Tout point $y\in Y$ a donc un voisinage affine
$U=\operatorname{Spec}(A)$ tel que $T_i^A(A)$ soit un $A$-module libre pour
$i\leq p$ ; en vertu de (7.3.7), on en conclut que
$\mathscr T_i^U=T_i^A$ est exact pour $i\leq p$.
\end{proof}

Nous allons surtout appliquer les critères de platitude cohomologique au cas où
le complexe $\mathscr P_\bullet$ est réduit à un \emph{seul}
$\mathscr O_X$-Module cohérent $\mathscr F$ \emph{plat} sur $Y$, pris égal à
$\mathscr P_0$ ; rappelons que l'on a alors
$\mathscr T_p(\mathscr M)=\mathrm R^{-p}f_*
(\mathscr F\otimes_{\mathscr O_Y}\mathscr M)$.

\begin{proposition}[7.8.6]
Soient $Y$ un préschéma localement noethérien, $f:X\to Y$ un morphisme propre
et plat, $y$ un point de $Y$ ; désignons par $X_y$ la fibre
$f^{-1}(y)=X\otimes_Y\kappa(y)$. Supposons que
$\Gamma(X_y,\mathscr O_{X_y})=R$ soit une $\kappa(y)$-algèbre séparable
(Bourbaki, \emph{Alg.}, chap. VIII, § 7, n\textsuperscript{o} 5) autrement dit
composée d'un nombre fini d'extensions séparables de degré fini de $\kappa(y)$.
Alors $\mathscr O_X$ est cohomologiquement plat sur $Y$ au point $y$ en
dimension 0.
\end{proposition}

\begin{proof}
En vertu de (7.8.4), on peut se borner au cas où $Y$ est le spectre de l'anneau
local $A=\mathscr O_y$ ; l'hypothèse que $f$ est plat entraîne
$\mathscr T_{-1}=0$, donc on voit déjà que $T_0^A$ est exact à gauche et tout
revient à voir qu'il est exact à droite ; en vertu de (7.7.10, c)), on est même
ramené au cas où $A=\mathscr O_y$ est \emph{artinien}. Soit $k'$ une extension
finie de $\kappa(y)$ qui soit un corps neutralisant de $R$, de sorte que
$R\otimes_{\kappa(y)}k'$ est composée directe d'un nombre fini de corps
isomorphes à $k'$. On sait qu'il existe un homomorphisme local de $A$ dans un
anneau local $A'$, faisant de $A'$ une $A$-algèbre \emph{libre} finie sur $A$,
et tel que le corps résiduel de $A'$ soit isomorphe à $k'$ ($0_{\mathrm I}$,
10.3.2). En vertu de (7.6.1), on est ramené à prouver que $T_0^{A'}$ est exact
à droite, autrement dit on peut supposer que $R$ est composé direct de $m$
corps isomorphes à $\kappa(y)$. Notons maintenant le lemme élémentaire
suivant :

\begin{lemma}[7.8.6.1]
Soit $Z$ un espace annelé en anneaux locaux ; pour que $Z$ soit connexe,

\oldpage[III]{207}
il faut et il suffit que l'anneau $\Gamma(Z,\mathscr O_Z)$ ne soit pas un produit
de deux anneaux non réduits à 0.
\end{lemma}

\begin{proof}
Il est clair en effet que si $Z$ est réunion de deux ouverts non vides disjoints,
$\Gamma(Z,\mathscr O_Z)$ est isomorphe au produit des deux anneaux
$\Gamma(Z_1,\mathscr O_Z)$ et $\Gamma(Z_2,\mathscr O_Z)$ non réduits à 0.
Inversement, dire que $\Gamma(Z,\mathscr O_Z)$ est un tel produit équivaut à
dire qu'il y a dans $\Gamma(Z,\mathscr O_Z)$ un idempotent $s$ distinct de 0 et
de 1 ; pour tout $z\in Z$, $s_z$ est alors un idempotent dans $\mathscr O_z$,
donc égal à 0 ou 1. Mais il est clair que l'ensemble des $z$ tels que $s_z=0$
est ouvert ; d'autre part, si $s_z=1$, on a par définition $s(z)\ne0$, donc
l'ensemble des $z$ où $s_z=1$ est aussi ouvert ($0_{\mathrm I}$, 5.5.2) ; d'où
la conclusion.
\end{proof}

Il résulte de ce lemme que $X_y$ a exactement $m$ composantes connexes $X'_i$
et que $\Gamma(X'_i,\mathscr O_{X'_i})=\kappa(y)$ pour tout $i$. Comme $A$ a été
supposé local et artinien, son spectre est réduit à un point, donc $X$ et $X_y$
ont même espace sous-jacent ; $X$ a donc $m$ composantes connexes $X_i$ telles
que $X'_i=X_i\otimes_Y\kappa(y)$. On est ainsi finalement ramené au cas où
$R=\kappa(y)$ ; en vertu de (7.7.10, b)), on est ramené à prouver que
l'homomorphisme canonique
\[
  \Gamma(X,\mathscr O_X)
  \longrightarrow
  \Gamma(X_y,\mathscr O_{X_y})
\]
est surjectif ; mais cela est trivial, car le composé
\[
  \Gamma(Y,\mathscr O_Y)=A
  \longrightarrow
  \Gamma(X,\mathscr O_X)
  \longrightarrow
  \Gamma(X_y,\mathscr O_{X_y})=\kappa(y)
\]
est déjà surjectif.
\end{proof}

\begin{corollary}[7.8.7]
Sous les hypothèses de (7.8.6), il existe un voisinage ouvert $U$ de $y$ tel
que :
\begin{enumerate}
  \item[(i)] $f_*(\mathscr O_X)|U$ soit isomorphe à un
    $(\mathscr O_Y|U)$-Module de la forme $(\mathscr O_Y|U)^m$.
  \item[(ii)] Pour tout $z\in U$, l'homomorphisme canonique
    \[
      (f_*(\mathscr O_X))_z
      \otimes_{\mathscr O_z}\kappa(z)
      \longrightarrow
      \Gamma(X_z,\mathscr O_{X_z})
    \]
    est bijectif.
\end{enumerate}
\end{corollary}

\begin{proof}
(i) résulte de (7.8.6) et (7.8.4).

(ii) résulte de ce que $\mathscr T_0^U$ est exact (pour $U$ convenablement
choisi), et de (7.7.5.3).
\end{proof}

\begin{corollary}[7.8.8]
Supposons vérifiées les conditions de (7.8.6) et en outre que
$\Gamma(X_y,\mathscr O_{X_y})=\kappa(y)$. Alors il existe un voisinage ouvert
$U$ de $y$ tel que l'homomorphisme canonique
$\mathscr O_Y|U\to f_*(\mathscr O_X)|U$ soit bijectif.
\end{corollary}

\begin{proof}
En effet, il résulte de (7.8.7, (ii)) que l'entier $m$ figurant dans (7.8.7,
(i)) est nécessairement égal à 1.
\end{proof}

\begin{corollary}[7.8.9]
Sous les hypothèses de (7.8.6), il existe un voisinage ouvert $U$ de $y$, un
$\mathscr O_U$-Module cohérent $\mathscr Q$ (déterminé à isomorphisme unique
près) et un isomorphisme de foncteurs en le $\mathscr O_U$-Module quasi-cohérent
$\mathscr M$ :
\[
  \mathrm R^1f_*(f^*(\mathscr M))
  \xrightarrow{\sim}
  \mathscr{H}\!om_{\mathscr O_U}(\mathscr Q,\mathscr M).
  \tag{7.8.9.1}\label{III.7.8.9.1.fr}
\]
\end{corollary}

\begin{proof}
En effet, l'hypothèse entraîne que $\mathscr T_0^U$ est exact pour un $U$
convenable ; il suffit donc d'appliquer l'équivalence de (7.7.5, a)) et
(7.7.5, b$'$)) au cas $p=0$ et en prenant pour $\mathscr P_\bullet$ le complexe
réduit à son terme de degré 0 égal à $\mathscr O_X$.
\end{proof}

\begin{remarks}[7.8.10]
\textnormal{(i)} Sous les conditions de (7.8.6), considérons la factorisation de
Stein de $f$ (4.3.3)
\[
\xymatrix{
  X \ar[r]^{f'} & Y' \ar[r]^{g} & Y .
}
\]

\oldpage[III]{208}
avec $Y'=\operatorname{Spec}(f_*(\mathscr O_X))$ ; le morphisme fini $g$ est alors
tel que $g_*(\mathscr O_{Y'})=f_*(\mathscr O_X)$ soit localement libre au
voisinage de $y$, et sa fibre en $y$ est le spectre d'une algèbre séparable sur
$\kappa(y)$ (II, 1.5.1). Nous en déduirons au chap. IV qu'il y a un voisinage
ouvert $U$ de $y$ dans $Y$ tel que pour la restriction $g^{-1}(U)\to U$ de $g$,
\emph{toute} fibre $g^{-1}(z)$ (où $z\in U$) soit spectre d'une algèbre séparable
sur $\kappa(z)$ (c'est ce que nous appellerons un \emph{revêtement étale} de
$U$) ; il résultera alors de (7.8.7, (ii)) que l'hypothèse faite sur le point
$y$ dans (7.8.6) est vérifiée aussi en tous les points d'un voisinage de $y$.

\medskip
\textnormal{(ii)} Nous verrons au chap. V que, même si $X$ est projectif sur $Y$
(et même s'il est en outre « simple » sur $Y$, propriété qui sera définie au
chap. IV), le $\mathscr O_U$-Module $\mathscr Q$ de (7.8.9) n'est pas
nécessairement localement libre ; en d'autres termes, $\mathscr O_X$ (sous ces
conditions) n'est pas nécessairement cohomologiquement plat en \emph{dimension
1} sur $Y$ au point $y$. Au chap. V, nous interpréterons $\mathscr Q$ comme le
faisceau des 1-différentielles du schéma de Picard de $X$ par rapport à $Y$ le
long de la section unité.
\end{remarks}

\subsection{Application aux morphismes propres : III. Invariance de la
caractéristique d'Euler-Poincaré et du polynôme de Hilbert}
\label{subsection:III.7.9.fr}
\label{III.7.9.fr}

\begin{env}[7.9.1]
\label{III.7.9.1.fr}
Soient $A$ un anneau, $M$ un $A$-module projectif de type fini ; rappelons
(Bourbaki, \emph{Alg. comm.}, chap. II, \S~5, n\textsuperscript{o}~2) qu'il
revient au même de dire que le $\mathscr O_X$-Module associé $\widetilde M$ sur
$X=\operatorname{Spec}(A)$ est \emph{localement libre de type fini}. Pour tout
$\mathfrak p\in\operatorname{Spec}(A)$ on appelle \emph{rang de $M$ en
$\mathfrak p$} et on note $\operatorname{rang}_{\mathfrak p}(M)$ le rang du
$A_{\mathfrak p}$-module libre $M_{\mathfrak p}$ (ou encore le rang en
$\mathfrak p$ du $\mathscr O_X$-Module localement libre $\widetilde M$). On a
donc
\[
  \operatorname{rang}_{\mathfrak p}M
  =\operatorname{rang}_{\mathfrak p}(M_{\mathfrak p})
  =\operatorname{rang}_{\kappa(\mathfrak p)}
  \bigl(M\otimes_A\kappa(\mathfrak p)\bigr).
  \tag{7.9.1.1}\label{III.7.9.1.1.fr}
\]
\end{env}

\begin{proposition}[7.9.2]
\label{III.7.9.2.fr}
Soit $P_\bullet$ un complexe fini de $A$-modules projectifs de type fini, et
pour tout $A$-module $M$, soit
$T_\bullet(M)=H_\bullet(P_\bullet\otimes_A M)$. Alors, pour tout
$\mathfrak p\in\operatorname{Spec}(A)$, on a
\[
  \sum_i(-1)^i\operatorname{rang}_{\kappa(\mathfrak p)}
  T_i\bigl(\kappa(\mathfrak p)\bigr)
  =\sum_i(-1)^i\operatorname{rang}_{\mathfrak p}(P_i).
  \tag{7.9.2.1}\label{III.7.9.2.1.fr}
\]
\end{proposition}

\begin{proof}
En effet, on a par définition
$T_i(\kappa(\mathfrak p))=H_i(P_\bullet\otimes_A\kappa(\mathfrak p))$ et,
compte tenu de (7.9.1.1), la formule (7.9.1.2) n'est autre que l'invariance de
la caractéristique d'Euler-Poincaré d'un complexe fini d'espaces vectoriels de
dimension finie par passage à l'homologie ($0_{\mathrm I}$, 11.10.2).
\end{proof}

\begin{corollary}[7.9.3]
\label{III.7.9.3.fr}
La fonction
\[
  \mathfrak p\longmapsto
  \sum_i(-1)^i\operatorname{rang}_{\kappa(\mathfrak p)}
  T_i\bigl(\kappa(\mathfrak p)\bigr)
\]
est localement constante dans $\operatorname{Spec}(A)$.
\end{corollary}

\begin{theorem}[7.9.4]
\label{III.7.9.4.fr}
Soient $Y$ un préschéma localement noethérien, $f:X\to Y$ un morphisme propre,
$\mathscr P_\bullet$ un complexe fini de $\mathscr O_X$-Modules cohérents et
$Y$-plats. Si l'on pose
\[
  \mathscr T_\bullet(\mathscr M)
  =\mathscr H^{-\bullet}
  \bigl(f,\mathscr P^\bullet\otimes_{\mathscr O_Y}\mathscr M\bigr)
  \qquad\text{(cf. (7.7.1.1))}
\]
la fonction

\oldpage[III]{209}
\[
  y\longmapsto
  \sum_i(-1)^i\operatorname{rang}_{\kappa(y)}
  \mathscr T_i\bigl(\kappa(y)\bigr)
  \tag{7.9.4.1}\label{III.7.9.4.1.fr}
\]
est localement constante dans $Y$.
\end{theorem}

\begin{proof}
On peut se borner au cas où $Y=\operatorname{Spec}(A)$ est affine d'anneau $A$
noethérien. Comme le complexe $\mathscr P_\bullet$ est fini, on sait (7.7.12,
(i)) qu'on a
\[
  \mathscr T_p(\mathscr M)
  =\mathscr H_p(\mathscr L_\bullet\otimes_{\mathscr O_Y}\mathscr M),
  \qquad \mathscr L_\bullet=\widetilde{L_\bullet},
\]
$L_\bullet$ étant un complexe \emph{fini} de $A$-modules projectifs de type
fini. Le théorème résulte alors de (7.9.3).
\end{proof}

\begin{env}[7.9.5]
\label{III.7.9.5.fr}
Sous les conditions de (7.9.4), la fonction (7.9.4.1) est \emph{constante}
lorsque $Y$ est \emph{connexe}. Lorsque $Y$ est connexe et non vide, on désigne
la valeur unique (entière) de (7.9.4.1) par
\[
  \operatorname{EP}(f,\mathscr P_\bullet)
  \quad\text{ou}\quad
  \operatorname{EP}(Y,\mathscr P_\bullet),
  \quad\text{ou simplement}\quad
  \operatorname{EP}(\mathscr P_\bullet)
\]
s'il ne peut en résulter de confusion, et l'on dit que cet entier est la
\emph{caractéristique d'Euler-Poincaré de $\mathscr P_\bullet$} relativement à
$f$ (ou à $Y$). Dans le cas général, on notera aussi
\[
  \operatorname{EP}(f,\mathscr P_\bullet;y)
  \quad\text{ou}\quad
  \operatorname{EP}(Y,\mathscr P_\bullet;y)
  \quad\text{ou}\quad
  \operatorname{EP}(\mathscr P_\bullet;y)
\]
le second membre de (7.9.4.1).
\end{env}

\begin{env}[7.9.6]
\label{III.7.9.6.fr}
Sous les hypothèses de (7.9.4) relativement à $X$, $Y$ et $f$, soit
\[
  0\longrightarrow\mathscr P'_\bullet
  \xrightarrow{u}\mathscr P_\bullet
  \xrightarrow{v}\mathscr P''_\bullet
  \longrightarrow0
\]
une suite exacte de complexes finis de $\mathscr O_X$-Modules cohérents et
$Y$-plats, les homomorphismes $u$ et $v$ étant de degrés \emph{pairs} $2d$,
$2d'$ respectivement. Comme $\mathscr T_\bullet$ est un foncteur homologique
(7.7.1), on a une suite exacte d'homologie
\[
  \cdots\longrightarrow
  \mathscr T_i(\mathscr P'_\bullet,\kappa(y))
  \longrightarrow
  \mathscr T_{i+2d}(\mathscr P_\bullet,\kappa(y))
  \longrightarrow
  \mathscr T_{i+2d+2d'}(\mathscr P''_\bullet,\kappa(y))
  \longrightarrow
  \mathscr T_{i-1}(\mathscr P'_\bullet,\kappa(y))
  \longrightarrow\cdots
\]
n'ayant d'ailleurs qu'un nombre fini de termes. En écrivant que la
caractéristique d'Euler-Poincaré de ce complexe est nulle ($0_{\mathrm I}$,
11.10.1), il vient aussitôt
\[
  \operatorname{EP}(\mathscr P_\bullet;y)
  =\operatorname{EP}(\mathscr P'_\bullet;y)
  +\operatorname{EP}(\mathscr P''_\bullet;y)
  \tag{7.9.6.1}\label{III.7.9.6.1.fr}
\]
pour tout $y\in Y$. Or, si par exemple $\mathscr P_\bullet=(\mathscr P_i)$ avec
$\mathscr P_i=0$ pour $i<0$, on a la suite exacte de complexes
\[
\xymatrix@C=2.4em{
  \cdots \ar[r] & 0 \ar[r]\ar[d] & 0 \ar[r]\ar[d]
    & 0 \ar[r]\ar[d] & 0 \ar[r] & \cdots \\
  \cdots \ar[r] & 0 \ar[r]\ar[d] & 0 \ar[r]\ar[d]
    & \mathscr P_1 \ar[r]\ar[d] & \mathscr P_2 \ar[r] & \cdots \\
  \cdots \ar[r] & 0 \ar[r]\ar[d] & \mathscr P_0 \ar[r]\ar[d]
    & \mathscr P_1 \ar[r]\ar[d] & \mathscr P_2 \ar[r] & \cdots \\
  \cdots \ar[r] & 0 \ar[r]\ar[d] & \mathscr P_0 \ar[r]\ar[d]
    & 0 \ar[r]\ar[d] & 0 \ar[r] & \cdots \\
  \cdots \ar[r] & 0 \ar[r] & 0 \ar[r]
    & 0 \ar[r] & 0 \ar[r] & \cdots
}
\]
les flèches verticales non nulles étant les automorphismes identiques ; on peut
appliquer (7.9.6.1) à cette suite exacte, d'où, par récurrence sur la longueur de
$\mathscr P_\bullet$, la formule
\[
  \operatorname{EP}(\mathscr P_\bullet;y)
  =\sum_i(-1)^i\operatorname{EP}(\mathscr P_i;y).
  \tag{7.9.6.2}\label{III.7.9.6.2.fr}
\]

\oldpage[III]{210}
où, pour tout $\mathscr O_X$-Module cohérent $\mathscr F$, plat sur $Y$, on
désigne par $\operatorname{EP}(\mathscr F;y)$ (ou
$\operatorname{EP}(f,\mathscr F;y)$ ou $\operatorname{EP}(Y,\mathscr F;y)$) la
fonction $\operatorname{EP}(\mathscr Q_\bullet;y)$ correspondant au complexe
$\mathscr Q_\bullet$ dont le seul terme $\ne0$ est de degré $0$ et égal à
$\mathscr F$. On voit donc qu'on peut se ramener à étudier les caractéristiques
d'Euler-Poincaré de complexes réduits à un seul terme.
\end{env}

\begin{proposition}[7.9.7]
\label{III.7.9.7.fr}
Sous les hypothèses de (7.9.4), soient $Y'$ un préschéma localement noethérien,
$g:Y'\to Y$ un morphisme,
\[
  X'=X\times_Y Y',
  \qquad f'=f_{(Y')}:X'\longrightarrow Y',
\]
$\mathscr P'_\bullet$ le complexe fini
\[
  \mathscr P_\bullet\otimes_{\mathscr O_Y}\mathscr O_{Y'}
\]
de $\mathscr O_{X'}$-Modules ; $\mathscr P'_\bullet$ est formé de
$\mathscr O_{X'}$-Modules cohérents et $Y'$-plats, et pour tout $y'\in Y'$, on a
\[
  \operatorname{EP}(\mathscr P'_\bullet;y')
  =\operatorname{EP}(\mathscr P_\bullet;g(y')).
  \tag{7.9.7.1}\label{III.7.9.7.1.fr}
\]
\end{proposition}

\begin{proof}
Les $\mathscr O_{X'}$-Modules $\mathscr P'_i$ étant images réciproques des
$\mathscr P_i$ par la projection $X'\to X$ sont cohérents, ils sont $Y'$-plats en
vertu de ($0_{\mathrm I}$, 6.2.1) et (I, 4.14.5), la question étant locale sur
$X$, $Y$ et $Y'$ ; enfin, on sait que $f'$ est propre (II, 5.4.2), donc le
premier membre de (7.9.7.1) est défini. La formule (7.9.7.1) résulte alors de
(6.10.4.2), (7.7.2) et du lemme (7.6.7), en se ramenant, comme on peut toujours
le faire, au cas où $Y$ et $Y'$ sont affines.
\end{proof}

\begin{proposition}[7.9.8]
\label{III.7.9.8.fr}
Supposons vérifiées les hypothèses de (7.9.4) et en outre qu'il existe un entier
$i_0$ tel que
\[
  \mathscr T_i(\kappa(y))=0
  \qquad\text{pour $i\ne i_0$ et tout $y\in Y$.}
\]
Alors
\[
  \mathscr T_{i_0}(\mathscr O_Y)
  =\mathscr H^{-i_0}(f,\mathscr P_\bullet)
\]
est un $\mathscr O_Y$-Module localement libre, dont le rang en $y\in Y$ est égal
à $(-1)^{i_0}\operatorname{EP}(f,\mathscr P_\bullet;y)$.
\end{proposition}

\begin{proof}
Remarquons d'abord que les hypothèses de (7.4.4) sont vérifiées par les
$\mathscr T_j^{\mathscr O_y}$, donc (7.4.7) leur est applicable, et l'hypothèse
entraîne que $\mathscr T_i^{\mathscr O_y}$ est \emph{nul} pour $i\ne i_0$ en
vertu de (7.5.3) ; en raison de (7.3.3), $\mathscr T_{i_0}$ est donc aussi exact,
et par suite (7.8.4), $\mathscr H^{-i_0}(f,\mathscr P_\bullet)$ est localement
libre et son rang en un point $y\in Y$ est
\[
  \operatorname{rang}_{\kappa(y)}
  \mathscr T_{i_0}(\kappa(y))
  =\operatorname{EP}(f,\mathscr P_\bullet;y)
\]
par définition, puisque $\mathscr T_i(\kappa(y))=0$ pour $i\ne i_0$.
\end{proof}

\begin{corollary}[7.9.9]
\label{III.7.9.9.fr}
Soient $Y$ un préschéma localement noethérien, $f:X\to Y$ un morphisme propre,
$\mathscr F$ un $\mathscr O_X$-Module cohérent et $Y$-plat ; on suppose qu'il
existe un entier $i_0$ tel que
\[
  H^i\bigl(f^{-1}(y),
  \mathscr F\otimes_{\mathscr O_Y}\kappa(y)\bigr)=0
  \qquad\text{pour tout $i\ne i_0$ et tout $y\in Y$.}
\]
Alors $\mathrm R^{i_0}f_*(\mathscr F)$ est un $\mathscr O_Y$-Module localement
libre, dont le rang en $y$ est égal à
$(-1)^{i_0}\operatorname{EP}(f,\mathscr F;y)$.
\end{corollary}

En particulier :

\begin{corollary}[7.9.10]
\label{III.7.9.10.fr}
Sous les conditions préliminaires de (7.9.9) pour $X$, $Y$ et $\mathscr F$,
supposons que l'on ait $\mathrm R^if_*(\mathscr F)=0$ pour tout $i>0$. Alors
$f_*(\mathscr F)$ est un $\mathscr O_Y$-Module localement libre, dont le rang en
$y$ est égal à $\operatorname{EP}(f,\mathscr F;y)$.
\end{corollary}

\begin{proof}
Il suffira, en vertu de (7.9.9) de prouver le lemme suivant :
\end{proof}

\begin{lemma}[7.9.10.1]
\label{III.7.9.10.1.fr}
Sous les hypothèses de (7.9.10), on a
\[
  H^i\bigl(f^{-1}(y),
  \mathscr F\otimes_{\mathscr O_Y}\kappa(y)\bigr)=0
\]
pour tout $i>0$ et tout $y\in Y$.
\end{lemma}

\begin{proof}
En effet, on peut se borner au cas où $Y=\operatorname{Spec}(A)$ est affine.
Avec les notations de (7.9.4), et $\mathscr P_\bullet$ étant réduit à son terme
de degré $0$ égal à $\mathscr F$, on a en effet
$\mathscr T_p(\mathscr O_Y)=0$ pour $p<0$ par hypothèse ; on conclut de (7.3.7)
que $\mathscr T_p$ est exact pour $p<0$, et le lemme résulte alors de
l'équivalence de (7.7.5, a)) et (7.7.5, d)).
\end{proof}

\oldpage[III]{211}
\begin{proposition}[7.9.11]
\label{III.7.9.11.fr}
Les hypothèses étant celles de (7.9.4), soit $\mathscr L$ un
$\mathscr O_X$-Module inversible très ample pour $Y$, et posons
$\mathscr P_\bullet(n)=\mathscr P_\bullet\otimes_{\mathscr O_X}
\mathscr L^{\otimes n}$ pour tout $n\in\mathbf Z$. Alors, pour tout $y\in Y$,
la fonction
\[
  n\longmapsto\operatorname{EP}(f,\mathscr P_\bullet(n);y)
  \tag{7.9.11.1}\label{III.7.9.11.1.fr}
\]
est un polynôme à coefficients dans $\mathbf Q$, qui est le même pour tous les
points d'une même composante connexe de $Y$.
\end{proposition}

\begin{proof}
Il est clair que $\mathscr P_\bullet(n)$ est un complexe de
$\mathscr O_X$-Modules $Y$-plats. En vertu de (7.9.6.2), on peut se borner au cas
où $\mathscr P_\bullet$ est réduit à un seul terme $\mathscr F\ne0$ de degré 0 ;
en outre, comme il s'agit de questions locales sur $Y$, on peut supposer $Y$
affine et $f$ projectif (II, 5.5.3) ; posons $X_y=f^{-1}(y)$, et soit
$\mathscr L_y=\mathscr L\otimes_{\mathscr O_Y}\kappa(y)$, qui est un
$\mathscr O_{X_y}$-Module très ample (II, 4.4.10) ; en vertu de (7.7.2), on a,
pour le foncteur $\mathscr T_\bullet$ relatif au complexe
$\mathscr P_\bullet(n)$,
\[
  \mathscr T_i(\kappa(y))
  =H^{-i}(X_y,\mathscr F_y\otimes\mathscr L_y^{\otimes n})
  \qquad
  \bigl(\text{où }\mathscr F_y=\mathscr F\otimes_{\mathscr O_Y}\kappa(y)\bigr) ;
\]
d'où résulte que $\operatorname{EP}(f,\mathscr F(n);y)$ n'est autre que la
caractéristique d'Euler-Poincaré $\chi_{\kappa(y)}(\mathscr F_y(n))$ définie dans
(2.5.1) ; le fait que (7.9.11.1) soit un polynôme résulte alors de (2.5.3) ; en
outre, pour chaque $n$, sa valeur est constante dans une composante connexe de
$Y$ (7.9.4), ce qui achève la démonstration.
\end{proof}

\begin{env}[7.9.12]
\label{III.7.9.12.fr}
Nous désignerons par $\operatorname{PH}(f,\mathscr P_\bullet;y)$ ou
$\operatorname{PH}(\mathscr P_\bullet;y)$ le polynôme (7.9.11.1), à coefficients
rationnels, et nous dirons que c'est le \emph{polynôme de Hilbert en $y$
relatif à $\mathscr P_\bullet$, $f$ et $\mathscr L$} (ou simplement le
\emph{polynôme de Hilbert en $y$ de $\mathscr P_\bullet$}, ou de $f$, s'il n'en
résulte pas de confusion) ; lorsque $Y$ est connexe non vide, on supprime la
mention de $y$ dans la notation et la terminologie. L'invariant ainsi obtenu
jouera un rôle essentiel au chap. V, dans la théorie des « modules » des
faisceaux cohérents quotients d'un faisceau cohérent donné.

Avec les notations de (7.9.6) et (7.9.11), on a
\[
  \operatorname{PH}(\mathscr P_\bullet;y)
  =\operatorname{PH}(\mathscr P'_\bullet;y)
  +\operatorname{PH}(\mathscr P''_\bullet;y)
  \tag{7.9.12.1}\label{III.7.9.12.1.fr}
\]
et en particulier
\[
  \operatorname{PH}(\mathscr P_\bullet;y)
  =\sum_i(-1)^i\operatorname{PH}(\mathscr P_i;y) ;
  \tag{7.9.12.2}\label{III.7.9.12.2.fr}
\]
cela résulte trivialement de (7.9.6.1) et (7.9.6.2). De même, avec les notations
et hypothèses de (7.9.7), on a
\[
  \operatorname{PH}(\mathscr P'_\bullet;y')
  =\operatorname{PH}(\mathscr P_\bullet;g(y')).
  \tag{7.9.12.3}\label{III.7.9.12.3.fr}
\]
\end{env}

La formule (7.9.12.2) ramène l'étude des polynômes de Hilbert d'un complexe à
celle des polynômes de Hilbert d'un seul $\mathscr O_X$-Module $Y$-plat. Ces
derniers admettent une interprétation remarquable indépendante de considérations
homologiques :

\begin{corollary}[7.9.13]
\label{III.7.9.13.fr}
Soient $Y$ un préschéma noethérien, $f:X\to Y$ un morphisme propre,
$\mathscr L$ un $\mathscr O_X$-Module inversible très ample pour $Y$,
$\mathscr F$ un $\mathscr O_X$-Module cohérent et $Y$-plat. Il existe un entier
$n_0$ tel que pour $n\geq n_0$, $f_*(\mathscr F(n))$ soit un
$\mathscr O_Y$-Module localement libre, de rang en $y\in Y$ égal à
$\operatorname{PH}(f,\mathscr F;y)(n)$.
\end{corollary}

\begin{proof}
Comme le morphisme $f$ est projectif (II, 5.5.3), il existe $n_0$ tel que pour
$n\geq n_0$

\oldpage[III]{212}
on ait $\mathrm R^if_*(\mathscr F(n))=0$ pour tout $i>0$ (2.2.1) ; la conclusion
résulte donc de (7.9.10).
\end{proof}

Le critère de platitude suivant sera important dans la théorie des « modules »
des faisceaux cohérents du chap. V :

\begin{proposition}[7.9.14]
\label{III.7.9.14.fr}
Soient $Y$ un préschéma noethérien, $f:X\to Y$ un morphisme projectif,
$\mathscr L$ un $\mathscr O_X$-Module inversible ample pour $f$, et posons
$\mathscr F(n)=\mathscr F\otimes\mathscr L^{\otimes n}$ pour tout
$\mathscr O_X$-Module $\mathscr F$ et tout $n\in\mathbf Z$. Pour qu'un
$\mathscr O_X$-Module cohérent $\mathscr F$ soit $Y$-plat, il faut et il suffit
qu'il existe un entier $n_0$ tel que, pour tout $n\geq n_0$,
$f_*(\mathscr F(n))$ soit un $\mathscr O_Y$-Module localement libre.
\end{proposition}

\begin{proof}
La nécessité de la condition se démontre comme dans (7.9.13) (le résultat de
(2.2.1) s'appliquant à un faisceau ample $\mathscr L$, puisque $f$ est
projectif). Pour démontrer la réciproque, on peut se borner au cas où $Y$ est
affine d'anneau $A$ ; en vertu de l'hypothèse et de (2.2.2, (i)), les
$A$-modules $\Gamma(X,\mathscr F(n))$ sont de type fini et projectifs (Bourbaki,
\emph{Alg. comm.}, chap. II, \S~5, n\textsuperscript{o}~2, th. 1). Soit $S$
l'anneau gradué
\[
  S=\bigoplus_{n\geq0}\Gamma(X,\mathscr L^{\otimes n}) ;
\]
on sait que $X$ s'identifie canoniquement à $\operatorname{Proj}(S)$ (II,
4.5.2, (b) et 5.4.4). Soit
\[
  M=\bigoplus_{n\geq n_0}\Gamma(X,\mathscr F(n)) ;
\]
remplaçant au besoin $\mathscr L$ par une puissance $\mathscr L^{\otimes d}$,
on peut supposer que $S$ est engendré par un nombre fini d'éléments de degré 1
(2.3.5.1), et il résulte alors de (II, 2.7.5 et 2.7.2) que $\mathscr F$
s'identifie à $\mathscr P\!roj_0(M)$. Pour tout élément homogène $g\in S$ de
degré $>0$, on a donc $\Gamma(X_g,\mathscr F)=M_{(g)}$ ; or, $M$, somme directe
de $A$-modules projectifs, est un $A$-module plat, donc il en est de même de
$M_g$ ($0_{\mathrm I}$, 6.3.2), et par suite aussi de $M_{(g)}$, qui est un
facteur direct de $M_g$ ($0_{\mathrm I}$, 6.1.2). On en conclut (I, 4.14.5) que
$\mathscr F$ est $Y$-plat en tout point de $X_g$, et comme les $X_g$ recouvrent
$X$, la proposition est démontrée.
\end{proof}

\bigskip
\noindent\emph{(A suivre.)}

\clearpage
\oldpage[III]{213}
\section*{INDEX DES NOTATIONS}

\begin{flushleft}
$\mathbf{Tor}_n^A(P_\bullet,Q_\bullet)$ : 6.3.1.

\medskip
$\mathscr{T}\!or_n^{\mathscr O_S}(\mathscr P_\bullet,\mathscr Q_\bullet)$,
$\mathscr{T}\!or_n^S(\mathscr P_\bullet,\mathscr Q_\bullet)$ : 6.4.1 et 6.5.1.

\medskip
$\mathscr{T}\!or_n^S(\mathscr F,\mathscr G)$ : 6.5.1.

\medskip
$\mathscr{T}\!or_q^S(\mathscr P_\bullet^{(1)},\mathscr P_\bullet^{(2)},\ldots,
\mathscr P_\bullet^{(m)})$ : 6.5.15.

\medskip
$\mathbf H^{-n}(\mathfrak U^{(i)},\mathscr P_\bullet^{(i)})$,
$\mathbf H^{-n}(X^{(i)},\mathscr P_\bullet^{(i)})$ : 6.6.1.

\medskip
$\mathbf{Tor}_n^A(L_{\bullet\bullet}^{(1)},L_{\bullet\bullet}^{(2)})$,
$\mathbf{Tor}_n^S(\mathfrak U^{(1)},\mathfrak U^{(2)};
\mathscr P_\bullet^{(1)},\mathscr P_\bullet^{(2)})$ : 6.6.2.

\medskip
$\mathscr{T}\!or_n^S(\mathfrak U^{(1)},\mathfrak U^{(2)};
\mathscr F^{(1)},\mathscr F^{(2)})$,
$\mathscr{T}\!or_n^S(\mathfrak U^{(1)},\mathfrak U^{(2)};
\mathscr P_\bullet^{(1)},\mathscr P_\bullet^{(2)})$ : 6.6.2.

\medskip
$\mathscr{T}\!or_n^S(f_1,f_2;\mathscr P_\bullet^{(1)},
\mathscr P_\bullet^{(2)})$ : 6.6.8, 6.7.1 et 6.7.3.

\medskip
$\mathscr{T}\!or_n^S((f_i)_{i\in I};(\mathscr P_\bullet^{(i)})_{i\in I})$,
$\mathscr{T}\!or_n^S(f_1,\ldots,f_m;\mathscr P_\bullet^{(1)},\ldots,
\mathscr P_\bullet^{(m)})$ : 6.7.11.

\medskip
$\mathbf{Ab}$, $\mathbf{Ab}_A$ : 7.1.1.

\medskip
$T_{(B)}$, $T^{(B)}$, $T\otimes_A B$
($T$ foncteur de $\mathbf{Ab}_A$ dans $\mathbf{Ab}$) : 7.1.3.

\medskip
$t_M$ : 7.2.2.

\medskip
$T_{\mathfrak p}$ : 7.1.4.

\medskip
$\mathscr T_\bullet^{Y'}(\mathscr M')$,
$\mathscr T_\bullet^U(\mathscr M|U)$ : 7.7.2.

\medskip
$\operatorname{EP}(f,\mathscr P_\bullet;y)$,
$\operatorname{EP}(Y,\mathscr P_\bullet;y)$,
$\operatorname{EP}(\mathscr P_\bullet;y)$ : 7.9.5.

\medskip
$\operatorname{PH}(f,\mathscr P_\bullet;y)$,
$\operatorname{PH}(\mathscr P_\bullet;y)$ : 7.9.12.
\end{flushleft}

\clearpage
\oldpage[III]{214}
\section*{INDEX TERMINOLOGIQUE}

\begin{flushleft}
Caractéristique d'Euler-Poincaré d'un complexe de Modules : 7.9.5.

\medskip
Cohomologiquement plat en un point $y\in Y$ en dimension $p$,
cohomologiquement plat sur $Y$ en dimension $p$, cohomologiquement plat sur
$Y$ au point $y$, cohomologiquement plat sur $Y$ (complexe de
$\mathscr O_X$-Modules $Y$-plats) : 7.8.1.

\medskip
Complexes homotopes : 6.1.4.

\medskip
Espace annelé de dimension cohomologique $\leq n$ : 6.5.5.

\medskip
Extension des scalaires dans un foncteur covariant additif : 7.1.3.

\medskip
Faisceau d'anneaux de dimension cohomologique $\leq n$ : 6.5.5.

\medskip
Formule de Künneth : 6.7.8.

\medskip
Homologiquement plat en un point $y\in Y$ en dimension $p$,
homologiquement plat sur $Y$ en dimension $p$, homologiquement plat sur $Y$
au point $y$, homologiquement plat sur $Y$ (complexe de $\mathscr O_X$-Modules
$Y$-plats) : 7.8.1.

\medskip
Homotopisme : 6.1.4.

\medskip
Hypertor de deux complexes de $A$-modules : 6.3.1.

\medskip
Hypertor local de deux complexes de Modules : 6.4.1.

\medskip
Hypertor global de deux complexes de Modules : 6.6.2, 6.6.8 et 6.7.3.

\medskip
Polynôme de Hilbert relatif à un complexe de Modules : 7.9.12.

\medskip
Suites spectrales de Künneth : 6.7.3.
\end{flushleft}

\clearpage
\oldpage[III]{215}
\section*{TABLE DES MATIÈRES}

\begin{flushleft}
\textsc{Chapitre III.} --- \textbf{Étude cohomologique des faisceaux cohérents
(suite)} \dotfill 5

\medskip
\S~6. Foncteurs Tor locaux et globaux ; formule de Künneth \dotfill 5

6.1. Introduction \dotfill 5

6.2. Hypercohomologie des complexes de Modules sur un préschéma \dotfill 6

6.3. Hypertor de deux complexes de modules \dotfill 9

6.4. Foncteurs hypertor locaux de complexes de Modules quasi-cohérents ; cas
des schémas affines \dotfill 14

6.5. Foncteurs hypertor locaux de complexes de Modules quasi-cohérents : cas
général \dotfill 16

6.6. Foncteurs hypertor globaux de complexes de Modules quasi-cohérents et
suites spectrales de Künneth : cas de la base affine \dotfill 21

6.7. Foncteurs hypertor globaux de complexes de Modules quasi-cohérents et
suites spectrales de Künneth : cas général \dotfill 25

6.8. Les suites spectrales d'associativité des hypertor globaux \dotfill 32

6.9. Les suites spectrales de changement de base dans les hypertor globaux
\dotfill 34

6.10. Structure locale de certains foncteurs cohomologiques \dotfill 39

\medskip
\S~7. Étude du changement de base dans les foncteurs homologiques covariants
de Modules \dotfill 43

7.1. Foncteurs de $A$-modules \dotfill 43

7.2. Caractérisation du foncteur produit tensoriel \dotfill 44

7.3. Critères d'exactitude des foncteurs homologiques de modules \dotfill 48

7.4. Critères d'exactitude pour les foncteurs
$\mathrm H_i(P\otimes_A M)$ \dotfill 53

7.5. Cas des anneaux locaux noethériens \dotfill 58

7.6. Descente des propriétés d'exactitude. Théorème de semi-continuité et
critère d'exactitude de Grauert \dotfill 60

7.7. Application aux morphismes propres : I. La propriété d'échange
\dotfill 65

7.8. Application aux morphismes propres : II. Critères de platitude
cohomologique \dotfill 72

7.9. Application aux morphismes propres : III. Invariance de la
caractéristique d'Euler-Poincaré et du polynôme de Hilbert \dotfill 76

\begin{flushright}
\emph{Reçu le 15 décembre 1962.}
\end{flushright}
\end{flushleft}

% La page imprimée 216 est blanche.
\clearpage
\thispagestyle{empty}
\null
\clearpage

\oldpage[III]{217}
\section*{ERRATA ET ADDENDA}
\begin{center}
(Liste 2)

\medskip
A) \emph{Erreurs typographiques}
\end{center}

\noindent\textbf{$(0_{\mathrm I},\,3.2.1)$} Ligne 2 de la p. 27, remplacer $X$
par $U$.

\noindent\textbf{$(0_{\mathrm I},\,3.2.6)$} Ligne 1 du bas de la p. 27 et ligne
1 de la p. 28, remplacer (3 fois) $\mathfrak S_\lambda$ par
$\mathscr H_\lambda$.

\noindent\textbf{$(0_{\mathrm I},\,3.4.5)$} Ajouter une parenthèse devant
3.4.5).

\noindent\textbf{$(0_{\mathrm I},\,4.2.1)$} Ligne 10 de la p. 39, remplacer
$\psi_*(A)$ par $\psi_*(\mathscr A)$.

\noindent\textbf{$(0_{\mathrm I},\,5.1.3)$} Ligne 16 de la p. 45, remplacer
$\mathscr F|V$ par $\mathscr F|U$.

\noindent\textbf{$(0_{\mathrm I},\,5.3.9)$} Ligne 12 de la p. 48, avant
« telle que », ajouter : au-dessus d'un voisinage ouvert $U$ de $x$.

\noindent\textbf{$(0_{\mathrm I},\,7.6.9)$} Ligne 3 du bas de la p. 73,
remplacer $I_\lambda$ par $J_\lambda$.

\noindent\textbf{(I, 1.1.15)} Ligne 3 de la p. 83, remplacer $\mathfrak a$ par
$\mathfrak a\ne A$.

\noindent\textbf{(I, 1.4.1)} Ligne 2 du bas de la p. 90, remplacer
$\widetilde N$ par $N$.

\noindent\textbf{(I, 2.3.1)} Ligne 3 du bas de la p. 101, remplacer (0, 4.1.6)
par (0, 4.1.7).

\noindent\textbf{(I, 2.3.2)} Ligne 11 de la p. 101, remplacer $B$ par $B_s$ et
$C$ par $C_t$.

\noindent\textbf{(I, 2.5.5)} Ligne 19 de la p. 104, remplacer $S$-morphisme
par $S$-préschéma.

\noindent\textbf{(I, 3.3.7)} Ligne 17 de la p. 109, remplacer $Y_{(S)}$ par
$Y_{(S')}$.

\noindent\textbf{(I, 3.4.5)} Ligne 9 de la p. 113, remplacer $\mathrm s$ par
$s$.

\noindent\textbf{(I, 3.4.8)} Ligne 4 de la p. 114, remplacer $p(x)$ par
$f(x)$.

\noindent\textbf{(I, 3.5.1)} Ligne 3 de la p. 115, remplacer $1_X\times g$
par $1_{X'}\times g$.

\noindent\textbf{(I, 3.7.2)} Ligne 8 de la p. 119, remplacer le premier $X$
par $X'$.

\noindent\textbf{(I, 4.2.2)} Ligne 8 de la p. 123, dans la flèche verticale de
gauche du diagramme, remplacer $\alpha_{\psi(y)}$ par $\varphi_{\psi(y)}$.

\noindent\textbf{(I, 4.4.3)} Ligne 16 de la p. 126, remplacer isomorphée par
isomorphes.

\noindent\textbf{(I, 4.5.5)} Ligne 17 du bas de la p. 127, remplacer (4.2.4)
par (4.2.5).

\noindent\textbf{(I, 5.1.4)} Ligne 14 du bas de la p. 128, remplacer (2.1.7)
par (2.1.8).

\noindent\textbf{(I, 5.3.13)} Ligne 20 de la p. 134, remplacer (4.2.4) par
(4.2.5).

\noindent\textbf{(I, 6.4.2)} Ligne 1 du bas de la p. 147, remplacer
recouvrement par recouvrement fini.

\noindent\textbf{(I, 9.1.13)} Ligne 15 de la p. 171, remplacer
$p^{-1}(\mathscr F)$ par $p^*(\mathscr F)$.

\noindent\textbf{(I, 9.5.11)} Ligne 7 de la p. 179, remplacer $Y$ par $Y'$.

\noindent\textbf{(I, 9.6.5)} Ligne 17 du bas de la p. 180, remplacer
$(0,\,4.1.4)$ par $(0,\,4.1.3)$.

\noindent\textbf{(I, 10.12.2)} Ligne 14 du bas de la p. 206, remplacer
$(X_n)_{(S_n)}$ par $(X_n)_{(S_m)}$.

\clearpage
\oldpage[III]{218}
\noindent\textbf{(I), Index terminologique :} Ligne 6 du bas de la p. 219,
remplacer 0, 4.1.4 par 0, 4.1.5. Lignes 16, 17, 18 du bas de la p. 222,
remplacer I, 2.1.7 par I, 2.1.8 ; ligne 19 du bas de la p. 222, remplacer I,
2.1.8 par I, 2.1.7.

\noindent\textbf{(II, 1.5.2)} Ligne 17 de la p. 12, insérer un $f$ à côté de
la flèche verticale de gauche du diagramme.

\noindent\textbf{(II, 4.2.7)} Ligne 16 du bas de la p. 75, remplacer (III,
2.1.14) par (III, 2.1.13).

\noindent\textbf{(II, 5.2.1)} Ligne 16 du bas de la p. 97, remplacer $X-U$ par
$U$. Ligne 8 du bas de la p. 97, remplacer $X'$ par $X_f$.

\noindent\textbf{(II, 5.3.6)} Ligne 16 de la p. 100, remplacer (4.6.18) par
(4.6.17).

\noindent\textbf{(II, 6.2.7)} Ligne 9 du bas de la p. 116, remplacer chap. V
par chap. IV.

\noindent\textbf{(II, 6.6.5)} Ligne 15 du bas de la p. 132, remplacer chap. V
par chap. IV.

\noindent\textbf{(II, 7.4.12)} Ligne 1 de la p. 151, remplacer chap. V par
chap. III.

\noindent\textbf{(II, 8.1.4)} Ligne 15 de la p. 154, remplacer (III, 2.3.8)
par (III, 2.3.7).

\noindent\textbf{(II, 8.10.5)} Ligne 17 du bas de la p. 187, remplacer
$g_{(j)x}$ par $g_{j(x)}$.

\noindent\textbf{$(0_{\mathrm{III}},\,10.2.7)$} Ligne 1 de la p. 20,
remplacer $u$-adique par $\mathfrak m$-adique. Ligne 13 du bas de la p. 20,
remplacer $k$ par $k_i$.

\noindent\textbf{$(0_{\mathrm{III}},\,11.1.1)$} Ligne 18 de la p. 23,
remplacer $\operatorname{Im}(d_r^{p+r,q-r+1})$ par
$\operatorname{Im}(d_r^{p-r,q+r-1})$.

\noindent\textbf{$(0_{\mathrm{III}},\,11.8.6)$} Ligne 14 du bas de la p. 45,
remplacer les flèches $\to$ par $\rightsquigarrow$ (4 fois).

\noindent\textbf{$(0_{\mathrm{III}},\,12.3.4)$} Ligne 4 de la p. 61,
remplacer $\rightsquigarrow$ par $\to$.

\noindent\textbf{$(0_{\mathrm{III}},\,13.2.4)$} Ligne 2 du bas de la p. 67,
remplacer $\varinjlim$ par $\varprojlim$ (2 fois).

\noindent\textbf{(III, 1.1.7)} Ligne 15 de la p. 85, remplacer $A^r$ par
$\bigwedge(A^r)$ (2 fois).

\noindent\textbf{(III, 1.2.4)} Ligne 7 de la p. 87, remplacer
$\Gamma(U,\mathscr F)$ par $\Gamma(U,\mathscr O_X)$ et remplacer $\mathscr G$
par $\mathscr F$. Ligne 8 de la p. 87, remplacer
$\mathrm H^p(\mathfrak U,\mathscr G)$ par $\mathrm H^p(\mathfrak U,\mathscr F)$.

\noindent\textbf{(III, 2.5.3.1)} Ligne 3 de la p. 111, remplacer $n>0$ par
$n\geq0$. Ligne 8 de la p. 111, remplacer $-r\leq n\leq0$ par
$-r\leq n<0$.

\noindent\textbf{(III, 3.1.2)} Ligne 13 de la p. 116, remplacer $\mathscr G$
par $\mathscr G\in\mathbf K'$.

\noindent\textbf{(III, 4.3.6)} Ligne 7 de la p. 133, remplacer
$K'=K\otimes_A\widehat A$ par $K'\supset K\otimes_A\widehat A$. Ligne 16 de
la p. 133, remplacer $R(X)$ par $R(Y)$.

\noindent\textbf{(III, 4.4.12)} Ligne 10 de la p. 138, remplacer chap. V par
chap. IV.

\noindent\textbf{(III, 4.6.6)} Ligne 19 de la p. 142, remplacer « chap. V » par
« dans un paragraphe ultérieur ».

\begin{center}
B) \emph{Modifications de texte}\footnote{Pour faciliter la recherche des
références, les modifications appartenant à la liste d'errata insérée dans le
chapitre N seront désormais désignées par $\mathbf{Err}_{\mathrm N}$ suivi d'un
numéro.}
\end{center}

\noindent\textbf{$(\mathbf{Err}_{\mathrm{III}},\,1)$} Dans
$(0_{\mathrm I},\,5.1.3)$, ligne 18 de la p. 45, remplacer « somme directe »
par « somme directe finie ».

\noindent\textbf{$(\mathbf{Err}_{\mathrm{III}},\,2)$} Dans
$(0_{\mathrm I},\,5.4.3)$, lignes 8 à 4 du bas de la p. 49, remplacer depuis
« Dans le cas général\ldots{} » par le texte suivant :

Dans le cas général où $X$ est un espace annelé tel que $\mathscr O_x$ soit un
anneau \emph{local} pour
\oldpage[III]{219}
tout $x\in X$, si $\mathscr L$ est un $\mathscr O_X$-Module \emph{de type fini}
tel qu'il existe un $\mathscr O_X$-Module $\mathscr F$ pour lequel
$\mathscr L\otimes_{\mathscr O_X}\mathscr F$ soit \emph{isomorphe à}
$\mathscr O_X$, le raisonnement précédent montre que pour tout $x\in X$,
$\mathscr L_x$ est isomorphe à $\mathscr O_x$. On en déduit que $\mathscr L$ est
alors \emph{inversible}. En effet, pour tout $x\in X$, soient $U$ un voisinage
ouvert de $x$ tel que $\mathscr L|U$ soit engendré par $n$ sections
$s_i\ (1\leq i\leq n)$ au-dessus de $U$ (5.2.1). On peut supposer par exemple
que $s=s_1$ est telle que $s_x\ne0$, et comme $\mathscr L_x$ est isomorphe à
$\mathscr O_x$, il existe pour chaque $i$ une section $t_i$ de $\mathscr O_X$
au-dessus d'un voisinage ouvert $V_i\subset U$ de $x$ telle que
$(t_i)_x s_x=(s_i)_x$. Il y a par suite un voisinage ouvert
$V\subset\bigcap_{i=1}^nV_i$ de $x$ tel que $t_i s=s_i$ dans $V$, autrement
dit $\mathscr L|V$ est engendré par \emph{l'unique} section $s|V$. En outre,
si $z$ est une section au-dessus d'un ouvert $W\subset V$ du noyau de
l'homomorphisme $\mathscr O_X|V\to\mathscr L|V$ défini par $s$ (5.1.1), $z_y$
annule $\mathscr L_y$ pour tout $y\in W$, donc $z_y=0$ par hypothèse et par
suite $z=0$, ce qui achève de prouver notre assertion. De plus, la
considération du produit tensoriel
$\mathscr L^{-1}\otimes\mathscr L\otimes\mathscr F$ montre aussitôt que
$\mathscr F$ est isomorphe à $\mathscr L^{-1}$.

\noindent\textbf{$(\mathbf{Err}_{\mathrm{III}},\,3)$} Dans
$(0_{\mathrm I},\,7.2.4)$, il faut imposer la condition que les puissances
$\mathfrak J^n$ sont des idéaux \emph{fermés} dans l'anneau admissible $A$ pour
que la proposition soit exacte. La démonstration donnée est incorrecte, et
l'énoncé rectifié résulte de Bourbaki, \emph{Top. gén.}, chap. III,
3\textsuperscript{e} éd., \S~3, n\textsuperscript{o}~5, cor. 1 de la prop. 9.
Il faut de même supposer les $\mathfrak J^n$ \emph{fermés} dans $A$ dans les
énoncés $(0_{\mathrm I},\,7.2.5)$ et $(0_{\mathrm I},\,7.2.6)$.

\noindent\textbf{$(\mathbf{Err}_{\mathrm{III}},\,4)$} Après la proposition
(I, 2.2.5), ajouter : avec les notations de (I, 2.2.5), si $\mathfrak J$ est un
idéal de $A$, on note $\widetilde{\mathfrak J\mathscr F}$ le sous-$\mathscr O_X$-Module
$\mathfrak J\widetilde{\mathscr F}$ défini dans $(0,\,4.3.5)$.

\noindent\textbf{$(\mathbf{Err}_{\mathrm{III}},\,5)$} Dans (I, 3.4.5), ligne 1
du bas de la p. 112, remplacer « \emph{un corps} $K$ » par « \emph{un corps
algébriquement clos} $K$ ». Ligne 1 de la p. 113, remplacer « \emph{extension} »
par « \emph{extension algébriquement close} ». Lignes 1 à 4 de la p. 113,
remplacer depuis « $K$ sera appelé\ldots{} » par le texte suivant :

Pour tout point de $X$ à valeurs dans un corps $K$, $K$ sera appelé le
\emph{corps des valeurs} du point correspondant, et si $x$ est la localité de
ce point on dit encore que ce dernier est \emph{localisé en} $x$. On définit
ainsi une application $X(K)\to X$, faisant correspondre à un point à valeurs
dans $K$ sa localité.

Lignes 7 et 16 de la p. 113, supprimer « géométrique » ; lignes 10 et 13 de
la p. 113, supprimer « géométriques ». Lignes 10 et 11 du bas de la p. 115,
supprimer « \emph{géométrique} ».

\noindent\textbf{$(\mathbf{Err}_{\mathrm{III}},\,6)$} À la fin de la
démonstration de (I, 3.6.1), ligne 12 du bas de la p. 117, ajouter : Si l'on
pose
$X'=X\times_Y\operatorname{Spec}(\mathscr O_y/\mathfrak a_y)$, alors, pour tout
point $x\in X'$, identifié par $p$ à un point de $X$, on a
$\mathscr O_{X',x}=\mathscr O_{X,x}/\mathfrak a_y\mathscr O_{X,x}$. La question
étant en effet locale sur $X$ et $Y$, on peut supposer que
$X=\operatorname{Spec}(B)$, $Y=\operatorname{Spec}(A)$, et l'on a
$(B/\mathfrak a_yB)_x=B_x/\mathfrak a_yB_x$ par platitude $(0,\,1.3.2)$.

\noindent\textbf{$(\mathbf{Err}_{\mathrm{III}},\,7)$} Dans (I, 5.1.1), ligne 3
du bas de la p. 127, remplacer « $\mathscr O_X$-Module quasi-cohérent » par
« \emph{Idéal quasi-cohérent de} $\mathscr B$ ».

\noindent\textbf{$(\mathbf{Err}_{\mathrm{III}},\,8)$} Après (I, 5.1.10),
ligne 16 de la p. 131, ajouter :

\emph{Remarque (5.1.11).} --- Une légère adaptation du raisonnement fait dans
(5.1.9) montre que la conclusion reste valable \emph{sans supposer a priori
que $X$ soit un préschéma}, mais

\oldpage[III]{220}
en supposant seulement que $(X,\mathscr O_X)$ soit un espace annelé en anneaux
locaux, $\mathscr J$ un Idéal de $\mathscr O_X$ tel que $\mathscr J^n=0$, que
l'espace annelé $X_0=(X,\mathscr O_X/\mathscr J)$ soit un schéma affine et enfin
que les Idéaux $\mathscr J^k/\mathscr J^{k+1}$ soient des
$\mathscr O_{X_0}$-Modules quasi-cohérents. Il suffit en effet d'utiliser
(1.8.1) (cf. $(\mathbf{Err}_{\mathrm{II}})$) au lieu de (2.2.4) dans le
raisonnement.

\noindent\textbf{$(\mathbf{Err}_{\mathrm{III}},\,9)$} Dans (I, 5.3.1), ligne
11 de la p. 132, insérer, après « ou $\Delta_X$ », « ou $\Delta_\varphi$ si
$\varphi:X\to S$ est le morphisme structural ».

\noindent\textbf{$(\mathbf{Err}_{\mathrm{III}},\,10)$} Dans (I, 5.3.9), la
démonstration est insuffisante, car elle ne prouve pas que $\Delta_X(X)$ soit
localement fermé dans $X\times_SX$. Pour avoir une démonstration correcte, il
suffit d'utiliser (4.2.4, $a$)) : pour tout $x\in X$ et tout voisinage affine
$U$ de $x$ dans $X$, $U\times_SU$ est un voisinage affine de $\Delta_X(x)$ ;
compte tenu de (5.3.16) (dont la démonstration n'utilise que la définition
(5.3.1.1)), on est ramené à démontrer (5.3.9) lorsque
$S=\operatorname{Spec}(B)$ et $X=\operatorname{Spec}(A)$ sont des schémas
affines ; il est clair alors en vertu de (5.3.1.1) que $\Delta_X$ correspond à
l'homomorphisme canonique $A\otimes_BA\to A$ qui transforme $x\otimes y$ en
$xy$ ; cet homomorphisme étant surjectif, $\Delta_X$ est dans ce cas une
immersion fermée (4.2.3), ce qui achève la démonstration.

\noindent\textbf{$(\mathbf{Err}_{\mathrm{III}},\,11)$} Dans (I, 7.1.15) et
(I, 7.1.16), lignes 7 et 15 du bas de la p. 158, supprimer
« \emph{géométriques} ».

\noindent\textbf{$(\mathbf{Err}_{\mathrm{III}},\,12)$} Remplacer (I, 7.4.7)
par : On étend (par abus de langage) les définitions de (7.4.1) au cas où
$X$ est un préschéma \emph{réduit} dont tout point admet un voisinage ouvert
n'ayant qu'un nombre \emph{fini} de composantes irréductibles ; il résulte
alors de (7.3.4) et (7.4.6) que, pour un $\mathscr O_X$-Module quasi-cohérent de
type fini $\mathscr F$, dire que $\mathscr F$ est un \emph{faisceau de torsion}
équivaut à dire que $\operatorname{Supp}(\mathscr F)$ \emph{ne contient aucune
composante irréductible de $X$}.

\noindent\textbf{$(\mathbf{Err}_{\mathrm{III}},\,13)$} Dans (I, 10.11.7),
lignes 17 à 23 de la p. 205, la fin de la démonstration depuis « Reste à
prouver\ldots{} » est inutile en vertu de (10.11.6), les faisceaux considérés
étant cohérents par $(0,\,5.3.5)$.

\noindent\textbf{$(\mathbf{Err}_{\mathrm{III}},\,14)$} Dans (II, 1.7.8),
après la ligne 10 du bas de la p. 16, ajouter : lorsque
$\mathscr E=\mathscr O_S^n$, on écrit aussi $\mathbf V_S^n$ au lieu de
$\mathbf V(\mathscr E)$ ; si en outre $S=\operatorname{Spec}(A)$ est affine,
on écrit $\mathbf V_A^n$ au lieu de $\mathbf V_S^n$ ; on a
$\mathbf V_A^n=\operatorname{Spec}(A[T_1,\ldots,T_n])$ où les $T_i$ sont des
indéterminées.

\noindent\textbf{$(\mathbf{Err}_{\mathrm{III}},\,15)$} Dans (II, 1.7.10),
ligne 10 de la p. 17, supprimer « \emph{géométriques} » ; lignes 12 et 16--17
de la p. 17, supprimer « géométrique ».

\noindent\textbf{$(\mathbf{Err}_{\mathrm{III}},\,16)$} Dans (II, 1.7.12),
ajouter : On pose de même :
\[
  \mathbf V(\mathscr O_S^n)=\mathbf V_S^n=S[T_1,\ldots,T_n].
\]

\noindent\textbf{$(\mathbf{Err}_{\mathrm{III}},\,17)$} Dans (II, 4.2.6),
ligne 7 de la p. 75, supprimer « \emph{géométriques} » ; ligne 8 de la p.
75, supprimer « \emph{géométrique} ».

\noindent\textbf{$(\mathbf{Err}_{\mathrm{III}},\,18)$} Dans (II, 4.6.13),
ligne 7 de la p. 92, le raisonnement doit être précisé, car avec les notations
de (4.4.10), il faut ici prouver que l'immersion $\Gamma_f$ est
\emph{quasi-compacte}, afin d'appliquer (i bis). En vertu de (4.6.4), pour
prouver que $\mathscr L$ est ample relativement à $f$, on peut se borner au
cas où $Y$ est affine. Notons d'autre part que dans chacune des hypothèses de
(v), $f$ est quasi-compact (I, 6.6.4). D'autre part, si $g$ est
\oldpage[III]{221}
séparé, $\Gamma_f$ est une immersion fermée (I, 5.4.3) donc quasi-compacte
(I, 6.6.4). Si au contraire $X$ est localement noethérien, comme $Y$ est affine
et $f$ quasi-compact, l'espace sous-jacent à $X$ est quasi-compact, donc
noethérien, et on peut de nouveau appliquer (I, 6.6.4) pour prouver que
$\Gamma_f$ est quasi-compacte.

\noindent\textbf{$(\mathbf{Err}_{\mathrm{III}},\,19)$} L'énoncé de (II,
5.2.2) est incorrect, les conditions $b)$, $c)$ et $c')$ n'étant pas locales
sur $Y$ lorsqu'on ne sait pas si pour un ouvert $U$ de $Y$, un
$(\mathscr O_X|f^{-1}(U))$-Module quasi-cohérent est restriction à $f^{-1}(U)$
d'un $\mathscr O_X$-Module quasi-cohérent. Il faut donc remplacer dans la ligne
16 de la p. 98 « \emph{Soit $f:X\to Y$ un morphisme séparé quasi-compact} »
par : « \emph{Soient $X$, $Y$ deux préschémas tels que $X$ soit un schéma ou
que l'espace sous-jacent à $X$ soit localement noethérien, et soit
$f:X\to Y$ un morphisme quasi-compact.} » Dans la démonstration, on observera
que les hypothèses entraînent que $f$ est \emph{séparé} lorsqu'on suppose que
$X$ est un schéma, en vertu de (I, 5.5.5 et 5.5.8) ; on peut donc appliquer
(I, 9.2.2) dans les deux cas.

\noindent\textbf{$(\mathbf{Err}_{\mathrm{III}},\,20)$} Dans (II, 6.2.3),
ajouter, après la ligne 18 du bas de la p. 115 :

\emph{On dit qu'un morphisme $f:X\to Y$ est quasi-fini en un point $x\in X$
s'il existe un voisinage ouvert affine $V$ de $y=f(x)$ et un voisinage ouvert
affine $U$ de $x$ tels que $f(U)\subset V$ et que le morphisme $U\to V$
restriction de $f$ soit quasi-fini. On dit qu'un morphisme $f:X\to Y$ est
localement quasi-fini s'il est quasi-fini en tout point de $X$.}

\noindent\textbf{$(\mathbf{Err}_{\mathrm{III}},\,21)$} Dans (II, 6.4.3), la
démonstration est incorrecte, les éléments d'un sous-$A$-module de type fini de
$K$ n'étant pas nécessairement entiers sur $A$. Remplacer les 9 dernières lignes
de la p. 121 par le texte suivant : Comme on peut supposer que $E$ est sans
torsion, donc fidèle, $E$ est aussi un $A[u]$-module fidèle ($A$ étant plongé
dans l'anneau des endomorphismes de $E$). Comme $E$ est un $A$-module de type
fini, il en résulte que $u$ est \emph{entier} sur $A$ (Bourbaki,
\emph{Alg. comm.}, chap. V, \S~1, n\textsuperscript{o}~1, lemme 1). On en
conclut immédiatement que les valeurs propres de $u\otimes1$ (dans une clôture
algébrique de $K$) sont des éléments \emph{entiers} sur $A$, et il en est donc
de même des $\sigma_i(u)$.

\noindent\textbf{$(\mathbf{Err}_{\mathrm{III}},\,22)$} Dans
$(0_{\mathrm{III}},\,9.1.1)$, ligne 18 du bas de la p. 12, supprimer
« ouverte ».

\noindent\textbf{$(\mathbf{Err}_{\mathrm{III}},\,23)$} Dans
$(0_{\mathrm{III}},\,11.7.3)$, ligne 10 de la p. 43, après $\mathsf C''$,
ajouter : tel que pour tout objet projectif $P$ de $\mathsf C'$ (resp. tout
objet projectif $P'$ de $\mathsf C''$) le foncteur
$A'\rightsquigarrow T(P,A')$ (resp. $A\rightsquigarrow T(A,P')$) soit exact
dans $\mathsf C'$ (resp. $\mathsf C$).

\noindent\textbf{$(\mathbf{Err}_{\mathrm{III}},\,24)$} L'énoncé de
$(0_{\mathrm{III}},\,13.7.7)$ est inexact, et doit être modifié comme suit :
ligne 6 de la p. 78, supprimer « et
$(\mathrm R^{n+1}T(A_k))_{k\in\mathbf Z}$ » ; ligne 7 de la p. 78, remplacer
« $\mathrm R'^nT(A)$ est un $S$-module de type fini » par
« $\mathrm R'^nT(A)$ et $\mathrm R'^{n+1}T(A)$ sont des $S$-modules de type
fini » ; lignes 12--13 de la p. 78, supprimer « et $p+q=n+1$ ». Dans la
démonstration, ligne 23 de la p. 78, supprimer « et pour $n+1$ ».

\noindent\textbf{$(\mathbf{Err}_{\mathrm{III}},\,25)$} Dans (III, 1.4.15),
ligne 6 du bas de la p. 92, remplacer « \emph{de type fini} » par
« \emph{quasi-compact} ».

\noindent\textbf{$(\mathbf{Err}_{\mathrm{III}},\,26)$} Dans (III, 2.2.4),
ligne 1 de la p. 102, remplacer « \emph{que les supports de $\mathscr F$ et de
$\mathscr H$ soient propres sur $Y$} » par « \emph{que le support de
$\operatorname{Im}(\mathscr F\to\mathscr G)$ soit propre sur $Y$} ». Dans la
démonstration, remplacer le texte des lignes 10 à 16, depuis « Cela
étant\ldots{} », par :

\oldpage[III]{222}
Posons
$\mathscr G_1=\operatorname{Im}(\mathscr F\to\mathscr G)
=\operatorname{Ker}(\mathscr G\to\mathscr H)$, qui est cohérent. Comme on a la
suite exacte
$0\to\mathscr G_1(n)\to\mathscr G(n)\to\mathscr H(n)$ et que le foncteur
$f_*$ est exact à gauche, la suite
$0\to f_*(\mathscr G_1(n))\to f_*(\mathscr G(n))\to f_*(\mathscr H(n))$ est
exacte pour tout $n$, et il suffit donc de montrer que, pour $n$ assez grand,
la suite $f_*(\mathscr F(n))\to f_*(\mathscr G_1(n))\to0$ est exacte. Autrement
dit, on peut se borner au cas où $\mathscr H=0$ et où le support de
$\mathscr G$ est propre sur $Y$, donc \emph{fermé dans $Z$} (II, 5.4.10) ;
$\mathscr G'=i_*(\mathscr G)$ est par suite un $\mathscr O_Z$-Module cohérent tel
que $\mathscr G'|X=\mathscr G$. On sait (I, 9.4.3) qu'il existe un
$\mathscr O_Z$-Module cohérent $\mathscr F'$ tel que
$\mathscr F'|X=\mathscr F$. En outre, comme $X$ est ouvert dans $Z$ et
$\operatorname{Supp}(\mathscr G)\subset X$ fermé dans $Z$, il est immédiat que
l'on définit un homomorphisme surjectif $u':\mathscr F'\to\mathscr G'$ de
faisceaux tel que $u'|X$ soit l'homomorphisme surjectif donné
$u:\mathscr F\to\mathscr G$, en prenant $u'|U=0$ pour tout ouvert $U$ de $Z$
ne rencontrant pas $\operatorname{Supp}(\mathscr G)$ et $u'|U=u|U$ pour tout
ouvert $U\subset X$. Cela étant, on voit comme au début de la démonstration de
(2.2.2) que l'on peut se borner au cas où $Y$ est affine, et il s'agit donc de
montrer que, pour $n$ assez grand, l'homomorphisme
$\Gamma(u):\Gamma(X,\mathscr F(n))\to\Gamma(X,\mathscr G(n))$ est surjectif. Or,
on a le diagramme commutatif
\[
\xymatrix@C=35mm@R=12mm{
\Gamma(Z,\mathscr F'(n)) \ar[r]^{\Gamma(u')} \ar[d]_v &
\Gamma(Z,\mathscr G'(n)) \ar[d]^w \\
\Gamma(X,\mathscr F(n)) \ar[r]_{\Gamma(u)} & \Gamma(X,\mathscr G(n))
}
\]
où $v$ et $w$ sont les homomorphismes de restriction. La définition de
$\mathscr G'$ montre en outre que $w$ est \emph{bijectif} ; par ailleurs, en
vertu de (2.2.3), $\Gamma(u')$ est \emph{surjectif} pour $n$ assez grand, donc
il en est de même de $\Gamma(u)$.

\noindent\textbf{$(\mathbf{Err}_{\mathrm{III}},\,27)$} Dans (III, 2.2.5),
après la ligne 7 du bas de la p. 102, ajouter :

(iii) Sous les hypothèses de (2.2.4) concernant $X$, $Y$, $f$ et $\mathscr L$,
il est immédiat que si $\mathscr H$ est un $\mathscr O_X$-Module cohérent dont
le support est \emph{propre sur $Y$}, un raisonnement analogue à celui de
(2.2.4) montre qu'il existe un entier $N$ tel que pour $n\geq N$, on ait
$\mathrm R^if_*(\mathscr H(n))=0$ pour tout $i>0$. On en conclut, par la suite
exacte de cohomologie, que si $u:\mathscr F\to\mathscr G$ est un homomorphisme
de $\mathscr O_X$-Modules cohérents tel que $\operatorname{Ker}(u)$ et
$\operatorname{Coker}(u)$ aient leurs supports \emph{propres sur $Y$}, alors il
existe $N$ tel que pour $n\geq N$, l'homomorphisme correspondant
$\mathrm R^if_*(\mathscr F(n))\to\mathrm R^if_*(\mathscr G(n))$ soit
\emph{bijectif} pour tout $i>0$.

\noindent\textbf{$(\mathbf{Err}_{\mathrm{III}},\,28)$} Dans (III, 4.4.9),
ligne 17 du bas de la p. 137, remplacer « \emph{Soient $Y$ un préschéma
intègre localement noethérien} » par « \emph{Soient $X$ et $Y$ deux
préschémas intègres localement noethériens} » ; ligne 16 du bas de la p.
137, remplacer « \emph{de type fini} » par « \emph{localement de type
fini} ». La démonstration est essentiellement inchangée,
$f^{-1}(y)$ étant localement de type fini sur $k(y)$, donc encore discret.

\noindent\textbf{$(\mathbf{Err}_{\mathrm{III}},\,29)$} dans (I, 9.3.4), ligne
19 du bas de la p. 173, remplacer « \emph{noethérien} » par
« \emph{quasi-compact} » ; lignes 18 et 19 du bas de la p. 173, remplacer
« \emph{cohérent} » par « \emph{quasi-cohérent de type fini} ». Dans la
démonstration, on se ramène au cas où

\oldpage[III]{223}
$X=\operatorname{Spec}(A)$, $\mathscr F=\widetilde M$, où $M$ est un $A$-module
de type fini, $\mathscr J=\widetilde{\mathfrak J}$, où $\mathfrak J$ est un
idéal de type fini de $A$ ; le reste du raisonnement est alors inchangé.

\noindent\textbf{$(\mathbf{Err}_{\mathrm{III}},\,30)$} Dans (I, 9.3.5),
remplacer les lignes 1 à 7 du bas de la p. 173 par le texte suivant :

\emph{Proposition (9.3.5). --- Soient $X$ un préschéma, $\mathscr F$ un
$\mathscr O_X$-Module quasi-cohérent de type fini. Alors il existe un
sous-préschéma fermé $Y$ de $X$, dont l'espace sous-jacent est égal à
$\operatorname{Supp}(\mathscr F)$, et un $\mathscr O_Y$-Module quasi-cohérent de
type fini $\mathscr G$ tels que, si $j:Y\to X$ est l'injection canonique,
$\mathscr F$ soit isomorphe à $j_*(\mathscr G)$.}

Il suffira de montrer que l'Idéal $\mathscr J$ de $\mathscr O_X$,
\emph{annulateur de $\mathscr F$}, est \emph{quasi-cohérent} ; on prendra alors
pour $Y$ le sous-préschéma fermé de $X$ défini par $\mathscr J$ (I, 4.1.2), et
comme $\mathscr J\mathscr F=0$, $\mathscr F$ est un
$(\mathscr O_X/\mathscr J)$-Module et on répondra à la question en prenant
$\mathscr G=j^*(\mathscr F)$. Pour voir que $\mathscr J$ est quasi-cohérent, on
peut (la question étant locale) se borner au cas où
$X=\operatorname{Spec}(A)$, $\mathscr F=\widetilde M$, où $M$ est un $A$-module
engendré par un nombre fini d'éléments $x_i\ (1\leq i\leq r)$ ; l'Idéal
$\mathscr J$ est alors l'intersection des annulateurs des $x_i$. Mais
l'annulateur de $x_i$ est le noyau de l'homomorphisme
$\mathscr O_X\to\mathscr F$ correspondant à l'homomorphisme $s\mapsto sx_i$ de
$A$ dans $M$ ; c'est donc bien un Idéal quasi-cohérent (I, 4.1.1) et toute
intersection finie de tels Idéaux est aussi un Idéal quasi-cohérent (I,
1.3.10).
